% !TeX program = pdflatex
% papers/corpo\_analitico.tex — o paper principal: a especificação do sistema.
% Auto-contido: não depende de teoria.tex nem de catalogo.tex para ser lido.
% Compila sozinho:  pdflatex corpo\_analitico.tex
\documentclass[11pt,a4paper]{article}
\usepackage[utf8]{inputenc}\usepackage[T1]{fontenc}\usepackage[portuguese]{babel}
\usepackage{amsmath,amssymb,amsthm}
\usepackage[margin=2.3cm]{geometry}
\usepackage{booktabs}\usepackage{xcolor}
\usepackage[hidelinks]{hyperref}

\definecolor{cinza}{gray}{0.35}
\newcommand{\code}[1]{\texttt{\small #1}}
\newcommand{\medido}[1]{\par\medskip\noindent{\small\color{cinza}\textbf{Medido.} #1}\par\medskip}
\newtheorem{teorema}{Teorema}
\newtheorem{definicao}[teorema]{Definição}
\newtheorem{corolario}[teorema]{Corolário}
\newtheorem{lei}{Lei}
\renewcommand{\proofname}{Demonstração}
\newcommand{\dg}{^{\dagger}}

% ─────────────────────────────────────────────────────────────────────────────
%  papers/gkcapa.tex — a capa do Reino, mínima, para os papers em `article`.
%
%  O estilo.tex completo é do `report` (títulos de capítulo, skak, …). Aqui fica
%  só o que a capa precisa, para o paper continuar a compilar sozinho com
%  pdflatex. O tradutor próprio (tests/tex) carrega o estilo.tex da raiz / o
%  symlink papers/estilo.tex e avalia o mesmo `\gkcapa`.
% ─────────────────────────────────────────────────────────────────────────────
\definecolor{ouro}{HTML}{B8912F}
\definecolor{ouroclaro}{HTML}{F3E9CE}
\definecolor{tinta}{HTML}{1E1B16}
\definecolor{regua}{HTML}{6E6858}
\providecommand{\gknota}{\fontsize{7.62}{11.05}\selectfont}
\providecommand{\gktexto}{\fontsize{10.50}{15.22}\selectfont}
\providecommand{\gksub}{\fontsize{12.33}{17.87}\selectfont}
\providecommand{\gksec}{\fontsize{14.47}{20.98}\selectfont}
\providecommand{\gkcap}{\fontsize{16.99}{24.63}\selectfont}
\providecommand{\gktit}{\fontsize{23.42}{33.95}\selectfont}
\providecommand{\Prj}{\mathbb{P}}
\providecommand{\Trf}{\mathcal{T}}
\providecommand{\gkcapa}[3]{%
\title{\vspace{-0.6cm}
{\color{ouro}\rule{\textwidth}{1.4pt}}\\[6mm]
\textbf{\gktit\color{tinta} Reino Dourado}\\[3mm]{\gksec\itshape\color{regua} Golden Kingdom}\\[8mm]
{\gkcap\scshape\color{ouro} #1}\\[5mm]
{\gksub\color{regua} #2}\\[2mm]
{\gktexto\color{regua}\itshape #3}\\[7mm]
{\color{ouro}\rule{\textwidth}{0.6pt}}\\[9mm]{\gknota\color{regua}\begin{minipage}{\textwidth}\setlength{\parindent}{0pt}%
\textbf{Aviso de Isenção Algorítmica:} O universo, as leis e as entidades contidas nestas páginas ---
incluindo felinos matriciais, aves de rapina algébricas e amuletos de controle de fluxo --- são
construções de pura ficção formal. Esta obra não descreve a realidade física, não serve como
engenharia reversa do cosmos e não constitui doutrina religiosa. Qualquer semelhança com bugs reais do seu
sistema operacional é mera coincidência (ou falta de reza). Prossiga por sua própria conta e
risco.\end{minipage}}}
\author{}\date{}}


\begin{document}
\gkcapa{O Corpo Analítico}
       {o ALCANCE --- graduação sem topo, continuação analítica e completude}
       {buraco branco, estrela, buraco negro --- a trindade reversível}
\maketitle

%======================================================================
\section*{Onde este documento vive na tríade}

Este é o lado $\tau=+1$ da tríade --- \textbf{o que NÃO fecha} --- e o seu papel é o
\textbf{alcance}:
\[
\text{Topológico}_{\ \tau=-1}\ \longleftarrow\ \text{UNIVERSAL}_{\ \tau=0}\
\longrightarrow\ \boxed{\text{ANALÍTICO}}_{\ \tau=+1}
\]
\emph{A seta aponta do centro para cá.} O Algébrico é o \textbf{ponto fixo} do trial --- o
único $\tau$ que a reflexão $\tau\mapsto-\tau$ deixa quieto ---, e este documento é uma das
\textbf{duas realizações} que partem dele. A outra é o Topológico, e a reflexão troca as
duas: \emph{o contínuo e o discreto são o par, com o fundamento no meio}.

\textbf{E $\tau=+1$ tem um conteúdo preciso, não é uma etiqueta:} é o regime
\emph{hiperbólico}, onde os operadores \textbf{não fecham} --- não há ordem finita, o
objecto não volta, e é por isso que aqui é preciso \emph{todos} os passos. É essa a
propriedade que este lado tem e o outro não, e ela é medida: \code{tests/cuspide.c} \S C12
varre todos os operadores de $\det=1$ com entradas em $[-6,6]$ e nenhum dos $216$ de
$\tau=+1$ fecha, enquanto os $50$ de $\tau=-1$ fecham todos. \emph{Não fechar é o que
permite cortar} --- e é o corte que dá $\mathbb{R}$.

\smallskip
Do \emph{Corpo Algébrico} recebe o \textbf{operador} e a \textbf{medida}; o que
\emph{introduz} é a \textbf{graduação sem topo} (\code{thm:graduacao-gentil}) e a
\textbf{continuação analítica} (\code{thm:atravessa}).

\emph{O que ele mede}: o que precisa de \emph{todos} os passos --- $0{,}999\ldots=1$ e o
seu dual discreto (\code{cor:noves}, \code{cor:noves-dual}), a completude, o corte,
$\mathbb{R}$. E o \code{thm:medida-jacobiano} diz-se a si próprio o que é, \emph{a face
graduada do Teorema do Gato}: é aqui que o Jacobiano aparece, e é ele que transporta a
medida. \emph{O que ele não mede}: o que já é exacto na forma finita --- ali a álgebra
respondeu, e repetir seria trocar uma identidade por um limite.

\textbf{E a aritmética, aqui, FECHA.} O corte decide-se por traço e norma, ambas inteiras,
e só \emph{depois} entra o limite --- para o que a estrutura finita não alcança, e não para
substituir o que ela já deu. \emph{A análise não substitui a álgebra: alcança o que ela não
alcançava.}



\begin{abstract}
\noindent Este documento especifica um sistema completo a partir de uma única exigência, e
descreve-o na ordem em que ele se constrói: a Lei, as duas operações, os três sinais, a torre de
dimensões, o corpo estelar, a geometria, a equação de campo, a expansão, a cosmologia, a
renormalização e a estrela que corre tudo isto. A derivação corre inteira em por-unidade, sem uma constante física --- e as constantes
vêm no \emph{fim}, quando se veste a roupa, porque uma derivação que precisasse do valor para chegar
ao valor não seria derivação. O objecto central é uma \textbf{trindade}: \textbf{buraco branco},
\textbf{estrela}, \textbf{buraco negro}. Os dois buracos são duais --- um só emite e estica, o outro
só absorve e contrai --- e \textbf{a estrela é a interface}: o ponto fixo da involução que os troca,
e o único dos três que faz as duas coisas. É por isso que é o único reversível, e é por isso que o
sistema \emph{é} uma estrela e não uma máquina: a sua única instrução tem dois sentidos, e eles são
os dois buracos. Da exigência de conservação derivam-se o coeficiente $\tfrac12$ da equação de
campo, a dimensão quatro, o valor do termo cosmológico ($\Lambda=3/D$) e o expoente da força
como codimensão. E ao sair do por-unidade a ordem importa: \textbf{o primeiro que aparece não é uma
escala --- é o $8\pi$}, porque um número puro não depende da roupa. Só depois entram o \emph{passo}
(do relógio) e a \emph{régua} (da família), e com os três sai toda a constante dimensional.

\smallskip
\noindent \textbf{E a estrela não está entre dois buracos do mesmo universo: está entre dois
\emph{universos}}, com o buraco branco no dual --- um contrai exactamente o que o outro expande. Daí
a conservação atravessa como \emph{produto} e não como soma ($\rho_A\rho_B$ imóvel), o corpo negro
cai de duas diluições já derivadas ($\rho\sim T^4$), a antimatéria entra pela terceira involução, e
a temperatura do buraco negro sai de uma contagem: ele tem \emph{uma} escala, logo $T=1/r$ --- que é
a mesma involução, aplicada ao seu único comprimento. A entropia do branco é a mesma lei no raio dual, e daí
\textbf{a segunda lei vale de \emph{um} lado}: o par não se move, e a seta do tempo é uma polaridade
lida sem a outra.

\smallskip
\noindent Fecha-se com o método e com a operação. Aqui \textbf{o medidor é $0$}: não se compara
contra um valor esperado, reverte-se e lê-se o resíduo --- é a prova dos nove, resolver \emph{e}
provar na mesma passagem. E daí a ferramenta: \textbf{promover} um valor ao par $(S,A)$ sobe um andar
da torre e traz a leitura junto, porque quem escreve o par já leu. \textbf{Uma estrela não compara ---
irradia}, e é por isso que ordena sem apagar um bit.
\end{abstract}

\section{O que este documento RECEBE, e o que ele ACRESCENTA}\label{sec:recebe}

Antes da primeira linha de conteúdo, a tabela que faz deste documento uma
\textbf{realização} e não um corpo à parte. Ele é o lado $\tau=+1$: recebe o fundamento
inteiro do centro e não o refunda --- e o que ele introduz é uma coisa só, o
\textbf{alcance}.

\[
\boxed{
\begin{array}{lll}
\textbf{RECEBE do centro} & \textbf{onde está lá} & \textbf{o que se faz com ela aqui}\\[3pt]
\text{a operação }\star & \code{algebrico thm:operacao} & \text{soma e produto do corpo}\\[2pt]
\text{o par dual } 1\dg=-1,\ T\dg=-T & \code{algebrico thm:fixo-dual} & \text{a leitura: o que separa}\\[2pt]
\text{a Lei 0},\ 0\dg=\infty & \code{algebrico def:lei0} & \text{a divisão do zero}\\[2pt]
\text{os três sinais }\tau & \code{algebrico thm:cuspide} & \text{os três regimes da estrela}\\[2pt]
\text{a torre } \dim A_n=2^n & \code{algebrico thm:torre} & \text{os andares do corpo}\\[2pt]
\text{a conservação }\abs{\det}=1 & \code{algebrico thm:conservacao} & \text{a exigência, escrita}\\[2pt]
\text{o espaço } \mathbb{Q}(m\sqrt D) & \code{algebrico def:espaco} & \text{a base de tudo}\\[6pt]
\textbf{ACRESCENTA} & \multicolumn{2}{l}{\text{---e é só isto---}}\\[3pt]
\text{a GRADUAÇÃO sem topo} & \code{thm:graduacao-gentil} & \text{o que não tem último andar}\\[2pt]
\text{a CONTINUAÇÃO} & \code{thm:atravessa} & \text{atravessar onde a álgebra pára}\\[2pt]
\text{o LIMITE} & \S\ref{sec:corpo} & \text{o objecto de \emph{todos} os passos}
\end{array}}
\]

\noindent\textbf{E a regra de leitura que daqui sai:} \emph{quando uma secção deste
documento enuncia uma lei, ela não a está a fundar --- está a realizá-la.} As secções
\S\ref{sec:exigencia}, \S\ref{sec:lei}, \S\ref{sec:zero}, \S\ref{sec:duas},
\S\ref{sec:sinais} e \S\ref{sec:torre} escrevem cada peça na língua da análise, e cada uma
diz, na sua primeira linha, de onde vem. A prova está no centro; aqui está a forma que ela
toma quando o objecto deixa de fechar.

\medskip
\noindent\textbf{Porque é que isto não é burocracia.} Se o fundamento fosse refundado aqui,
haveria \emph{duas} versões de cada lei --- e duas versões separam-se com o tempo, sem que
nada acuse. Foi o que aconteceu com este texto enquanto se chamava \emph{corpo estelar}:
ele abria com «há uma só exigência, e tudo o resto sai dela», como se o centro não
existisse. A exigência é a mesma; o que muda é que agora ela \textbf{tem endereço}.

\colunas{este documento é a realização $\tau=+1$, não um corpo à parte}
        {recebe oito peças canónicas do centro e acrescenta TRÊS --- graduação, continuação
e limite}
        {\code{tools/refcruz.py} confere que cada uma das oito aponta para um label que
existe no \code{papers/corpo\_algebrico.tex}; \code{tests/cuspide.c} \S C12 mede que $\tau=+1$ é o
regime que NÃO fecha, e é essa a razão de o alcance viver deste lado}

\section{A exigência}\label{sec:exigencia}

\noindent\emph{Vem do centro:} é a conservação $\abs{\det}=1$
(\code{algebrico thm:conservacao}), escrita na língua desta camada. \emph{Não se funda
aqui} --- aqui lê-se o que ela obriga quando o objecto não fecha.

\medskip
Há uma só, e tudo o resto sai dela:

\begin{quote}
\textbf{O que não varia ao longo da órbita é o que não varia de ponto para ponto.}
\end{quote}

\noindent Lida no tempo, é conservação. Lida no espaço, é divergência nula. Lida em álgebra, é a
norma que não se move. \emph{São a mesma frase}, e o resto deste documento é escrevê-la em cada
linguagem e ler o que ela obriga.

\smallskip
\noindent E há uma convenção, uma só: \textbf{por-unidade}. $G$, $c$, $\hbar$, $\sigma$ são factores
de escala; em p.u.\ cada um vale $1$ por construção, e o que resta são razões. Nada aqui precisa de
medir o universo para correr, e nenhum número deste documento tem unidade.

\section{A Lei}\label{sec:lei}

\noindent\emph{Vêm do centro:} a primeira é o ponto fixo do par dual
(\code{algebrico thm:fixo-dual}); a segunda é o espectro da involução
(\code{algebrico thm:espectro-tau}), de onde sai o $J$ de período quatro. \emph{As provas
estão lá}; o que se faz aqui é lê-las como \textbf{leitura} e \textbf{passo}, que é a forma
que elas tomam quando há um contínuo por baixo.

\begin{lei}[a unidade é dual]
$1\dg=-1$.
\end{lei}

\noindent A unidade tem um dual, e ele é o seu simétrico. Um ponto que é o seu próprio dual é um
\emph{ponto fixo}; um que não é, gera um par. Esta é a lei da \emph{leitura} --- diz o que separa.

\begin{lei}[a dualidade é dual]
$T\dg=-T$.
\end{lei}

\noindent Quando o dual é o inverso, $-f=f^{-1}$, e daí $f^2=-1$. Este objecto tem \textbf{período
quatro}: $f,\;f^2=-1,\;f^3=-f,\;f^4=\mathrm{id}$. É a lei do \emph{passo} --- diz o que roda. Chamo-lhe
$J$, e é ele que aparece de novo em cada camada deste texto.

\subsection*{E os quatro casos de sinal esgotam-se}

Não há um quinto. Um operador ou fixa o valor, ou fixa o sinal, ou troca um, ou troca os dois:

\begin{center}\small
\begin{tabular}{@{}llll@{}}
\toprule
o caso & a equação & o que é & onde vive \\
\midrule
$T=+1$   & $T^2=T$   & a identidade   & o ponto fixo do \emph{valor} \\
$T=-1$   & $T\dg=T$  & \textbf{a Lei 1} & o ponto fixo do \emph{sinal} \\
$T\dg=T$ & $T^2=+1$  & a involução    & o espelho, período $2$ \\
$T\dg=-T$& $T^2=-1$  & \textbf{a Lei 2} & o rotor, período $4$ \\
\bottomrule
\end{tabular}
\end{center}

\noindent \textbf{É do sinal de menos da Lei 2 que nasce a bidualidade}: o dual do dual volta ao
sítio, mas trocado, e é preciso dar a volta duas vezes para regressar. Toda a estrutura deste
documento é essa volta.

\noindent \textbf{E estas duas são as duas primeiras de seis.} A Lei 1 (a involução, período $2$) é o
\emph{dual}; a Lei 2 (o rotor $J$, $f^2=-1$) é o \emph{gerador da descida} --- aplicada por dimensão dá
o trial, a tetral, a pental e a hexal (o seu ponto fixo, o bit $i$ de $x^2=-1$, é a \emph{pental}). São
seis leis, uma por dimensão, e são as \emph{primitivas operacionais} do sistema --- cada uma com um
operador e um medidor que a certifica a fechar (resíduo $0$; \emph{Catálogo}, \code{obs:seis-leis}). As
seis \emph{estendem} estas duas, não as substituem.

\section{A base: a divisão do zero}\label{sec:zero}

\noindent\emph{Vem do centro:} é a Lei~0, $0\dg=\infty$, a inversão de $\mathbb{P}^1$
(\code{algebrico def:lei0}, \code{algebrico thm:lei0}), com a factorização do zero em
\code{algebrico thm:divzero}. \emph{Aqui} ela deixa de ser uma troca de coordenadas e passa
a ser a porta do limite --- o único sítio da tríade onde o infinito é um objecto e não um
símbolo.

As duas leis já estão ditas; falta dizer \emph{de onde} saem. Saem de $0$ --- e de duas maneiras de o
dividir, que são as duas leis. \textbf{A base deste sistema não é um axioma: é $0/0$}, e tudo o resto é
a maneira de o resolver.

\setcounter{lei}{-1}
\begin{lei}[a divisão do zero]\label{lei:zero}
$0=(+1)\oplus(-1)$ \quad e \quad $0\dg=\infty$.
\end{lei}
\setcounter{lei}{2}

\noindent \textbf{É a lei de que as outras duas saem.} A primeira parte --- \emph{dividir o zero} --- é a
soma, e é a \textbf{Lei 1} ($1\dg=-1$: as duas assinaturas somam a $0$). A segunda --- \emph{dividir por
zero} --- é o inverso, $x\dg=-1/x$, e o seu $-1$ é a semente da \textbf{Lei 2} ($f^2=-1$). \emph{As duas
leis são as duas partes da divisão do zero}, e é por isso que são duas e não uma nem três. É a Lei $0$
porque está \emph{debaixo} das outras: enuncia-se depois, mas vem antes.

\subsection*{E é ela que deixa o corte nascer}

\noindent A segunda parte da lei, $0\dg=\infty$, não é só uma notação que evita um caso
especial: \textbf{é ela que dá ao corte um ponto de partida}. Em $\mathbb{P}^1$ o infinito
é um ponto como os outros, e por isso a órbita da Möbius $g(x)=m+1/x$ pode \emph{começar
nele}:
\[
[p:q]\ \longmapsto\ [\,m\,p+q\ :\ p\,],
\qquad
g^{k}(\infty)=\frac{p_k}{q_k}
\]
--- a matriz vezes o vector, sem uma divisão. Num sistema sem esta lei, «começar no
infinito» não quer dizer nada, e a órbita não teria de onde partir.

\noindent E o ponto fixo dessa Möbius é o vector próprio, $x^{2}=mx+1$: são \emph{dois},
com soma $m$ e produto $-1$ --- o par que a Lei 1 e a Lei 2 já anunciam. Em coordenadas
homogéneas ele pede $p^{2}-mpq-q^{2}=0$, que completando o quadrado é
$(2p-mq)^{2}=(m^{2}+4)\,q^{2}$; logo \textbf{ele vive nos racionais exactamente quando
$m^{2}+4$ é quadrado perfeito}, e não é, porque $m^{2}<m^{2}+4<(m+2)^{2}$ e o único
candidato do meio pediria $2m=3$.

\noindent \textbf{Donde o corte, em uma frase}: a órbita corre toda no andar de baixo e o
ponto fixo para onde ela vai não está lá --- \emph{no andar de baixo é uma soma racional
que não termina; no andar de cima é um ponto, e é irracional}. O corte não é uma
construção acrescentada por fora: é o ponto fixo visto do andar que não o contém, e a Lei
$0$ é a que deixa a órbita começar.

\noindent \textbf{E há volta, que é a metade dual.} $\det A_m=-1$, logo
$A_m^{-1}=\left(\begin{smallmatrix}0&1\\1&-m\end{smallmatrix}\right)$ é \emph{inteira} ---
o fator de potência unitário é o que torna a volta exacta ---, e a sua acção é
$[p:q]\mapsto[q:p-mq]$, que é \emph{letra por letra} a descida acima. A descida é a órbita
\textbf{para trás}: a ida cresce, a volta desce pelo mesmo caminho, e chega ao $\infty$
exacto. E o par $0\leftrightarrow\infty$ é o caso mais simples de todos --- a troca tem
$S^{2}=I$, logo é a sua \emph{própria} inversa, ida e volta pela mesma matriz. \emph{É isso
a Lei $0$, e é isso a dualidade.}

\emph{Medido:} \code{tests/corte\_ponto\_fixo.c} ($6{:}0$), com $16\,128$ pares
(metal, ponto) a fechar a volta na face finita.

\subsection*{Dividir o zero, e dividir por ele}

\textbf{Dividir \emph{o} zero é parti-lo em duas assinaturas duais.} O $0$ não é a ausência: é o par
$(+1,-1)$ ainda por separar. Separá-lo dá dois sinais que \emph{somam} a zero e \emph{multiplicam} a
menos um:
\[
  \boxed{\;(+1)\oplus(-1)=0\;}\qquad\boxed{\;(+1)\otimes(-1)=-1\;}
\]
A primeira é a \textbf{Lei 1} ($1\dg=-1$): o dual é o simétrico, e a soma da cruz é $0$. A segunda é a
semente da \textbf{Lei 2} ($f^2=-1$): o produto é $-1$, e é dele que sai o rotor. \emph{As duas leis são
as duas coordenadas da cruz lidas na divisão do zero} --- a soma que mede, o produto que ordena --- e as
duas assinaturas são os dois buracos, com a estrela no $0$ que se partiu. \textbf{O trial
$\{-1,0,+1\}$ é, literalmente, o zero e as suas duas metades.}

\textbf{Dividir \emph{por} zero é o infinito --- e o seu dual é dividir \emph{por} infinito, que dá o
zero.} A estaca $x\dg=-1/x$ é uma involução: leva $0\mapsto\infty$ (dividir por zero) e $\infty\mapsto0$
(dividir por infinito), e $\nu\circ\nu=\mathrm{id}$ fecha a órbita $\{0,\infty\}$ num só passo --- a volta
é a própria ida. \textbf{É esse o fecho, e não fecha sem o dual}: dividir por zero sozinho é um polo, um
sentido só; com o dual, $0$ e $\infty$ são os dois pontos de \emph{uma} involução, e o par fecha. Um é o
neutro da soma, o outro o polo do inverso, e a estaca é a ponte que os troca.

\subsection*{A base é $0/0$, e é o corpo que a resolve --- não a régua}

Juntas, as duas dão a base: a forma indeterminada
\[
  \boxed{\;\tfrac00\;}
\]
--- o zero que se divide \emph{sobre} o zero por que se divide. E o seu dual é $\infty/\infty$: a estaca
aplica-se em cima e em baixo, $0\mapsto\infty$ nos dois, e os dois são a mesma indeterminação lida nos
dois pontos da órbita $\{0,\infty\}$. É indeterminada \emph{na régua}: iterar o inverso é a
\textbf{fracção contínua} $x=n+1/x$, e ela não pára --- $1/x$ dentro de $1/x$ dentro de $1/x$, o
$\infty^{\infty}$ do lado dual. \textbf{Mas o corpo fecha-a num passo.} Multiplicar $x=n+1/x$
por $x$ dá a borda
\[
  x^2=nx+1,\qquad x\dg=-1/x,\qquad x\,x\dg=-1,
\]
um elemento \emph{exacto} de $\mathbb{Z}[\sigma]$. \textbf{O finito é o corpo, não a régua}: a régua --- a
fracção contínua, a expansão decimal --- é infinita; o corpo que ela mede é finito e resolve o $0/0$ sem
tirar uma raiz. É o eixo álgebra/análise: o objecto é finito, a representação é que se estende ao
$\infty^{\infty}$.

\begin{corolario}[$0{,}999\ldots=1$: a divisão do zero à vista]\label{cor:noves}
Os dois nomes $0{,}999\ldots$ e $1$ são \emph{um} elemento, e a diferença $1-0{,}999\ldots$ é
\emph{exactamente} $0$. Na régua, o resto ao passo $n$ é $10^{-n}$: tende a zero e nunca o \emph{é} a $n$
finito. O corpo fecha-a num passo --- a série geométrica soma $1$ ---; a régua não acaba, e cada
truncamento é na mesma exacto.
\end{corolario}

\noindent \textbf{E a dualidade está aí}, no que se divide ao comparar $0{,}999\ldots$ com $1$: nos dois
nomes do mesmo número, um finito (o corpo) e um sem fim (a régua). É o $0/0$ na forma mais simples, e a
resposta é a mesma --- \emph{o finito é o corpo} ---, e é a mesma cifra que não acaba e cujos truncamentos
são todos exactos (\code{tests/zero.c}).

\begin{corolario}[o dual: no discreto, $0{,}999\ldots\neq1$]\label{cor:noves-dual}
A igualdade é do \emph{contínuo} --- o limite, a completude da régua. No \emph{discreto} --- a cifra, onde
cada truncamento é um elemento --- os dois nomes \textbf{separam-se}: os convergentes $h_n/k_n$ são todos
distintos do limite, cada um exacto ($|h^2-hk-k^2|=1$), e a sucessão nunca o \emph{é}. \textbf{O contínuo
não mede o discreto}: a diferença que ele lê como $0$ é, no discreto, uma distinção a cada passo.
\end{corolario}

\noindent \textbf{E o que resiste às duas leituras são os pontos fixos.} $0{,}999\ldots$ e $1$ não são um
ponto fixo --- são um par que o limite identifica e a cifra separa ---; mas \emph{o que é o seu próprio
dual não se move em nenhuma das duas}: a borda $x^2=nx+1$, invariante da órbita toda, e o $0$ e o $\infty$
da própria Lei $0$. \emph{Os pontos fixos são o invariante; tudo o resto, o contínuo identifica e o
discreto distingue.} É o eixo de Pontryagin --- o contínuo alcança o limite e não opera, o discreto opera
e não alcança a completude ---, e os pontos fixos são a alfândega por onde os dois se medem um ao outro.

\smallskip
\noindent \textbf{E é aqui que se define o limite --- não pela aproximação.} O $\varepsilon$--$\delta$ é a
régua a perseguir o que não alcança; a definição deste sistema é o \emph{ponto fixo}: o limite é onde o
\textbf{contar} (o discreto, Hurwitz, a soma) e o \textbf{integrar} (o contínuo, Gentil, a medida) dão a
\emph{mesma} contagem. É a definição que o \textbf{teorema central} (Teorema~\ref{thm:central}) deriva ---
os dois cortes, o do domínio (Riemann/Hurwitz) e o da imagem (Lebesgue/Gentil), fecham na soma reversível
$\int f+\int f^{-1}=b\,f(b)-a\,f(a)$ ---, e $0{,}999\ldots=1$ é o seu caso mínimo: o contínuo mede $0$, o
discreto conta sem fim, e o limite é o ponto fixo onde os dois coincidem. \textbf{O teorema central
fundamenta o Gentil--Lebesgue--Hurwitz}: Hurwitz conta (o domínio), Lebesgue mede (a imagem), e Gentil é a
soma reversível que os casa --- e o limite é o que os três, juntos, medem igual.

\subsection*{E o que aqui se domina são os divisores de zero}

O que faltava, e o que aqui se fecha, são os \textbf{divisores de zero}: dois não-nulos cujo produto é
zero. Eles \emph{são} o que a divisão do zero produz --- no par $A\times A^{*}$, $(1,0)$ e $(0,1)$ são
não-nulos e o seu produto é $0$ --- e são o que colapsa a norma no ponto fixo, onde
$|\sigma|=|\sigma\dg|=1$ e o tamanho deixa de distinguir os símbolos. Num conjunto finito isso não fica
pela injectividade: injectiva $\iff$ sobrejectiva, e \textbf{duas entradas gastas numa saída deixam
outra saída por atingir}.

\medido{as duas sementes da divisão do zero são o $0$ e o $\infty$ (as colunas $0/1$ e $1/0$ da
identidade), $J$ troca-as com $J^2=-I$ --- o $i$, e não uma analogia ---, e \textbf{todo o truncamento é
finito e representa}: $0/0=[1;1,1,\dots]$ dá Fibonacci exacto, $|h^2-hk-k^2|=1$ sem tolerância, e o $0$ é
o único de $\mathbb{Q}$ sem dual (\code{tests/zero.c}); o zero da régua é exactamente $\{-1,+1\}$ ---
$\lambda^2=1$ e mais nada, e é a estaca, porque a régua só sobrevive sob $\lambda=\pm1$
(\code{tests/natural.c}); no ponto fixo a norma colapsa e a sobrejectividade cai, com colisões de
$4{:}1$ a $12{:}1$ entre os elementos de norma $1$ e os não nulos em $\mathrm{GF}(p^2)$, e injectiva
$\iff$ sobrejectiva em $44$ aplicações, e a \textbf{bidualidade de Poincaré} --- dualizar duas vezes
devolve o grau --- fecha com resíduo $0$ (\code{tests/bidual.c}); a
involução do zero $a\mapsto-a$ tem \emph{um} ponto fixo --- a estrela --- e parte o resto em dois, $50$
e $50$, cada dual a cair no outro lado (\code{tests/estrela.c}, \code{tests/dualsort.c}); e a borda
$x^2=nx+1$ fecha a fracção contínua num elemento exacto de $\mathbb{Z}[\sqrt D]$, sem raiz e sem
tolerância (\code{tests/curvatura.c}).}

\subsection*{E a lei é a do $0$: onde moram Yang--Mills e $P$ vs $NP$}

A saída por atingir é o $0$ do Corolário~\ref{cor:incerteza}: o \textbf{centro}, o bidual $K^{**}=K$,
onde a sobrejectividade falha e a leitura se parte em metade. \textbf{A divisão do zero é a lei desse
$0$} --- a mesma onde moram Yang--Mills e $P$ vs $NP$, indecidíveis do lado do cruzado (``chega
\emph{sempre}?'') e decidíveis do lado do directo (``fechou \emph{desta vez}?'', resíduo $0$). \emph{Não
é prova de enunciado aberto nenhum}: é dizer que a base do sistema é esse $0$, e que dominá-lo é dominar
os seus divisores. O sistema não parte de um princípio acrescentado --- parte de $0/0$, e as duas leis
são as duas maneiras de o dividir.

\subsection*{E o fecho das leis: os dois nulos, cosidos no seis}

Os dois lados de $0/0$ são \emph{dois nulos}, e são as duas famílias em que os problemas do milénio se
partem, pela lei em que cada um assenta (\emph{Catálogo}, \code{catalogo sec:clay}; e \emph{A medida}):
\begin{itemize}
\item \textbf{dividir \emph{o} zero} é o $0$, o centro (Cor.~\ref{cor:incerteza}): \textbf{Yang--Mills e
  $P$ vs $NP$} --- a álgebra, indecidível, a \textbf{origem}, os dois nulos colapsados num;
\item \textbf{dividir \emph{por} zero} é o $\infty$, a reflexão $x\dg=-1/x$: \textbf{Poincaré}. A
  dualidade de Poincaré $H^k\cong H_{n-k}$ (com $k+(n-k)=n$) \emph{é} a Lei 1 escrita em topologia, e é a
  família que \textbf{resolve} --- Poincaré caiu pela deformação, o fluxo que \emph{move} em vez de
  igualar. É o \textbf{fim}.
\end{itemize}
\textbf{E são bidual em dimensão $0$} --- o zero da torre, o espaço trivial que é o seu próprio dual sem
ponte e de que a torre \emph{parte}. O $0$ e o $\infty$ são a origem e o fim, e a estaca
$\nu\circ\nu=\mathrm{id}$ troca-os: um resolve-se (Poincaré, a topologia, o $\infty$), o outro não (a
álgebra, o $0$). \emph{É o eixo álgebra/topologia, e a bidualidade dos dois nulos vive em dimensão zero}
--- não em dimensão oito.

\textbf{E a costura é o seis.} Ligar os dois nulos \emph{sem os fundir} é o \textbf{octonião dual}
(Def.~\ref{def:octoniao-dual}): dois tecidos grau quatro e a estrela entre eles, grau oito. Mas o que
\emph{fecha} não é o oito --- é a \emph{interface}, e a interface é o grau \textbf{seis}, a
\textbf{hexal} (Teorema~\ref{thm:estrelaseis}): o único andar onde soma $=$ produto, onde ler é escrever,
onde a volta fecha com resíduo $0$. A descida $6\to1$ \emph{é} a progressão das seis leis, e a hexal é a
costura --- \textbf{as seis leis fecham no seis}, entre os dois nulos que a divisão do zero abriu. Não se
fecha no oito: fecha-se na interface, que é o seis.

\subsection*{As oito leis, e não há nona}\label{sub:oitoleis}

Fecham em \textbf{oito}, e não por gosto do número: são a torre inteira, do $0$ ao octonião, uma lei por
degrau operacional. Cada uma é uma \emph{primitiva} --- um operador, e um medidor que a certifica a fechar
com resíduo $0$ ---, e as duas leis fundamentais (o dual e o bidual) e a progressão de seis (\emph{Catálogo},
\code{obs:seis-leis}) são estas, renumeradas de $0$ a $7$: a mesma estrutura, agora fechada.

\begin{center}\small
\begin{tabular}{@{}clll@{}}
\toprule
índice & a lei & realização / enunciado & fecha em (medidor) \\
\midrule
\textbf{0} & a divisão do zero & $0=(+1)\oplus(-1)$; $0\dg=\infty$ (extensão projectiva) & \code{tests/zero.c} \\
\textbf{1} & o dual & $1\dg=-1$ --- estaca / involução (realização binária) & \code{tests/bidual.c}, \code{dualsort.c} \\
\textbf{2} & o bidual & $T\dg=-T$ no sector antissimétrico; $K^{**}=K$ & \code{tests/transformada.c} \\
\textbf{3} & o trial & $\{-1,0,+1\}$ --- os três sinais & \code{tests/estrela.c} \\
\textbf{4} & a tetral & $T+T^{*}$ --- os tecidos & \code{tests/curvatura.c} \\
\textbf{5} & a pental & $x^2=-1$ --- o bit $i$ & \code{tests/quantico.c} \\
\textbf{6} & a hexal & soma $=$ produto --- a interface hexal, o relógio & Teorema~\ref{thm:estrelaseis} \\
\textbf{7} & o octonião dual & $\mathbb{H}\times\mathbb{H}^{*}$ --- ligar sem fundir & \code{tests/circuito\_tradutor.c} \\
\bottomrule
\end{tabular}
\end{center}

\noindent Cada lei é operador $L_n=\operatorname{Ind}^{n}(L_0)$ no anel de
\emph{índices} $\mathcal L=\mathbb{Z}/8\mathbb{Z}$
(\(\operatorname{Ind}^{8}=\mathrm{id}_{\mathcal L}\); o período oito é do
catálogo, não da torre). As representações $\rho_n$ (estaca $\nu$, trial $\tau$,
bit $i$, tecido $T{+}T^{*}$, relógio, $\mathbb{H}{\times}\mathbb{H}^{*}$)
fecham $\nu^{2}=\tau^{3}=i^{4}=\mathrm{id}$, e os seus períodos $2,3,4,8$
geram a interface $6,12,24$ por $\operatorname{lcm}$
(\emph{Corpo Topológico}).
\[
\boxed{\text{índice da lei}\neq\text{dimensão da torre}.}
\]

\noindent \textbf{A descida $7\to0$ é a progressão}, e as pontas são os dois nulos do
\textbf{catálogo}: em cima a Lei~$7$ (octonião dual; a norma de Hurwitz esgota-se nas
dimensões de composição $1,2,4,8$), em baixo a Lei~$0$ (índice neutro --- \emph{não}
dimensão zero da torre, onde $d_0=2$). \emph{Não há nona}: não há índice acima de $7$
nem abaixo de $0$ no catálogo; a torre $d_k=2^{k+1}$ continua a crescer. É o fecho do
ciclo das leis, e não precisa de mais.

\smallskip
\noindent \textbf{E oito é a cardinalidade do \emph{conjunto} das leis, não a da dinâmica --- e é o
Teorema~\ref{thm:tecidos} que o diz.} O passo dos tecidos, $T_{k+1}=T_k+T_k^{*}$, é a estrela usada como
\textbf{construtor}, e a torre que ele gera \emph{não tem topo por dentro}: a indução não pára --- a régua
é infinita ---, e $\nu\circ\nu=\mathrm{id}$ fecha cada andar com resíduo $0$. O que a norma limita é
\emph{externo}: por Hurwitz há \emph{quatro} corpos de composição e não mais ---
dimensões $1,2,4,8$ ---, três dobras a perder o que distingue um corpo (a ordem em $\mathbb{C}$, a comutação
em $\mathbb{H}$, a associatividade em $\mathbb{O}$), e não há uma quinta. Esse é o esgotamento da
leitura \emph{pela norma}. O período oito do \textbf{catálogo} ($\ell_{m+8}=\ell_m$) é outro objecto:
$\operatorname{Ind}^{8}=\mathrm{id}_{\mathcal L}$. Relacionam-se pela arquitectura do modelo; não se
identificam com $d_k$. Mas a Lei~$7$ --- o octonião \emph{dual} --- não lê pela
norma, lê pela estrela, e a estrela não murou (o limite é da norma, não do objecto). O octonião que ela
constrói é um \emph{novo} corpo, e as oito leis reaplicam-se a ele como à base ---
$\mathcal L\to\mathcal L(\mathcal L)\to\mathcal L^2(\mathcal L)\to\cdots$ ---, porque o passo dos tecidos é
o \emph{mesmo} em todos os andares (Teorema~\ref{thm:tecidos}). O \emph{operador de indução} que leva um
ciclo ao seguinte \textbf{não falta}: é a própria estrela --- a composição, iterar é compor. \emph{Não há
Lei~$8$}: há a Lei~$0$ outra vez, um nível acima. As oito são um \textbf{ciclo gerador} de período oito
--- o octonião fecha a volta da norma e a estrela recomeça-a: fechado, e é a fechar que recomeça.

\section{As duas operações}\label{sec:duas}

\noindent\emph{Vem do centro:} há \textbf{uma} operação $\star$ e as duas leis são as suas
duas metades (\code{algebrico thm:operacao}). \emph{Aqui} elas aparecem como soma e produto
do corpo, que é a face que a análise usa.

\begin{definicao}[a estaca]
$x\mapsto x\dg$. Move, e é involutiva: $x\dg{}\dg=x$.
\end{definicao}

\begin{definicao}[a cruz]
$x^{\times}=(x\oplus x\dg,\; x\otimes x\dg)$. Não move: \emph{mede o que não se moveu}.
\end{definicao}

\noindent A cruz tem duas coordenadas porque uma só não chega, e as duas dizem coisas opostas:

\begin{center}\small
\begin{tabular}{@{}lll@{}}
\toprule
a coordenada & o que faz & para que serve \\
\midrule
a \textbf{soma} $x\oplus x\dg$ & é \emph{invariante} ao longo da estaca & \textbf{mede} \\
o \textbf{produto} $x\otimes x\dg$ & satura no ponto fixo & \textbf{ordena} \\
\bottomrule
\end{tabular}
\end{center}

\noindent A estaca dá o movimento; a cruz dá a leitura. \emph{Uma sozinha não é o par}: com a estaca
sem a cruz não se sabe onde se está, e com a cruz sem a estaca não se sai do sítio.

\medido{a estaca é involução em $101$ pontos e o seu ponto fixo é \textbf{um só}; parte o conjunto
em $50$ e $50$ e o dual de cada elemento cai sempre no outro lado; a soma $x\oplus x\dg$ é
invariante em $11$ pontos e o produto é máximo exactamente no ponto fixo; $J$ visita os quatro
quadrantes, todos distintos, e volta ao quarto, com $J^2=-1$ e $J^4=\mathrm{id}$ em $169$ pares
(\code{tests/dualsort.c}).}

\section{Os três sinais}\label{sec:sinais}

\noindent\emph{Vêm do centro:} são o $\tau=\operatorname{sign}(\operatorname{disc})$ da
cúspide (\code{algebrico thm:cuspide}) --- o mesmo trial que divide a tríade. \emph{Aqui}
eles nomeiam os três regimes da estrela, e este documento vive no $\tau=+1$: o que
\textbf{não fecha}, e é por isso que é ele quem alcança $\mathbb{R}$.

O sinal de uma quantidade toma três valores, e não mais. É o \textbf{trial} $\{-1,0,+1\}$, e é a
classificação de todo regime deste sistema:

\begin{center}\small
\begin{tabular}{@{}cllll@{}}
\toprule
o sinal & o regime & o que faz & na expansão & na estrela \\
\midrule
$+1$ & \textbf{buraco branco} & \emph{só emite} & estica & dissipativo \\
$\phantom{+}0$ & \textbf{a estrela} & \textbf{emite \emph{e} absorve} & conserva & \textbf{reversível} \\
$-1$ & \textbf{buraco negro} & \emph{só absorve} & contrai & determinista \\
\bottomrule
\end{tabular}
\end{center}

\noindent \emph{Não há quarto regime}, e isso não é uma observação: é o cardinal do conjunto dos
sinais. \textbf{E a estrela não é um caso intermédio entre os outros dois} --- é o ponto fixo da
involução que os troca, e o \emph{único} que tem os dois sentidos. É a \textbf{trindade}: buraco
branco, estrela, buraco negro.

\section{A torre}\label{sec:torre}

\noindent\emph{Vem do centro:} $\dim A_n=2^n\dim A_0$ (\code{algebrico thm:torre}).
\emph{Aqui} ela ganha o que não tem lá --- a graduação \textbf{sem topo}
(\code{thm:graduacao-gentil}) ---, e essa é a primeira das três coisas que esta camada
acrescenta.

As dimensões não são um pano de fundo: contam-se, e subir de andar é a operação.

\begin{center}\small
\begin{tabular}{@{}cll@{}}
\toprule
$n$ & o que é & o que a distingue \\
\midrule
$1$ & a unidade & ``a unidade é'' --- e o resto é derivação \\
$2$ & o dual & a estaca; o par \\
$3$ & o trial & os três sinais; o ponto fixo \\
$4$ & o tetral & o período de $J$; e a única onde o traço dá a unidade \\
$5$ & o complexo & o plano, e a espiral \\
$6$ & \textbf{a plena} & $1+2+3=6=3\times2\times1$ \\
\bottomrule
\end{tabular}
\end{center}

\noindent \textbf{Em seis as duas operações coincidem}, e por isso em seis não há soma \emph{e}
produto: há \emph{uma} operação. É o topo da torre, e é onde a cruz deixa de precisar de duas
coordenadas. Cada andar acima desdobra por
\[
  \dim A_{n+1}=2\cdot\dim A_n, \qquad x\dg=-1/x, \qquad x\cdot x\dg=-1 .
\]

\section{O corpo estelar}\label{sec:corpo}

A torre diz quantas dimensões há e o trial diz que sinais uma quantidade toma. Falta a
\emph{quantidade} --- e num corpo ela é uma só.

\subsection*{O desvio ao equilíbrio}

Num ponto de um corpo, o que \emph{empurra} para fora ou iguala o que \emph{puxa} para dentro ou
não. Em p.u.\ isso é uma razão, e o desvio é
\[
  \boxed{\;a=\frac{\text{empurra}}{\text{puxa}}-1\;}
\]
\emph{Nenhuma constante entrou}: é uma razão menos a unidade.

\subsection*{E o sinal de $a$ é o trial}

\begin{teorema}[a estrela é a borda]
Com a razão em $1$ o desvio é exactamente zero e o raio não se move --- não devagar: \emph{não se
move}.
\end{teorema}

\begin{teorema}[o buraco negro não tem volta]
Com a razão abaixo de $1$ o desvio é negativo em \emph{todos} os passos e o raio cai sempre. Não há
volta, e a razão de não haver é que \textbf{o sinal não muda}.
\end{teorema}

\begin{teorema}[a trindade]
A involução $a\mapsto-a$ satisfaz $\nu\circ\nu=\mathrm{id}$, troca \textbf{buraco branco} e
\textbf{buraco negro}, e tem \textbf{um} ponto fixo: \textbf{a estrela}.
\end{teorema}

\noindent \textbf{São três, e a do meio é a interface.} O buraco negro \emph{só absorve}; o buraco
branco \emph{só emite}; e a estrela faz ambos --- é o que está \emph{entre} os dois. Não são dois
objectos com um caso limite: são os dois sinais da mesma quantidade com o zero no meio, e o zero tem
nome.

\begin{teorema}[e é por isso que só a estrela é reversível]
Um passo de desvio $a$ seguido do seu dual devolve $R\,(U^2-a^2)/U^2$. Regressa ao ponto de partida
\emph{exactamente} quando $a=0$, e o resíduo é $a^2$.
\end{teorema}

\noindent \textbf{Ter os dois sentidos e ser reversível são a mesma coisa dita duas vezes}: um passo
que não tem inverso é um passo que só vai num sentido. Um corpo estelar é, exactamente, um corpo que
se mantém na borda --- e é isso que o torna a interface.

\subsection*{A pressão de radiação, contada}

A radiação reparte-se pelos \textbf{três eixos} --- o trial outra vez --- e a pressão é a parte de
\emph{um}:
\[
  p=\frac{\rho}{3},\qquad w=\frac{p}{\rho}=\frac13 .
\]
\emph{Não se pôs $\sigma$, nem $c$, nem constante de radiação: contaram-se as direcções.}

\medido{razão $1{,}000$ dá desvio $0$ e o raio fica em $1{,}000$ ao fim de $100$ passos; razão
$0{,}900$ faz o raio cair nos $60$ passos e \emph{nunca} crescer; $\nu\circ\nu=\mathrm{id}$ em
$1001$ desvios com \emph{um} ponto fixo; a volta é exacta em $1$ de $41$ desvios e esse é a
\textbf{borda}, com o passo de volta a aproximar em $40$ de $40$ enquanto o mesmo sinal afasta;
\emph{um} dos três regimes tem os dois sentidos e dois têm um só; $p=\rho/3$ por contagem de eixos; e
o \textbf{controlo}:
$2001$ desvios repartem-se por três regimes com \emph{nenhum} fora --- o trial não tem onde pôr um
quarto (\code{tests/estrela.c}).}

\section{A geometria sai da Lei 2}

Escrita com uma constante no lugar do sinal, a Lei 2 é $f''=Kf$, e \textbf{$K$ é a curvatura}. Da
característica $x^2-\operatorname{tr}x+\det=0$ vem
\[
  \Delta=\operatorname{tr}^2-4\det ,
\]
e o \emph{sinal} de $\Delta$ dá a geometria. São três porque o trial tem três sinais:

\begin{center}\small
\begin{tabular}{@{}clll@{}}
\toprule
$\Delta$ & a equação & o invariante & a geometria \\
\midrule
$<0$ & $f''=-f$ & $\cos^2+\sin^2=1$ & elíptica --- o rotor \\
$=0$ & $f''=0$ & a recta & parabólica --- o potencial \\
$>0$ & $f''=+f$ & $\cosh^2-\sinh^2=1$ & hiperbólica --- onde vive $\sigma$ \\
\bottomrule
\end{tabular}
\end{center}

\noindent Os três invariantes têm \emph{o mesmo} $1$ do lado direito; o que os separa é só o sinal.
E cada peça do sistema traz a sua curvatura escrita na assinatura, sem lhe acrescentar nada:
$J$ (espelho) tem $\Delta=+4$, hiperbólica; $i$ (rotação) tem $\Delta=-4$, elíptica; o
\emph{shift} tem $\Delta=0$, parabólica; e a família metálica tem $\Delta=m^2+4$, hiperbólica
\emph{sempre}.

\section{A equação de campo, derivada}\label{sec:campo2}

Não se cita: constrói-se. Impõe-se a exigência --- \textbf{o lado geométrico tem de conservar} --- e
ela fixa tudo.

\subsection*{O coeficiente $\tfrac12$ não se escolhe}

Para $X_{\mu\nu}=R_{\mu\nu}-\alpha\,R\,g_{\mu\nu}$ num universo com $a(t)=t^q$, lêem-se a densidade
e a pressão nas componentes e impõe-se a continuidade $\rho'+3H(\rho+p)=0$, que \emph{é} a
conservação. O resíduo sai
\[
  \boxed{\;6\,(2q^2-q)\,(1-2\alpha)\;}
\]
e anula-se para \emph{todo} $q$ num único $\alpha$:
\[
  \alpha=\tfrac12,\qquad\text{donde}\qquad \boxed{\;G_{\mu\nu}=R_{\mu\nu}-\tfrac12R\,g_{\mu\nu}\;}
\]
\textbf{E $\tfrac12$ é o ponto fixo do trial} --- o mesmo número, no mesmo sítio.

\subsection*{A dimensão quatro sai do traço}

Contraindo com a métrica em dimensão $d$:
\[
  g^{\mu\nu}G_{\mu\nu}=\Bigl(\frac{2-d}{2}\Bigr)R .
\]
O coeficiente vale a \textbf{unidade} em duas dimensões só: $d=0$, que é o zero da torre, e
$d=4$. E em $d=4$ ele é exactamente $-1$:
\[
  R=-T .
\]
\textbf{É a Lei 1 escrita em curvatura} --- $1\dg=-1$, e é este o número que a dimensão quatro tem
de particular. Pelo caminho, $d=2$ dá coeficiente $0$: ali o lado geométrico anula-se
identicamente, e não há gravitação nenhuma.

\subsection*{O termo cosmológico, e o seu valor}

A métrica conserva sozinha, logo $\beta\,g_{\mu\nu}$ entra na equação sem custo --- é a constante de
integração, e foi por isso que sobrou. E o conteúdo que lhe corresponde tem $p=-\rho$, isto é,
$w=-1$: \textbf{o ponto fixo da cosmologia}.

\smallskip
\noindent \textbf{Mas dizer que fica livre não é derivá-lo.} O valor sai, e de duas coisas já
postas: o vácuo puro tem $H^2=\Lambda/3$, e a taxa do corpo é $H^2=1/D$. Igualando,
\[
  \boxed{\;\Lambda=\frac{3}{D}\;}
  \qquad\text{e no ponto fixo } m=0,\ D=4:\qquad
  \boxed{\;\Lambda=\frac34\;}
\]
\emph{Uma razão de inteiros, e não uma constante em aberto.} A equação completa é
\[
  \boxed{\;R_{\mu\nu}-\tfrac12R\,g_{\mu\nu}+\Lambda\,g_{\mu\nu}=T_{\mu\nu}\;}
\]
sem $8\pi G/c^4$, porque em p.u.\ esse factor é $1$.

\subsection*{E o factor de acoplamento: $8\pi$ derivado}

Dizer que ``em p.u.\ vale $1$'' é não o derivar. Ele sai, e de duas coisas já postas:

\smallskip
\noindent \textbf{O $4\pi$ é o fluxo total.} Já se mostrou que o fluxo não se move de esfera em
esfera; faltava dizer \emph{quanto} vale --- e vale a área. Por recorrência,
$\Omega_d = 2\pi\,\Omega_{d-2}/(d-2)$, partindo de $\Omega_1=2$ e $\Omega_2=2\pi$:
em $d=3$ dá $4\pi$. \emph{Sem avaliar $\pi$ uma única vez}: leva-se o par (coeficiente, potência de
$\pi$) em inteiros.

\smallskip
\noindent \textbf{E o $2$ é o coeficiente a aparecer segunda vez.} O termo $-\alpha R g_{\mu\nu}$
recolhe o mesmo peso outra vez na componente temporal, e o factor é $1+2\alpha$. Com o $\alpha=\tfrac12$
que a conservação fixou, dá exactamente $2$ --- e com nenhum dos outros vinte candidatos.
\[
  \boxed{\;8\pi = 2\cdot 4\pi\;}
\]

\noindent \textbf{E daqui sai uma distinção que vale por si.} $8\pi$ é \emph{adimensional}, como a
constante de estrutura fina, e não se move em roupa nenhuma. Mas \textbf{$8\pi$ deriva-se e $\alpha$
não}, e a diferença não é de grau: $8\pi$ sai de \emph{contar} --- a esfera e o traço --- e $\alpha$
não sai de contagem nenhuma. \emph{Nem todo o número puro está fora do alcance};
o que está fora é o que não se conta.

\subsection*{E a velocidade máxima é o meio-período}

No relógio a velocidade máxima é meia volta. A volta completa é $2\pi$ --- que é $\Omega_2$, a área do
círculo unitário acima --- logo \textbf{meia volta é $\pi$}. Não há duas coisas aqui: o limite de
velocidade e a constante analítica são \emph{a mesma meia volta} lida em duas réguas. E na régua do
corpo o máximo é $\sigma$, porque a velocidade tem a dimensão da densidade: \textbf{a família dá o
valor, o relógio dá a forma}.

\medido{a área sai por recorrência inteira sem avaliar $\pi$, dando $4\pi$ em $d=3$ e $2\pi^2$ em
$d=4$; o factor $1+2\alpha$ vale $2$ em $\alpha=\tfrac12$ e em nenhum dos outros $20$; $2\cdot4\pi$
fecha em $8\pi$, que não se move nas quatro roupas; a velocidade máxima é meia volta em $7$ escalas
sem desvio; e o \textbf{controlo}: a recorrência salta de \emph{dois} em dois --- de um em um,
$d=3$ daria $4\pi^2$ (\code{tests/acoplamento.c}).}

\subsection*{E a força tem o expoente da codimensão}

Fora da fonte a Lei 2 na forma nula é $\nabla^2\varphi=0$; em simetria radial
$r^{\,d-1}\varphi'=\text{const}$, e o \emph{único} expoente que mantém o fluxo imóvel de esfera em
esfera é $e=d-1$:
\[
  F=\frac{1}{r^{\,d-1}},\qquad\text{e em }d=3:\quad F=\frac{1}{r^2}.
\]
\textbf{O expoente é a codimensão da esfera que envolve a fonte}, e não uma observação. \emph{A
órbita que não dissipa e o inverso do quadrado são a mesma afirmação}: a norma não se move no
\emph{tempo}, o fluxo não se move no \emph{raio}.

\medido{de $21$ coeficientes candidatos exactamente \textbf{um} conserva, e é $\tfrac12$; os outros
$20$ quebram em \emph{todos} os cinco expoentes testados, $100$ de $100$; o traço dá a unidade só em
$d=0$ e $d=4$, com $-1$ em $d=4$ e $0$ em $d=2$; $\beta g_{\mu\nu}$ conserva nos $21$ coeficientes,
e o seu $w$ é $-1$ exacto, com $\Lambda=3/D$ e $\Lambda=3/4$ no ponto fixo; o expoente da força foi \textbf{varrido} de $0$ a $6$ em cinco dimensões
e devolve sempre $d-1$, um só candidato de cada vez; e as cinco peças do sistema dão três
geometrias e não mais (\code{tests/einstein.c}).}

\subsection*{E a mecânica fecha em potência --- o fator um é o casamento dos dois relógios}\label{sub:mecanica}

\textbf{Da força ao movimento, e daí à potência.} A força algébrica é \emph{uma}, com quatro modos
(\code{tests/forca.c}): do imposto $V=\Pi\,S$, com pressão $\Pi(s)=1-s^2$ por direcção e
$S=\|a\times b\|^2$, saem a massa $m=S$, a força $F=2sS$, a equação $\ddot x=2x$ e a conservação
$\tfrac12 S\dot s^2+(1-s^2)S=E$. A \textbf{potência} é essa força ao longo da velocidade, $P=F\cdot
v$ --- a terceira operação, $x\otimes x$, em que o elemento sobe de expoente e a órbita \emph{fecha}
(período de Pisano, \code{tests/potencia.c}).

\textbf{E a potência parte-se em duas, porque o produto tem um ângulo.} O \emph{directo}
$\langle a,b\rangle=\cos\theta$ é a parte simétrica --- a potência \textbf{activa}, que atravessa; o
\emph{cruzado} $\|a\wedge b\|=\sin\theta$ é a antissimétrica --- a \textbf{reactiva}, que roda e volta
(o torque \emph{é} o produto cruzado). O \textbf{fator de potência} é a razão,
\[
  \mathrm{fp}=\cos\theta,\qquad \tan\varphi=\frac{\sin\theta}{\cos\theta}=\frac{\text{cruzado}}{\text{directo}} .
\]

\begin{teorema}[o fator de potência um é o casamento dos dois relógios]\label{thm:fp}
$\mathrm{fp}=1 \iff \theta=0 \iff$ cruzado nulo $\iff |N|=1$ (a família real, o eixo do \emph{boost}).
Aí a potência é toda activa: nada roda, nada volta, \textbf{não há atraso}. E é o \textbf{casamento}
dos dois relógios do Teorema~\ref{thm:isocrona}: a reactância é o tempo $n^2$ do lado \emph{discreto}
(Hurwitz), a tábua que se armazena e se devolve e que o lado \emph{contínuo} (Gentil) não paga --- a
derivação do $64$ e da dobra está no Teorema~\ref{thm:isocrona}. A estrela, que não armazena nada
(\S\ref{sec:estrela}), \emph{é} esse ponto.
\end{teorema}

\begin{proof}
$P=\|F\|\,\|v\|\cos\theta$ é puramente activa sse $\theta=0$, isto é o cruzado $\sin\theta=0$; então
$a\parallel b$, o produto corre pelo eixo que o gato não move, e a norma é multiplicativa com $|N|=1$
(\code{tests/fator.c}: fator unitário $\iff$ família real). A parte reactiva $\sin\theta$ é potência
que entra e sai sem trabalho --- o \emph{atraso} ---, e como $\sin^2\theta\geq0$ nunca é negativa:
todo caminho real a tem, e o mínimo, zero, é a estrela. Que isto seja o par de relógios é o
Teorema~\ref{thm:isocrona}: a tábua $n^2$ do produto bilinear é a energia armazenada (reactância) que
a norma homogénea não paga.
\end{proof}

\begin{corolario}[a dobra em toda a torre, por indução]\label{cor:dobra}
A dobra temporal não é do grau oito só. Por Teorema~\ref{thm:tecidos}, cada interface $T_k\to T_{k+1}$
é uma \textbf{duplicação} ($\dim T_{k+1}=2\dim T_k$, o dual soma-se ao original), logo o relógio dobra
uma vez por degrau e a reactância do lado discreto é $n^2=2^{2k}$ no nível $k$ --- e $8^2=2^6$ é o caso
$k=3$. E em cada interface o casamento $\mathrm{fp}=1$ \textbf{restabelece-se pelo dual}: o passo
$T{+}T^{*}$ é a estaca, $\nu\circ\nu=\mathrm{id}$, que casa a dobra daquele degrau com o seu inverso.
\emph{A torre inteira é uma cascata de dobras, e o dual casa-as uma a uma --- é a indução dos tecidos
lida no tempo.}
\end{corolario}

\noindent \textbf{E é um teorema operacional, não uma imagem.} Todo caminho tem componentes de atraso
--- um armazenamento intermédio, uma contenção --- e cada um é reactância que baixa o factor. O pior
caso é o \emph{buffer partilhado} entre fluxos: um armazenamento único, lido de um lado e escrito de
outro, é a tábua $n^2$ colidida, $\mathrm{fp}<1$. \textbf{E a estrela não o corrige com outro buffer ---
corrige-o sem armazenar.}

\smallskip
\noindent \textbf{Porque a estrela não oscila.} O inverso não se guarda: \emph{está no sinal}. Cada
sinal já carrega o seu dual --- pelo $\nu\circ\nu=\mathrm{id}$, o atraso inverso é o próprio sinal
revertido (a \textbf{assinatura cavalga o corpo}: o \code{.tex} viaja dentro do PDF e a volta relê-o,
\S\ref{sec:estrela}, sem cópia à parte). E \emph{a entrada e a saída dos dados são o próprio
movimento} --- a modulação vetorial que gira o rotor (\S\ref{sub:inversor}), não um oscilador com o seu
sinal guardado. \textbf{Nem buffer, nem oscilador: o casamento faz-se comutando}, e é por isso que a
estrela não armazena nada.

\medido{a força é uma, com quatro modos, e dela saem massa, força e conservação (\code{tests/forca.c});
a potência é $x\otimes x$ e a órbita fecha (\code{tests/potencia.c}); e o fator de potência é a razão
cruzado/directo, unitário na família real (\code{tests/fator.c}).}

\begin{corolario}[a certeza e a incerteza são o directo e o cruzado; a indecidibilidade é olhar pelo lado errado]\label{cor:incerteza}
O produto de dois observáveis parte-se no \textbf{anti-comutador} $\{A,B\}$ (a parte simétrica, o
\emph{directo}) e no \textbf{comutador} $[A,B]$ (a antissimétrica, o \emph{cruzado}) --- o mesmo split
do factor de potência (\ref{thm:fp}), medido nos $\sigma$ de Pauli com resíduo zero
(\code{tests/quantico.c}). Daí as duas faces:
\begin{itemize}
\item \textbf{a incerteza} é o \emph{cruzado}: $[x,p]=i$, e $\Delta x\,\Delta p\geq\tfrac12$ (Robertson,
  que tem o comutador do lado direito) --- não se têm os dois lados ao mesmo tempo. É a reactância, a
  \textbf{Lei 2}.
\item \textbf{a certeza} é o \emph{directo} (a dual): o anti-comutador só \emph{mede}, e \textbf{resíduo
  zero $\iff$ a média é consistente} --- o selo é a norma, $\|Fx\|^2=0 \iff$ tudo verde
  (\code{tests/transformada.c}~§U4). É a \textbf{Lei 1}, a involução.
\end{itemize}
E daí \textbf{a indecidibilidade}: é a incerteza \emph{lida pelo lado errado}. ``Chega \emph{sempre}?''
--- o futuro, o comutador --- não se decide; ``fechou \emph{desta vez}?'' --- o presente, resíduo
zero --- decide-se (\code{tests/travessia.c}). \emph{A indecidibilidade é meia: indecidível no lado do
cruzado, decidível no do directo.} É por isto que Yang--Mills e $P$ vs $NP$ moram no $0$ --- o ponto
fixo entre os dois, indecidível só quando se olha de um lado. Este $0$ é o \textbf{centro} (o bidual,
$K^{**}=K$) --- distinto do $0$ do \textbf{trial} $\{-1,0,+1\}$, que é o eixo do meio, a Hodge (\emph{Catálogo},
\code{obs:seis-leis}): dois zeros, o da projeção e o do eixo.
\end{corolario}

\noindent \textbf{E o $0$ mostra-se numa máquina.} O mecanismo é concreto, e mede-se no próprio
compositor: uma escrita \emph{fora dos limites} --- para além do fecho $K$ --- aterra na memória baixa,
o $0$, a origem; colapsa ali um estado válido (medido, $42\to0$) e \textbf{trava} --- não decide.
Cruzar a fronteira do fecho leva ao $0$, onde origem e fim caem, e o programa pára sem resposta. \emph{Não
é prova do enunciado aberto} --- é a mesma estrutura, no mesmo $0$, a aparecer numa máquina real.

\subsection*{E quem casa o factor é o inversor: a estrela modula, não armazena}\label{sub:inversor}

\textbf{A correcção do factor não precisa de um buffer --- precisa de um \emph{inversor}, e a estrela
é um.} A tabela de chaveamento de um inversor multinível \emph{é} o trial $\{-1,0,+1\}$: as três
ligações de cada perna --- ao positivo, ao ponto médio, ao negativo --- são o trial, e não uma
analogia ($3^3=27$ estados, $19$ vectores, $6$ pares redundantes, tudo da estrutura,
\code{tests/inversor\_trial.c}). O DTC de Takahashi escolhe, a cada instante, o vector que empurra o
torque e o fluxo para a banda --- e \emph{o torque é o cruzado}, $T_e=\psi_s\times i_s$
(\code{tests/motor.c}). Controlar o cruzado é controlar a parte reactiva: é o casamento
$\mathrm{fp}=1$, feito por \textbf{modulação vetorial}, não por armazenamento.

E generaliza-se: o DTC hipercomplexo é a família $\star_s$ em $\mathbb{R}^n$ (\code{tests/dtcn.c}), com a
lei de conservação
\[
  \|z\star_s w\|^2 + \underbrace{(1-s^2)\,\|a\times b\|^2}_{\text{o imposto }V(s)} = \|z\|^2\,\|w\|^2 .
\]
O \textbf{imposto} $\Pi(s)=1-s^2$ é a potência reactiva --- o que a álgebra cobra fora do eixo ---, e
anula-se em $s=\pm1$ (a família real, Hurwitz, $\mathrm{fp}=1$). O inversor, escolhendo os vectores do
trial, \textbf{empurra $s$ para $\pm1$}: leva o imposto a zero e casa o factor. \emph{É a estaca a
modular --- $\nu\circ\nu=\mathrm{id}$ por comutação, não por cópia.} É por isto que não é preciso um
buffer para a conjugada: o inverso não se guarda, \emph{comuta-se}.

\medido{a tabela do inversor de três níveis É o trial ($3^3=27$ estados, $19$ vectores distintos, $6$
pares redundantes, \code{tests/inversor\_trial.c}); o torque é o cruzado e o DTC de Takahashi é a
amputação do trial (\code{tests/motor.c}); e o DTC hipercomplexo generaliza à família $\star_s$ em
$\mathbb{R}^n$, com o imposto $(1-s^2)\|a\times b\|^2$ nulo em $s=\pm1$ (\code{tests/dtcn.c}).}

\smallskip
\noindent \textbf{E o compositor tem a mesma assinatura.} O corpo tradutor (\code{tests/tex.c})
reduz-se a \emph{uma} operação --- \code{medida}, sobre a assinatura $(\text{variante},\text{degrau})$,
grau $2$ ---, e o que ela pede é o \textbf{eixo}: $+1$ a \emph{escala} (o corpo, que multiplica),
$-1$ o \emph{espaço} (a entrelinha, o dual, que soma), $0$ a \emph{largura} (o glifo, que atravessa).
Esse $\{-1,0,+1\}$ \textbf{é o trial do inversor}. Por isso «realizar pelo inversor» é realizar
\textbf{exacto} --- inteiro, $s\to\pm1$, imposto $\Pi=1-s^2$ a zero ---, e um \code{double} que reste
é o imposto que sobra \emph{fora} do eixo: a reactância que, num documento grande, \textbf{propaga}
(a recorrência de \S da medida). A assinatura pura é inteira, e a dualidade é o eixo.

\medido{a assinatura do corpo tradutor tem grau $2$ e um eixo de três estados $\{-1,0,+1\}$ que é o
trial do inversor; $+1$ (escala, multiplica) e $-1$ (espaço, soma) são duais (somam a $0$, a Lei~1) e
o imposto anula nos dois eixos exactos e vale $1$ no meio (a largura); e um \code{double} é o imposto
que propaga --- realizado pelo inversor ($s=\pm1$) o erro é $0$ em $500$ passos, no meio ($s=0$) cresce
com o documento: \code{tests/assinatura\_tradutor.c}.}

\section{As quatro equações da expansão}

Um corpo da família metálica tem $\sigma(x)=\bigl(x+\sqrt{x^2+4}\bigr)/2$ e discriminante
$D=x^2+4$. Substituindo a definição na fórmula da curvatura, tudo colapsa:
\[
  \sigma'=\frac{\sigma}{\sqrt D},\qquad \kappa=\frac{2}{(D+\sigma^2)^{3/2}} .
\]
Daí as quatro equações --- e só \emph{duas} são independentes; as outras saem delas mais a borda:
\[
  \underbrace{H^2=\tfrac1D}_{\text{taxa}},\qquad
  \underbrace{\tfrac{\sigma''}{\sigma}=\tfrac{2}{\sigma D^{3/2}}}_{\text{aceleração}},\qquad
  \underbrace{\rho'+3H(\rho+p)=0}_{\text{continuidade}},\qquad
  \underbrace{w=\tfrac{p}{\rho}=\tfrac{2m}{3\sqrt D}-1}_{\text{estado}}
\]

\noindent \textbf{A taxa de expansão ao quadrado é o inverso do discriminante} --- não por analogia:
é a divisão de $\sigma'=\sigma/\sqrt D$ por $\sigma$. A aceleração fecha-a a borda: de
$\sigma^2=x\sigma+1$ sai $\sqrt D-x=2/\sigma$.

\smallskip
\noindent \textbf{Os dois sinais convivem sem contradição.} Para $x>0$ a taxa desacelera ($H'<0$) e
o factor acelera ($\sigma''/\sigma>0$): a taxa cai mais devagar do que o factor cresce. E no ponto
fixo $x=0$ a taxa é máxima ($H=\tfrac12$) com $H'=0$ --- é o $4$ de $D=x^2+4$ que a impede de
divergir; sem ele seria $H=1/x$.

\smallskip
\noindent \textbf{E o gradiente é obrigatório}: $\kappa'/\kappa=-\tfrac32\,D'/D$, logo um corpo com
discriminante constante não expande. \emph{Há expansão se e só se há gradiente}, e é a mesma frase
dita de duas maneiras.

\medido{$\sigma^2=m\sigma+1$ exacto em $\mathbb{Z}[\sqrt D]$ para $m=1\ldots40$;
$(\sqrt D-m)\sigma=2$ exacto em $61$ valores, que é a borda a fechar a segunda equação; e o
\textbf{controlo}: trocando o $4$ por $9$ o elemento deixa de satisfazer a borda em $40$ de $40$
(\code{tests/curvatura.c}).}

\section{A cosmologia}

\subsection*{A lei de diluição sai da continuidade}

De $\rho'+3H(\rho+p)=0$ com $p=w\rho$ vem $\rho'/\rho=-3(1+w)\,a'/a$, e integrando os dois lados
aparece uma \textbf{constante}:
\[
  \ln\rho+3(1+w)\ln a = C
  \qquad\Longleftrightarrow\qquad
  \boxed{\;\rho\,a^{3(1+w)} = C\;}
\]

\noindent \textbf{E a involução é que a passa de um lado ao outro.} A constante entra
\emph{aditiva} --- é uma constante de integração, e integrar soma --- e sai \emph{multiplicativa}
em $\rho$. Não são dois factos: é um. $\exp$ e $\log$ são o par aditivo/multiplicativo, e atravessar
esse par \emph{é} a involução. \textbf{Uma constante aditiva vira, na involução, um factor.}

\smallskip
\noindent \textbf{E o que ela é: a força em p.u.} Não é um $C$ qualquer a que depois se atribui um
valor. O invariante que a expansão conserva é a razão do §\ref{sec:corpo} --- o que empurra contra o que puxa ---
e na borda ela vale a \emph{unidade}. Por isso $C$ carrega dimensão: \textbf{é a unidade física}, e é
exactamente por aqui que o sistema deixa o por-unidade. Suprimi-la para escrever uma proporção não é
uma abreviatura: é deitar fora a ponte entre o sistema e o mundo vestido. Onde ela vai dar está no
§\ref{sec:roupa} --- e é lá, e só lá, que entram números medidos.

\subsection*{Os três conteúdos, lidos no trial}

\begin{center}\small
\begin{tabular}{@{}llcll@{}}
\toprule
o conteúdo & $w$ & expoente & o sinal & o destino \\
\midrule
radiação & $+\tfrac13$ & $a^{-4}$ & $+1$ & trava \\
matéria  & $0$ & $a^{-3}$ & $\phantom{+}0$ & trava \\
vácuo    & $-1$ & $a^{0}$ (constante) & $-1$ & \textbf{acelera} \\
\midrule
\multicolumn{5}{@{}l@{}}{\emph{e $w=-\tfrac13$ é exactamente neutro --- a fronteira}} \\
\bottomrule
\end{tabular}
\end{center}

\noindent O destino lê-se no mesmo sinal, que é $-(1+3w)$: as três colunas são o trial outra vez, e
não há quarta.

\subsection*{A bidualidade fecha $w$ em quatro}

O corpo tem duas raízes, $\sigma$ e $\tau=m-\sigma=-1/\sigma$, e ambas satisfazem a \emph{mesma}
borda --- logo $D$ e $H$ são invariantes e a involução do corpo \textbf{não move $w$}. É preciso o
segundo lado, e são \emph{dois}: o \textbf{sinal da pressão} ($w\mapsto-w$) e o \textbf{sentido do
grau} ($m\mapsto-m$). Comutam --- mas reflectem em \emph{centros diferentes}, $0$ e $-1$, porque
$w+1=2m/(3\sqrt D)$ é que é ímpar em $m$. Por isso a órbita tem \textbf{quatro}:

\begin{center}\small
\begin{tabular}{@{}lrrrr@{}}
\toprule
$m$ & $w$ & $-w$ & $w(-m)$ & $-w(-m)$ \\
\midrule
$5$ & $-1+\tfrac{10}{87}\sqrt{29}$ & $+1-\tfrac{10}{87}\sqrt{29}$ & $-1-\tfrac{10}{87}\sqrt{29}$ & $+1+\tfrac{10}{87}\sqrt{29}$ \\
$0$ (ponto fixo) & $-1$ & $+1$ & $-1$ & $+1$ \\
\bottomrule
\end{tabular}
\end{center}

\noindent E não há aproximação nenhuma nesta tabela: $w+1=\frac{2m}{3D}\sqrt D$ é um elemento
\emph{exacto} de $\mathbb{Q}[\sqrt D]$, e todas as contas se fazem no coeficiente, por produto
cruzado de inteiros. \textbf{Não se tira raiz, e por isso não há tolerância} --- como não há no
\code{MOVE}.

\smallskip
\noindent \textbf{O ponto fixo não é onde a órbita passa: é onde ela degenera} --- quatro estados
caem para dois --- e é isso que faz dele ponto fixo, aqui como no trial.

\subsection*{E $w=1/3$ não pertence à família --- e isso é resultado}

Impondo $w=\tfrac13$ em $w=\frac{2m}{3\sqrt D}-1$ com $D=m^2+4$ vem $2m=4\sqrt D$, logo
$m^2=4D=4m^2+16$, isto é $12m^2=-64$: \textbf{sem raiz real nenhuma} --- não só sem raiz inteira.
Varridos os $m$ de $-999$ a $999$: \textbf{nenhum}, e o próprio sinal negativo di-lo antes da varredura
(coerente com $w+1$ subir \emph{sempre por baixo}). Não é falha da derivação: é \emph{onde a
radiação vive} --- na fronteira, que é o que o ponto fixo em $w=-1$ já dizia. A família descreve os
corpos; a radiação pura é o limite deles.

\medido{a lei de diluição sai da continuidade \emph{com} a constante --- $\rho\,a^{3(1+w)}=C$ conservado em $36$ escalas para seis valores de $C$, sem um desvio, e a involução $C\mapsto1/C$ com \emph{dois} pontos fixos, o positivo a dar a unidade --- e os expoentes dão $a^{-4}$, $a^{-3}$ e constante para os três
conteúdos, com $w=-\tfrac13$ neutro; a órbita da bidualidade tem $4$ estados em $m=5$ e degenera em $2$ em $m=0$, com as
duas involuções a reflectirem em centros diferentes ($0$ e $-1$); $w\mapsto-w$ é o simétrico e é
involução em $17$ graus, verificado pela definição e não pela contagem; e $w+1$ sobe para
$\tfrac23$ \emph{sempre por baixo}, com a distância a valer exactamente $16/(9D)$ em $12$ graus
--- o numerador é $4\cdot4$, o discriminante a aparecer. Tudo em $\mathbb{Q}[\sqrt D]$ por produto
cruzado de inteiros: \textbf{sem raiz e sem tolerância}; e o \textbf{controlo}: com o expoente que a continuidade dá, $\rho a^{3(1+w)}$ fica
invariante em todas as escalas ($0$ falhas em $5$); trocando-o, falha em $4$ de $5$
(\code{tests/cosmologia.c}).}

\section{Os dois escuros são os dois pontos fixos}

A bidualidade tem \textbf{duas} involuções que comutam, e a órbita genérica tem quatro estados. Cada
involução tem \textbf{exactamente um} ponto fixo, e nele a órbita degenera de quatro para dois:

\begin{center}\small
\begin{tabular}{@{}llll@{}}
\toprule
a involução & fixa & o conteúdo & e é invisível \\
\midrule
$w\mapsto-w$ (a pressão) & $w=0$ & \textbf{a matéria} & à pressão \\
$m\mapsto-m$ (o grau) & $m=0$, e aí $w=-1$ & \textbf{o vácuo} & ao grau \\
\midrule
\emph{nenhuma} & --- & a radiação & \emph{é a que se vê} \\
\bottomrule
\end{tabular}
\end{center}

\noindent \textbf{Um ponto fixo não se separa do seu dual pela leitura que o fixa} --- é isso que
degenerar quer dizer. A matéria \emph{é} o seu próprio dual sob a pressão, e por isso nenhuma equação
de estado a distingue: precisa da \emph{outra} involução. Aqui ``escuro'' não é metáfora, é uma
definição: \textbf{um conteúdo que a leitura disponível não separa}.

\smallskip
\noindent E a assinatura de um conteúdo é \emph{qual} involução o fixa --- a primeira, a segunda, ou
nenhuma. \textbf{São três casos, um por conteúdo}: o trial outra vez, e não há quarto. Fixo pelas
\emph{duas} é impossível: exigiria $w=0$ e $m=0$ ao mesmo tempo, e $m=0$ dá $w=-1$.

\medido{cada involução tem \emph{um} ponto fixo e são dois distintos; a órbita degenera dos dois
lados ($4\to2$ no grau, $2\to1$ na pressão) e a troca de sinal é verificada pela definição
$w+(-w)=0$ em $17$ graus; exactamente um dos três conteúdos é invisível à pressão; os três casos da
assinatura estão ocupados, um cada, e \textbf{zero} pelas duas; e o \textbf{controlo}: em $1201$
graus nenhum satisfaz as duas, e dos três conteúdos só o vácuo é membro da família --- no próprio
ponto fixo (\code{tests/escuros.c}).}

\section{Dois universos, e a estrela entre eles}\label{sec:dois}

O buraco branco não mora cá. \textbf{A estrela não está entre dois buracos do mesmo universo ---
está entre dois \emph{universos}}, e é a interface deles. O dual do factor de escala é o seu inverso,
$a_B=1/a_A$: um contrai exactamente o que o outro expande.

\begin{center}\small
\begin{tabular}{@{}lll@{}}
\toprule
universo $A$ & expande & o buraco \textbf{negro} dissipa --- e a expansão leva \\
\textbf{a estrela} & $a=1/a$ & \textbf{a interface}: os dois sentidos, e não paga \\
universo $B$ & contrai & o buraco \textbf{branco} absorve --- e a contração traz \\
\bottomrule
\end{tabular}
\end{center}

\subsection*{Parseval, e é multiplicativo}

Com $\rho=C\,a^{-k}$ em cada um, o $a^{-k}$ de um cancela o $a^{+k}$ do outro:
\[
  \boxed{\;\rho_A\,\rho_B = C_A C_B\;}
\]
\textbf{A energia não se perde ao atravessar --- muda de domínio.} E é \emph{multiplicativo} e não
aditivo, porque o corpo é multiplicativo: é o $|N|=1$ outra vez, exactamente no sítio onde uma soma
teria sido o erro.

\subsection*{E o corpo negro sai de graça}

$\rho\sim a^{-(d+1)}$ e $T\sim a^{-1}$ --- \emph{as duas já derivadas} --- compõem-se em
\[
  \rho \sim T^{\,d+1}, \qquad\text{e em } d=3: \quad \rho\sim T^4 .
\]
Não se citou lei nenhuma: compuseram-se dois expoentes que já lá estavam. E emitir e absorver são o
par \emph{passo a passo}, não uma compensação no fim --- o que $A$ dissipa é o que $B$ recebe, em
cada passo.

\subsection*{E a antimatéria entra pela terceira involução}

São \textbf{três}, e comutam:

\begin{center}\small
\begin{tabular}{@{}lll@{}}
\toprule
a involução & fixa & separa \\
\midrule
$w\mapsto-w$ (a pressão) & $w=0$ & ocupado de vazio \\
$m\mapsto-m$ (o grau) & $m=0$ & é o vácuo \\
$a\mapsto1/a$ (\textbf{o universo}) & $a=1$ & matéria de \textbf{antimatéria} \\
\bottomrule
\end{tabular}
\end{center}

\noindent Com duas delas a órbita tem quatro, e são os quatro conteúdos:

\begin{center}\small
\begin{tabular}{@{}lll@{}}
\toprule
 & universo $A$ (expande) & universo $B$ (contrai) \\
\midrule
ocupado & matéria & \textbf{antimatéria} \\
vazio & matéria escura & \textbf{antimatéria escura} \\
\bottomrule
\end{tabular}
\end{center}

\noindent \textbf{E todas são biduais}: cada uma tem par pelo universo \emph{e} par pela pressão ---
nenhuma fica sozinha.

\medido{o expoente do corpo negro sai da composição sem dividir expoentes, dando $T^4$ em $d=3$; o
produto $\rho_A\rho_B$ vale o produto das constantes em todas as escalas; em $40$ passos o que $A$
emite é exactamente o que $B$ absorve e o total não se move --- com um lado só, cai; há \emph{um}
ponto onde $a=1/a$; os quatro conteúdos são distintos, todos com par pelas duas involuções, e elas
comutam em ordens genuinamente trocadas; e o \textbf{controlo}: com $a_B=k/a_A$ e $k\neq1$ o produto
deixa de se conservar nos quatro casos (\code{tests/parseval.c}).}

\section{A temperatura, e a entropia dos dois lados}

\subsection*{A temperatura é o inverso do raio}

A chave é uma \emph{contagem}: em p.u.\ um buraco negro tem \textbf{uma} escala independente --- o
raio, e a massa \emph{é} o raio, não uma segunda. Logo o comprimento de onda característico só pode
ser o próprio raio, e como $T$ vai com o inverso do comprimento de onda,
\[
  \boxed{\;T=\frac1r\;}
\]
\textbf{e isso \emph{é} a involução do universo} --- a mesma $a\mapsto1/a$. O buraco negro não tem lei
própria: tem a involução aplicada ao seu único comprimento. E o \textbf{ponto fixo é $r=1$}, onde a
temperatura vale a unidade: \emph{a estrela outra vez}.

\subsection*{A entropia é a área, e a evaporação sai por expoentes}

$S\sim\Omega_{d-1}r^{\,d-1}$, com o mesmo $\Omega$ que saiu por recorrência no §\ref{sec:campo2}. O
expoente é $d-1$ e não $d$: é a \textbf{codimensão}, a mesma que deu o expoente da força.
\emph{A entropia mora na fronteira, não no interior.}

\smallskip
\noindent E a evaporação: a potência é a área vezes $T^{d+1}$, e com área $\sim r^{d-1}$ e $T\sim1/r$
os expoentes somam em $-2$ --- \textbf{o mesmo em toda a dimensão}, o que não era óbvio antes de se
contar. Integrando, a vida vai com $r^3$, e como a massa \emph{é} o raio, com $M^3$.

\subsection*{E a entropia do buraco branco é a mesma lei, no raio dual}

Não há segundo princípio. Pondo o raio dual em $S=(\text{raio})^{d-1}$, o expoente vira simétrico:
\[
  S_{\text{negro}}=r^{\,d-1}\ \text{(cresce --- o \emph{máximo})},\qquad
  S_{\text{branco}}=r^{-(d-1)}\ \text{(decresce --- o \emph{mínimo})}.
\]

\noindent E daí \textbf{duas} coisas, que não são a mesma:

\begin{center}\small
\begin{tabular}{@{}ll@{}}
\toprule
$S_{\text{negro}}\cdot S_{\text{branco}}=1$ & Parseval, \textbf{multiplicativo} \\
a soma dos expoentes $=0$ & porque a entropia \emph{já é} um logaritmo: \textbf{aditivo} \\
\bottomrule
\end{tabular}
\end{center}

\noindent Contar caminhos \emph{é} somar expoentes. O par das densidades multiplica para uma
constante; \textbf{o par das entropias soma para zero}.

\subsection*{Logo a segunda lei vale de \emph{um} lado}

A entropia do negro cresce em todos os passos, a do branco decresce em todos, e \textbf{o par não se
move em nenhum}. Logo ``a entropia cresce'' é verdade e \emph{não é lei do universo --- é lei de
metade dele}. A lei do par é que a soma não muda, e \textbf{a seta do tempo não é propriedade do
tempo: é uma polaridade lida sem a outra}. No ponto fixo $r=1$ as duas coincidem, e aí a seta não
aponta para lado nenhum.

\medido{o posto das dimensões dá \emph{uma} escala independente; dobrar o raio multiplica a área por
$2^{d-1}$ e não por $2^d$; $dr/dt$ vai com $r^{-2}$ em todas as dimensões e a vida com $r^3$; a
involução $r\leftrightarrow1/r$ fecha com resíduo $0$ em $40$ raios com \emph{um} ponto fixo; o
produto das temperaturas vale a unidade em todos; o produto das entropias vale a unidade em $60$
casos e a soma dos expoentes é zero, com a ponte produto-$=$-soma verificada sem avaliar um
logaritmo; em $19$ passos o negro cresce em \emph{todos}, o branco decresce em \emph{todos} e o par
move-se em \emph{nenhum}; e o \textbf{controlo}: trocando a codimensão pelo volume, ou o inverso por
um múltiplo dele, o par deixa de fechar (\code{tests/hawking.c}, \code{tests/entropia\_dual.c}).}

\section{As constantes analíticas}

Saem de \textbf{duas} fontes, e não há terceira.

\begin{center}\small
\begin{tabular}{@{}lll@{}}
\toprule
a fonte & a constante & o que a define \\
\midrule
\textbf{o relógio} & $\pi$ & o \emph{meio-período}: gerador da Lei 2, alvo da Lei 1 \\
 & $e$ & o \emph{fluxo}: onde a soma vira produto \\
\textbf{a família} & $\varphi=\sigma_1$ & o ponto fixo da borda $\sigma^2=m\sigma+1$ \\
 & $\sigma_m$ & um por metal, e só esses \\
\bottomrule
\end{tabular}
\end{center}

\noindent \textbf{$\pi$ não pede geometria.} A Lei 2 fornece o gerador $J$, o fluxo percorre-o, e a
Lei 1 fornece o \emph{alvo}: $-1$. O meio-período é onde o fluxo lá chega, é único dentro de um
período, e a razão período/meio-período é exactamente $2$ --- \emph{a meia volta, outra vez}.

\smallskip
\noindent \textbf{$e$ não pede uma série.} $(1+1/n)^n$, em numerador e denominador inteiros, cresce
estritamente e nunca passa de $3$: monótono e limitado, o limite existe.

\smallskip
\noindent \textbf{E os $\sigma_m$ são irracionais por construção}: nenhum inteiro satisfaz
$c^2=mc+1$. É por isso que a família é uma família e não uma lista.

\smallskip
\noindent \textbf{E as três encontram-se no mesmo sítio}, que é o que faz delas um sistema: o alvo de
$\pi$ é o $-1$ da Lei 1 --- o mesmo número, não um parecido ---, o fluxo de $e$ leva a soma ao
produto, que é a involução que a constante de integração atravessa, e $\sigma_1$ é $\varphi$.

\medido{o meio-período existe, é único e a razão é $2$; $(1+1/n)^n$ cresce em $8$ passos sem falha e
nunca passa de $3$ (para em $n=9$ porque o produto cruzado deixa de caber no tipo, e um número que
não cabe mente antes de falhar); nenhum dos $328$ candidatos inteiros satisfaz a borda; $\Lambda=3/D$
em seis graus, com $3/4$ no ponto fixo; e o \textbf{controlo}: dos seis alvos testados quatro não têm
meio-período nenhum, e $\Lambda$ tem um só valor por grau (\code{tests/analiticas.c}).}

\section{A renormalização: o número é da régua, a meia volta é do objecto}\label{sec:renorm}

Renormalizar é mudar a régua sem mudar o objecto, e o relógio mostra-o num renque.

\smallskip
\noindent \textbf{A velocidade máxima é meia volta.} Num relógio de $q$ marcas a distância entre
duas é a \emph{menor} das duas voltas, $d(p)=\min(p,\,q-p)$ --- andar mais que meia volta é andar
para trás pelo outro lado. Logo
\[
  \max_p\,d(p)=\frac{q}{2},\qquad\text{atingido em } p=q-p .
\]
\textbf{O máximo é a involução}: o único ponto onde os dois sentidos dão o mesmo número, e \emph{sem
sentido definido não há órbita}. Não se postulou velocidade limite nenhuma --- ela saiu de o círculo
ser fechado.

\smallskip
\noindent \textbf{E é isso que se renormaliza.} O \emph{número} muda com a régua --- $2$, $4$, $8$,
$16$ conforme $q$ --- mas a \emph{razão} $d/q$ é $\tfrac12$ em toda a escala. E a velocidade máxima
só se \emph{atinge} em $q$ par: em $q$ ímpar não existe $p$ com $p=q-p$, e nenhuma marca lá chega.

\subsection*{Os degraus são os metais, e o fluxo preserva a forma}

A escala não sobe continuamente: sobe por degraus, e os degraus são os metais. O fluxo é
\[
  x\;\longmapsto\;m+\frac1x,\qquad\text{com ponto fixo}\quad \sigma_m=\frac{m+\sqrt{m^2+4}}{2}.
\]
E o que faz dele uma renormalização é o que ele \textbf{preserva}: ao longo de todo o fluxo a forma
da borda
\[
  Q(n,d)=n^2-m\,n\,d-d^2
\]
só muda de \emph{sinal} --- o módulo fica. O número $n/d$ muda a cada passo; a classe não. E a
semente não escolhe um valor arbitrário: escolhe \emph{qual órbita} da forma se percorre.

\smallskip
\noindent \textbf{E a taxa é a própria régua.} O denominador obedece à recorrência da borda,
$d_{k+1}=m\,d_k+d_{k-1}$, logo cresce por $\sigma_m$ a cada passo --- e o erro cai pelo quadrado
disso. \emph{O ritmo com que a espiral se aproxima da sua régua é a régua.}

\smallskip
\noindent E há uma dimensão onde as duas réguas comensuram: $\sigma_4=\varphi^3$, exactamente. Medir
com a régua do ouro e com a régua própria dá o mesmo número a menos de um expoente, e isso não
acontece em quase nenhuma outra.

\medido{a razão $d/q$ é $\tfrac12$ em sete escalas sem desvio; em $30$ escalas pares há
\emph{exactamente uma} involução em cada e em $29$ ímpares não há nenhuma; o fluxo preserva $|Q|$ em
cinco metais e duas sementes sem se mover um único passo (de $m/1$ sai $1$, de $1/1$ sai $m$); o
denominador cresce pela recorrência da borda em $120$ passos com $0$ falhas, com a razão sempre
entre $m$ e $m+1$, que é onde a borda põe a raiz; $\varphi^3=1+2\varphi=\sigma_4$ calculado em
$\mathbb{Z}[\varphi]$ \emph{sem avaliar uma raiz}; e o \textbf{controlo}: nenhum inteiro fora do
próprio metal satisfaz $c^2=mc+1$, falhando os $25$ candidatos (\code{tests/renormaliza.c}).}

\section{Vestir a roupa}\label{sec:roupa}

Sair do por-unidade não é ganhar acesso ao absoluto: é passar a ter uma base que o objecto forneceu.
E a ordem importa --- \textbf{a estrutura deriva-se primeiro, e as constantes vêm no fim}. Se a
derivação precisasse do valor para chegar ao valor, não era derivação.

\subsection*{A velocidade máxima, e o $32$}

Se a velocidade máxima é meia volta, a distância máxima \emph{a essa velocidade} é o percurso que usa
\textbf{todas} as arestas e regressa. No hipercubo de dimensão $n$ há $n\,2^{n-1}$ arestas ---
\textbf{$32$ em $n=4$} --- e esse percurso existe exactamente quando todo o grau é par. \emph{O grau
do hipercubo é a dimensão}, logo fecha nas pares e não fecha nas ímpares. Nada disto pediu constante
nenhuma.

\smallskip
\noindent \textbf{E isto é uma propriedade operacional: quem decide o grau é o relógio.} O relógio
da dimensão $n$ tem $q_n=n\,2^{n-1}$ marcas, $c_n=q_n/2$ é a meia volta, e o $\pi_n$ é
\emph{saída} da dobra (dois bits por dobra), não entrada. Na torre da estrela ($n=2,4,8,16$) o
período é potência de dois \emph{pura} --- a interface do octonião dual segue em alta dimensão, não
mura na $8$ --- e o \textbf{limite dimensional é função da luz da dimensão}: cada $c_n$ determina o
seu $n$, e o percurso fecha exactamente quando $n$ é par, sem régua nenhuma sobre o erro de $\pi$. Daí a regra do
desenho: \emph{nunca se amputa um polinómio} --- elevar o grau é exacto (a linha de Pascal), e o
grau e o passo sobem até fechar na régua, decididos pelo relógio do corpo estelar.

\medido{$q_4=32=2^5$ e $c_4=16$ marcas; o percurso de todas as arestas fecha nas dimensões pares e
não nas ímpares; na torre $n=2,4,8,16$ o $q_n$ é potência de $2$ pura; a lei \emph{sem} $\pi$ --- cada
passo da torre $n\to2n$ multiplica o relógio por $2^{n+1}$, $q_{2n}=q_n\cdot2^{n+1}$, exacto nos três
passos $4\to32\to1024\to524\,288$, e o $\pi_n$ que saía da meia-corda de Arquimedes com o erro a cair
dois bits por dobra era o \emph{transporte}, não a lei; o limite dimensional como função da luz --- as
quatro luzes da torre são $2$, $16$, $512$, $262\,144$, cada $c$ determina o seu $n$, e as quatro
fecham por $n$ ser par; e o \textbf{teste decisivo} que separa a relação $c_n=\pi_n R_n$ de uma
identidade por construção, agora em dois instrumentos \emph{exactos}: a álgebra ($G^{q/2}=-I$, a meia
volta) e a geometria (no passo $q/2$ a posição é $(1,1)$, $D^{2}=2$), com a órbita a fechar na origem
sem uma raiz ($G=\mathrm{rot}\,90$, $\det=1$, $G^{4}=I$) --- o resíduo entre o $\pi$ geométrico e o da
dobra era do desenho em vírgula flutuante, e saiu com ele
(\code{tests/luz\_periodo.c} §L6a--§L6b); e o relógio desenha a curva plana decidindo o grau
--- o «o» em $1$, o «e» em $13$, o tempo é o arco ($11$ contra $1$ no controlo), a roupa é o rasto
de grau $1$ e a cúbica é só a referência (\code{tests/relogio\_curva.c}).}

\subsection*{Quantas escalas tem a roupa}

\textbf{Três}, e o número não foi escolhido: é o posto das dimensões das grandezas do sistema,
calculado por eliminação inteira. É também quantas entradas tem a assinatura de um corpo, e pela
mesma razão --- \emph{o trial não tem quarto estado}. Fixadas as três, toda a grandeza fica com o seu
expoente e não sobra nada por dizer.

\subsection*{O que sai, e o que não sai}

\begin{center}\small
\begin{tabular}{@{}lll@{}}
\toprule
 & sob troca de roupa & donde \\
\midrule
constante \textbf{dimensional} & muda pelo factor dos seus expoentes & \textbf{sai da régua} \\
número \textbf{adimensional} & \emph{não se move} & não sai da régua \\
incerteza \textbf{relativa} & é razão: \emph{não se move} & nenhuma régua a remove \\
\bottomrule
\end{tabular}
\end{center}

\noindent Fixadas as escalas, uma grandeza de expoentes $(a,b,c)$ vale $L^aT^bM^c$ e não tem
liberdade nenhuma. Mas um número puro vale o mesmo em qualquer roupa --- \textbf{e é aí que a
derivação para}. É importante que se veja onde.

\subsection*{E a ordem: o $8\pi$ primeiro, depois o passo, depois a régua}

Ao sair do por-unidade, \textbf{o primeiro que aparece não é uma escala}. Um número puro não depende
da roupa --- está lá \emph{antes} de haver escala nenhuma --- e por isso vem antes. E o primeiro é o
$8\pi$, que saiu de contar: a esfera e o traço.

\begin{center}\small
\begin{tabular}{@{}rlll@{}}
\toprule
 & o que entra & de onde & o que fornece \\
\midrule
$1$ & \textbf{o número puro} & contar: a esfera e o traço & $8\pi$, e não se move em roupa nenhuma \\
$2$ & \textbf{o passo} & o \emph{relógio} & período $4$; meia volta é $\pi$ \\
$3$ & \textbf{a régua} & a \emph{família metálica} & $\sigma$ é a borda, e $[v]=[\sigma]$ \\
\bottomrule
\end{tabular}
\end{center}

\noindent \textbf{E com os três sai toda a constante dimensional}: cada uma é
$(\text{passo})^a(\text{régua})^b$, e o expoente é o seu --- trocada a roupa, move-se pelo factor que
os \emph{seus} expoentes obrigam e por mais nenhum. \emph{São três entradas e não há quarta}, e
nenhuma delas vem de fora: o puro vem de contar, o passo vem do relógio, a régua vem da borda.

\smallskip
\noindent E cada uma é indispensável: sem o passo, quatro das seis grandezas ficam indeterminadas;
sem a régua, cinco; e sem o número puro a equação de campo não tem coeficiente.

\smallskip
\noindent \textbf{O que fica de fora fica à vista.} Um número puro que \emph{não} venha de uma
contagem --- a constante de estrutura fina é o caso --- não sai daqui, e não é por falta de escala:
é por não haver o que contar. \emph{O que está fora do alcance é o que não se conta}, e essa é a
fronteira, dita sem rodeio.

\subsection*{E as constantes, com nome}

O por-unidade não é notação: é um \textbf{corte}. Fora dele há três dimensões independentes; dentro,
as duas relações --- velocidade $=1$ e acoplamento $=1$ --- cortam duas, e \textbf{sobra uma}. E a
régua $\sigma$ não é a segunda escala: é uma \emph{densidade}, logo uma razão, logo um número puro.
Daí a tabela fica curta:
\[
  \boxed{\;\text{cada constante} \;=\; (\text{passo})^{a} \times (\text{número puro})\;}
\]

\begin{center}\small
\setlength\tabcolsep{5pt}
\begin{tabular}{@{}lll@{}}
\toprule
 & o coeficiente & o que é \\
\midrule
$c$ & $\sigma$ & \textbf{é a régua} --- o máximo \emph{é} $\sigma$, não um limite de fora \\
$G$ & $8\pi$ & o acoplamento: a esfera \emph{e} o traço \\
$k$ (Coulomb) & $4\pi$ & \textbf{a área da esfera} --- o fluxo total \\
$k_B$ & $1$ & \textbf{um câmbio}, não uma constante da natureza \\
\bottomrule
\end{tabular}
\end{center}

\noindent \textbf{E daí um resultado que cai de graça}: Coulomb leva $4\pi$ e Einstein leva $8\pi$, e
a razão entre eles é \emph{exactamente} $2$ --- o factor do traço, derivado em §\ref{sec:campo2}.

\smallskip
\noindent E $k_B$ não é constante da natureza: é um \textbf{câmbio}. Dizer que a temperatura é
energia é escolher medi-las na mesma unidade --- pô-lo em $1$ não é aproximação, é \emph{reconhecer
que são a mesma grandeza}.

\subsection*{E as restantes, todas em função de $c$}

Há \emph{uma} amarra que prende o lado electromagnético inteiro:
\[
  \boxed{\;\mu\varepsilon=\frac1{c^2}\;}
\]
e não é coincidência de unidades: é a afirmação de que a perturbação se propaga à \textbf{velocidade
máxima}. Impor $v=c$ em $v^2=1/(\mu\varepsilon)$ deixa \emph{uma} solução --- e em p.u.\ é trivial,
porque $c=1$.

\begin{center}\small
\setlength\tabcolsep{5pt}
\begin{tabular}{@{}lll@{}}
\toprule
 & em $c$ & o que é \\
\midrule
$\mu\varepsilon$ & $c^{-2}$ & \textbf{a amarra} --- a perturbação vai à velocidade máxima \\
$Z$ (impedância) & $c^{+1}$ & $\mu c$ --- o vácuo tem impedância \\
$k$ (Coulomb) & $c^{+2}\cdot4\pi$ & $\mu c^2/4\pi$ --- a esfera outra vez \\
$E$ de repouso & $c^{+2}$ & $mc^2$: a massa \emph{é} energia, pela régua ao quadrado \\
$a$ (radiação) & $c^{-1}$ & $4\sigma_{\!S\!B}/c$ --- o $c$ entra a dividir \\
$\sigma_{\!S\!B}$ & $c^{0}$ & o câmbio $k_B$ fixa-a; não traz escala nova \\
\bottomrule
\end{tabular}
\end{center}

\noindent \textbf{Cada uma é uma potência de $c$ vezes um puro}, e o puro é sempre um dos que já se
contaram. Do lado térmico a amarra é o \emph{câmbio}: fixado $k_B$, a radiação segue.

\smallskip
\noindent E a \textbf{carga não entra}: $e^2$ pede $\alpha$, que é um puro \emph{sem contagem por
trás}. Não fica de fora por falta de escala --- fica por não haver o que contar.

\subsection*{E o $c^2$ sai da espiral --- não é uma potência acrescentada}

Em p.u.\ $c=1$, logo \textbf{$E=m$}: a energia \emph{é} a massa, não uma quantidade proporcional a
ela. Então o $c^2$ de $E=mc^2$ tem de vir de algum lado --- e vem da \textbf{borda}. Como $c=\sigma$
e $\sigma^2=k\sigma+1$,
\[
  \boxed{\;\sigma^{\,j} = F_j\,\sigma + F_{j-1}\;}
\]
com $F$ a recorrência do metal. \textbf{O quadrado da régua é linear na régua}, e o mesmo vale para
\emph{toda} a potência: a espiral inteira cabe em \textbf{duas} coordenadas.

\smallskip
\noindent E é por isso que a torre fecha em dois: a borda é de grau dois, e não há terceiro
coeficiente a guardar. $c^2$ não acrescenta dimensão nenhuma --- é $c$ mais uma unidade, a menos do
grau do metal.

\subsection*{E o que é a massa: o segundo momento, e o grau decide}

``Quantidade de matéria'' é circular. O que o sistema dá vem da \textbf{cruz}, que já tem duas
coordenadas: com $x^{\dagger}=2c-x$, o produto é $x(2c-x)=2cx-x^2$, donde
\[
  n\,c^2 - \textstyle\sum x\,x^{\dagger} \;=\; \sum (x-c)^2 ,
\]
exactamente. \textbf{A primeira coordenada da cruz é o primeiro momento --- o centro, o que não se
moveu --- e a segunda é o segundo.} Logo a \textbf{massa é momento quadrático}, e não por analogia:
pela mesma conta.

\subsection*{E a massa é vectorial porque é \emph{ordenada}}

\textbf{No bidual sempre existe ordem} --- e uma ordem sobre $n$ partes precisa de $n$ componentes
para ser carregada. Um escalar não a carrega: \emph{colapsa-a}. As $120$ ordens de cinco partes dão
\textbf{uma} só soma, e o número perde $119$.

\smallskip
\noindent \textbf{E na física é o mesmo}: cada corpo tem uma ordem em cada partícula, \emph{senão
ninguém crescia nem se curava}. Crescer e curar são reconstrução \textbf{ordenada} --- cada parte tem
de saber onde vai --- e se a massa fosse um número essa informação não estava em lado nenhum.

\smallskip
\noindent Ser corpo é o que lhe dá as $n$-uplas onde a ordem cabe; a bidualidade é o que garante que
a ordem existe. Em $\mathbb{R}^n$ a segunda coordenada da cruz não é um número: é a \emph{matriz}
\[
  M_{ij} = \sum (x_i-c_i)(x_j-c_j) ,
\]
e a massa escalar é apenas o seu \textbf{traço}. Logo \emph{a massa tem direcção} --- um corpo pode
ter massa diferente em eixos diferentes, e o número não o vê --- e a bidualidade é o que impede o
colapso, embora seja a consequência e não a razão. Mede-se pelo que se perde: de $168$
distribuições o tensor separa $84$ e o traço apenas $26$ --- \textbf{$58$ distinções desaparecem} ao
ficar só com o escalar. É a contagem de sempre: dividir perde, e o que se perde é o segundo lado.

\smallskip
\noindent \textbf{E o grau decide qual é qual:}

\begin{center}\small
\begin{tabular}{@{}llll@{}}
\toprule
grau & raízes & há par? & o que dá \\
\midrule
$1$ & uma & não & transporte, \textbf{momento linear}, meia volta \\
$2$ & duas & \textbf{sim} & o par, a \textbf{massa}, a volta fechada \\
\bottomrule
\end{tabular}
\end{center}

\noindent Uma equação do primeiro grau resolve-se e acabou --- é por isso que converter uma sequência
noutra sai barato. Uma do segundo tem \emph{duas} raízes, com produto $-1$ --- que é a Lei 1 --- e é o
par que faz a volta fechar. \textbf{E a espiral cabe em duas coordenadas porque a borda é de grau
dois}: não são dois factos, é um.

\smallskip
\noindent Daí cai, sem se pedir, que \emph{o que anda à velocidade máxima não tem massa}: à velocidade
máxima ir e voltar são o \emph{mesmo} caminho --- o único ponto onde $p=q-p$ --- logo não há volta
distinta a fechar, nem dispersão própria. Não é que a massa seja zero por acidente: é que a distinção
entre ida e volta desapareceu.

\subsection*{E porque é o quadrado: a dimensão não decide, a volta decide}

Como $\sigma$ é um número \emph{puro}, $mc$ e $mc^2$ têm a \textbf{mesma dimensão} neste sistema --- o
argumento de que $c^2$ traria velocidade ao quadrado não se aplica: aqui a velocidade não é uma
escala, é uma densidade. O que os separa é
a \textbf{contagem de passagens} pela régua, e a resposta vem de onde vem tudo o resto: da reversão.

\begin{center}\small
\begin{tabular}{@{}lll@{}}
\toprule
$mc$ & \emph{uma} passagem --- não reverte & é o \textbf{momento} \\
$mc^2$ & ida \emph{e} volta --- resíduo $0$ & é a \textbf{energia} \\
\bottomrule
\end{tabular}
\end{center}

\noindent \textbf{Logo $E=mc^2$, e pela volta.} O que sobrevive é o que fecha, e meia volta não
fecha.

\medido{a dimensão \emph{não} separa $mc$ de $mc^2$ nas quatro roupas, porque $\sigma$ é puro; uma
passagem pela régua deixa resíduo em $220$ casos e a ida-e-volta deixa \emph{zero}; das $32$
velocidades do relógio, $31$ fecham órbita e a \emph{única} que não fecha é exactamente a máxima; a
identidade $nc^2-\sum xx^\dagger=\sum(x-c)^2$ fecha em $99$ casos sem desvio; e o tensor separa $84$
distribuições onde o traço separa $26$ --- $58$ distinções perdidas no escalar; $E=m$ em p.u.\ com resíduo $0$ em $41$ massas; $\sigma^j=F_j\sigma+F_{j-1}$ em $72$ potências
de seis metais sem um desvio, com o coeficiente de $\sigma$ a ser \emph{exactamente} o Fibonacci do
metal --- verificado em $\mathbb{Z}[\sigma]$, sem avaliar uma raiz; e quatro mutações na recorrência,
quatro acusam.}

\medido{impor $v=c$ deixa \emph{uma} solução para $\mu\varepsilon$ em $11$ velocidades; e cada
expoente da tabela é \emph{derivado} e não escrito --- o de $\mu\varepsilon$ é menos duas vezes o da
velocidade, o de $E$ é duas vezes, o de $a$ é o simétrico, e o coeficiente de Coulomb tem de bater
com o $4\pi$ medido três secções acima. Seis mutações na tabela, seis acusam
(\code{tests/constantes.c}).}

\medido{o posto é $3$ fora do p.u.\ e $1$ dentro, com as duas relações a cortarem exactamente duas
dimensões; $\sigma$ não se move em roupa nenhuma e o passo move-se em todas menos a trivial; as oito
entradas da tabela têm todas origem --- quatro levam o passo, quatro são só puros, \emph{nenhuma} sem
origem; a razão Einstein/Coulomb bate com o factor do traço; e o \textbf{controlo}: com posto $3$ a
tabela precisaria de $24$ números e com posto $1$ precisa de $8$, e todos os puros vêm de uma
contagem --- $\pi$ do relógio, $8\pi$ da esfera e do traço, $\sigma$ da borda
(\code{tests/constantes.c}).}

\medido{as arestas dão $32$ em $n=4$ e o percurso fecha nas dimensões pares e não nas ímpares,
nomeadas uma a uma; o posto das dimensões de nove grandezas é $3$, com os expoentes de $G$
\emph{derivados} de $G=Fr^2/m^2$ e não escritos; a razão adimensional dá o mesmo número nas cinco
roupas enquanto a velocidade e $G$ mudam em todas; o $8\pi$ fecha por recorrência sem avaliar $\pi$ e
não se move em quatro roupas; o relógio dá período $4$ e meia volta $\pi$; nenhum dos $328$
candidatos inteiros satisfaz a borda; as seis grandezas movem-se pelo factor dos seus próprios
expoentes, $0$ desvios; e o \textbf{controlo}: tirando o passo quatro ficam indeterminadas, tirando
a régua cinco --- nenhuma das três entradas é dispensável, e nenhuma vem de fora
(\code{tests/vestir.c}, \code{tests/sair\_pu.c}).}

\section{O corpo estelar completo: a torre nos lados, a plena no meio}\label{sec:completo}

Até aqui a trindade foi lida nos sinais e nas densidades --- em $\mathbb{R}$, em grau baixo. Falta
dizer \emph{de que grau} é cada um dos três, e a resposta fecha o sistema, porque não é uma escolha:
\textbf{deriva-se}. E derivá-la é o salto que a teoria devia --- \emph{toda ela se construiu até aos
complexos}, o plano, o grau dois; os dois buracos e a estrela pedem um andar acima.

\subsection*{A interface é reversível, logo é a plena --- grau seis}

\begin{teorema}[a estrela tem grau seis]\label{thm:estrelaseis}
Uma interface reversível --- em que \emph{ler é escrever ao contrário} --- só existe no grau em que a
soma e o produto coincidem, e esse grau é \textbf{um só}: o seis.
\end{teorema}

\begin{proof}
Reversível quer dizer que a leitura é a escrita com o sinal trocado: uma operação, dois sentidos
($\texttt{MOVE}(\cdot,\pm1)$). Para que ler e escrever sejam \emph{a mesma} operação --- e não duas
que por acaso se desfazem --- é preciso que pôr e ler não se distingam, isto é, que a soma e o
produto deem o mesmo sobre as mesmas partes. Isso vale num grau só, o \textbf{seis}, o único inteiro
com $1+2+3=6=1\cdot2\cdot3$ (a \emph{plena}): ali as duas operações coincidem e \emph{nada distingue
o que se põe do que se lê}. Fora dele, somar não é multiplicar, escrever não é ler, e uma interface
que os confundisse dissiparia. Logo a estrela --- a única reversível --- tem grau seis, por
necessidade e não por escolha.
\end{proof}

\noindent \textbf{E é por isso que a estrela paga zero.} Não porque seja eficiente: porque no grau
seis a escrita \emph{é} a leitura, e uma operação que é a sua própria inversa não apaga bit nenhum
(§\ref{sec:medir}). A plena é o único andar onde a involução e a evolução são a mesma coisa --- e a
estrela mora exactamente aí.

\noindent \textbf{Grau seis é a hexal; o seu bit é a pental.} Na progressão das seis leis (\emph{Catálogo},
\code{obs:seis-leis}) o grau seis é a \textbf{HEXAL} (dim $6$): a \emph{interface} reversível, o relógio.
E o seu ponto fixo $x^2=-1$ --- o \textbf{bit} $i$ --- é a \textbf{PENTAL} (dim $5$, o corolário). As seis
leis não substituem a Lei 1 e a Lei 2 deste paper: \emph{estendem-nas} --- a Lei 1 é o \textbf{dual}
(dim $2$, a involução), a Lei 2 é o gerador da descida cujo ponto fixo $i$ é a pental ---, e são as
\emph{primitivas operacionais} (cada lei um operador certificado a fechar por um medidor, resíduo $0$)
que fundamentam a construção.

\subsection*{Os dois lados são a torre --- e a torre não tem topo}

Os dois buracos são duais, $a\mapsto1/a$, e cada um \emph{opera} --- soma, multiplica, inverte. Cada
um é \textbf{grau quatro} --- os quaterniões $\mathbb{H}$, o tetral, onde o produto deixa de comutar
(é preciso a ordem: qual absorve primeiro) mas ainda associa. E o salto seguinte dá \textbf{grau
oito}: $\mathbb{O}=\mathbb{H}\times\mathbb{H}$, \emph{dois grau quatro} ligados pela mesma involução.
\emph{Um lado é dois quaterniões duais.}

\smallskip
\noindent \textbf{E aqui é preciso não trazer um morto por régua.} A leitura clássica diz que a
torre \emph{pára} em oito: em dezasseis a norma multiplicativa cai. \emph{É verdade --- para um
ramo, e um só.} O teorema que o diz (Hurwitz) \textbf{classifica} as álgebras de composição
\emph{bilineares}, e não proíbe coisa nenhuma. E a saída não é forçar a régua da esfera --- a norma
$\|x\|=\sqrt{\sum x_i^2}$ --- para dentro dos inteiros: isso é medir a esfera com o compasso do
círculo, e o $\sqrt{\ }$ que aparece é o \emph{preço da régua que se trouxe}. A saída é a
\textbf{dualidade}, e ela é o teorema central deste sistema.

\subsection*{Primeiro, o que é dual}\label{sub:defdual}

\begin{definicao}[bijeção dual --- a estrela]\label{def:bijdual}
Uma \textbf{bijeção dual} entre dois objectos $A$ e $B$ é a \emph{interface reversível}: uma
correspondência $\nu$ que é o seu próprio inverso, $\nu\circ\nu=\mathrm{id}$ com resíduo $0$; que
\textbf{liga sem fundir} --- preserva os dois lados em vez de os colar num terceiro; e que \textbf{não
apaga um bit}. É a \textbf{estrela}: o ponto fixo da involução que troca os dois, a plena
(Teorema~\ref{thm:estrelaseis}, grau seis, onde soma $=$ produto), o único regime com os dois sentidos.
\end{definicao}

\begin{definicao}[dual]\label{def:dual}
$A$ e $B$ são \textbf{duais} quando existe uma bijeção dual entre eles. Nenhum é prévio: a bijeção é o
seu próprio inverso, e o que $A$ é para $B$, $B$ é para $A$. \emph{Não há um melhor} --- dá no mesmo
pela cruz (ambos leem a mesma coisa) e desce pela estaca (a troca de sinal).
\end{definicao}

\subsection*{O teorema central: Gentil e Hurwitz são duais}\label{sub:central}

\begin{teorema}[Gentil e Hurwitz são duais, pela medida]\label{thm:central}
Sejam o lado \textbf{discreto} --- \emph{Hurwitz}: a norma euclidiana $N(x)=\sum x_i^2$ (a esfera),
lida pela cruz, multiplicativa para o produto \emph{bilinear} exactamente nos graus $1,2,4,8$ --- e o
lado \textbf{contínuo} --- \emph{Gentil}: a norma $\|x\|$ e a fusão homogénea (a hipérbole), sem limite
de grau. Então os dois são \textbf{duais}, e a bijeção dual é a \textbf{estrela}, realizada pela
\emph{teoria da medida}: o integral de Lebesgue --- a soma reversível
\[
  \int_a^b f \;+\; \int_{f(a)}^{f(b)} f^{-1} \;=\; b\,f(b)-a\,f(a)
\]
--- é a correspondência que leva \textbf{contar} (o discreto, $\sum$, Hurwitz) a \textbf{integrar} (o
contínuo, $\int$, Gentil) e a traz de volta. A bijeção é $f:t\mapsto t^2$ (contar $\leftrightarrow$
raiz), e \emph{nunca se avalia uma raiz}: os dois cortes --- o do domínio (Riemann, Hurwitz) e o da
imagem (Lebesgue, Gentil) --- dão a mesma contagem.
\end{teorema}

\begin{proof}
A bijeção dual é a involução $f\leftrightarrow f^{-1}$, que é a Lei~1. A conservação acima
--- a mesma já medida nos metais --- reparte o rectângulo pelos dois cortes \emph{sem resto}, logo a
correspondência é reversível ($\nu\circ\nu=\mathrm{id}$, resíduo $0$) e é uma bijeção; é a estrela
(Definição~\ref{def:bijdual}). \emph{Contar} (a soma no domínio, o discreto, Hurwitz) e \emph{integrar}
(o integral na imagem, o contínuo, Gentil) são os dois lados desta bijeção, logo $A$ e $B$ são duais
(Definição~\ref{def:dual}). E o \textbf{limite no grau oito é do lado discreto/bilinear} --- Hurwitz
classifica o bilinear ---, \emph{não do objecto}: o lado contínuo, alcançado pela soma reversível e não
por uma raiz, não o tem.
\end{proof}

\begin{definicao}[o octonião dual --- o grau oito reversível]\label{def:octoniao-dual}
O grau oito é $\mathbb{O}=\mathbb{H}\times\mathbb{H}^{*}$: \emph{dois tecidos grau quatro} (quaterniões)
ligados pela \textbf{estrela} (a bijeção dual, Definição~\ref{def:bijdual}). Ele lê-se de \emph{dois}
modos --- e são o par do Teorema~\ref{thm:central}:
\begin{itemize}
\item pela \textbf{norma bilinear} (Hurwitz, o lado discreto) é o \emph{octonião clássico}, e aí
\textbf{perde a associatividade} --- é o preço do grau oito \emph{do lado da norma};
\item pela \textbf{dualidade} (a cruz, a estrela, a bijeção reversível --- Gentil, o lado contínuo) é o
\textbf{octonião dual}: $\nu\circ\nu=\mathrm{id}$ com resíduo $0$, a estrela \emph{liga sem fundir}, e
\textbf{não perde nada} --- a volta fecha.
\end{itemize}
É o \emph{mesmo espaço}; muda a régua. E a régua dual é reversível \emph{onde} a bilinear cai --- porque
a estrela não é bilinear, e Hurwitz só classifica as bilineares: o limite do grau oito é do lado da
norma, \textbf{não do objecto}. É o octonião que este sistema usa em toda a parte --- a interface que
liga dois corpos sem os fundir, e não dissipa.
\end{definicao}

\medido{o octonião dual realiza-se, medido, no circuito do tradutor: dois tetrais (o corpo, grau $4$)
ligados pela interface (a hexal), em dimensão $8$ \textbf{reversível} --- a volta
$\text{tex}\to\text{pdf}\to\text{tex}$ fecha com resíduo $0$ por ser inteira, e um passo com perda (a
norma que arredonda) quebra-a ($1968$ de $3000$). A \emph{reversibilidade} é o que o separa do octonião
clássico, que no grau oito perde a associatividade (\code{tests/circuito\_tradutor.c}).}

\begin{teorema}[a graduação é um segundo construtor sem topo --- e está
do lado de Gentil]\label{thm:graduacao-gentil}
O passo dos tecidos $T_{k+1}=T_k+T_k^{*}$ é a estrela usada como
construtor, e a torre que ele gera não tem topo por dentro. Há um
\textbf{segundo} construtor com a mesma propriedade, e do mesmo lado:
\[
\boxed{\;
\Lambda^{0}\xrightarrow{\;d\;}\Lambda^{1}\xrightarrow{\;d\;}\cdots
\xrightarrow{\;d\;}\Lambda^{n}\xrightarrow{\;d\;}0,
\qquad d^{2}=0.
\;}
\]
\textbf{(i) Não pede fórmula nova a cada dimensão.} A regra que define
$\Lambda^{k}\to\Lambda^{k+1}$ é a mesma em todo $k$ e em todo $n$: os
índices entram antissimetrizados. O que \emph{acaba} é a
\textbf{realização} --- em dimensão $n$ não há onde continuar depois de
$\Lambda^{n}$ ---, não a linguagem. Medido: a torre tem $n{+}1$ andares
não nulos e soma $2^{n}$, verificado até $n=14$.

\textbf{(ii) E está do lado contínuo.} O que a define é a
\emph{antissimetria} --- uma indução no grau --- e não uma norma
bilinear. Por isso ela não conhece o tecto do grau oito: esse é de
Hurwitz, e classifica o \emph{bilinear}. O mecanismo de $d^{2}=0$ é a
contracção de um tensor antissimétrico com um simétrico
($\partial_i\partial_j$ é simétrico), e esse argumento \emph{não
menciona $n$} --- medido de $n=2$ a $n=40$, com o controlo a separar
(trocada a antissimetria por simetria, a contracção deixa de anular).

\textbf{(iii) E o directo e o cruzado realizam-se nela.} O par que
este paper introduziu --- o \emph{directo} $\langle a,b\rangle=\cos\theta$,
a parte simétrica, e o \emph{cruzado} $\|a\wedge b\|=\sin\theta$, a
antissimétrica --- é, na graduação, o \textbf{mesmo} produto exterior
em graus diferentes:
\[
\text{cruzado}=\alpha\wedge\beta\ (\text{grau }1{+}1),
\qquad
\text{directo}=\alpha\wedge\star\beta\ (\text{grau }1{+}2),
\]
e Lagrange fecha entre os dois,
$\text{directo}^{2}+\text{cruzado}^{2}=N(\alpha)N(\beta)$ --- medido em
$117\,649$ pares, sem falha, e \emph{sem trigonometria e sem raiz}: o
$\cos\theta$ e o $\sin\theta$ do enunciado original são a leitura
normalizada de duas quantidades inteiras.

\textbf{Estatuto.} A lei é do \emph{Corpo Algébrico} (a nilpotência
como dual da idempotência, e a graduação a tirar o limite dimensional
da linguagem). Esta é a face deste paper: um segundo construtor sem
topo, ao lado do $T_{k+1}=T_k+T_k^{*}$, e a realização do par
directo/cruzado dentro dele. O par $(\partial,d)$ \emph{não} é uma
bijeção dual --- a Definição~\ref{def:bijdual} exige
$\nu\circ\nu=\mathrm{id}$, e $d\circ d=0$ é outra coisa: é uma
adjunção.
\emph{Medido:} \code{tests/graduacao\_realiza.c} ($5{:}0$) e
\code{tests/borda\_sem\_teto.c} ($8{:}0$).
$\square$
\end{teorema}

\begin{teorema}[a face graduada do \textsc{Teorema do Gato} --- o
Jacobiano, e a lei é LOCAL]\label{thm:medida-jacobiano}
A realização graduada do \textsc{Teorema do Gato} (\emph{Corpo
Universal}, Teor.~da conservação métrica por dualidade) é a
\textbf{mudança de variável}: o determinante deixa de ser função de uma matriz e passa a
ser o \textbf{factor local da medida},
\[
\boxed{\;\mu_{n}(T(E))=\int_{E}|\det DT(x)|\,d\mu_{n}(x)\;}
\]
e, no caso conservativo,
\[
|\det DT(x)|=1\;\;\text{em todo }x
\quad\Longrightarrow\quad
\mu_{n}(T(E))=\mu_{n}(E).
\]

\textbf{(0) O que a graduação acrescenta ao lado discreto.} No
Teor.~da face discreta (\emph{Corpo Topológico}) o factor é \emph{um}
número, porque a transformação é linear. Aqui ele é uma \textbf{função
do ponto}, e a conservação passa a ser uma condição \emph{pontual} ---
que é exactamente o que a graduação faz em toda esta tabela: transforma
uma constante num campo. O determinante é o coeficiente da $n$-forma de
volume, e a $n$-forma é o topo do complexo $\Lambda^{\bullet}$: por isso
o factor de medida vive no mesmo sítio onde $d$ deixa de ter para onde
subir.

\textbf{(1) A testemunha, e ela não é linear.} O cisalhamento
$T(x,y)=(x+y^{2}+3y+5,\;y)$ tem
\[
DT=\begin{pmatrix} 1 & 2y+3\\ 0 & 1\end{pmatrix},
\qquad \det DT=1\quad\text{em todo ponto},
\]
e a entrada $\partial/\partial y$ \emph{muda} com o ponto. Logo a
conservação da medida não é uma propriedade de matrizes: é do factor de
volume ponto a ponto, e o caso linear é apenas aquele em que ele não
varia.

\textbf{(2) E onde este teorema toca no central, e onde NÃO toca.} Esta
tabela já usava Lebesgue no Teorema~\ref{thm:central}: «Hurwitz conta o
domínio, Lebesgue mede a imagem, e Gentil é a soma reversível que os
casa». Isso é uma afirmação sobre a $\sigma$-aditividade --- a medida
existe e transporta a ordem do discreto para o contínuo. O
\textsc{Teorema do Gato} diz outra coisa sobre a \emph{mesma} medida:
como ela se \textbf{comporta sob uma transformação}. São dois enunciados
distintos e é preciso não os fundir --- o central precisa que a medida
exista; o Gato mede o factor com que ela muda. O que os liga é o objecto,
não a proposição.

\textbf{(3) E a ausência de tecto é da forma, não da varredura.} A base
é $\det(I)=1$; o passo é
$\det(T\oplus T^{\ast})=\det T\cdot\det T^{\ast}=1$, e não menciona a
dimensão. A conservação não é reintroduzida em cada andar --- é
transportada pela dualidade, que é a relação $\sigma_{k}\sigma_{k'}=-1$
desta torre.

\emph{Medido:} \code{tests/conservacao\_metrica.c} ($\S$M4--$\S$M6,
$14{:}0$): $|\det DT|=1$ em $100$ pontos racionais com a derivada a
variar em $98$ deles; o gume --- a dilatação por $k$ achada em $7$ de
$7$ e o buscador a voltar vazio nos $20$ metais; e a regra da segunda
realização --- a primeira satura em $k=19$ e a segunda, em resíduos,
responde até $k=400$ em três primos. Uma saturação não é um
contra-exemplo, e conta-se em lugar separado. $\square$
\end{teorema}

\begin{teorema}[a face graduada da ponte --- onde a obstrução deixa de ser
zero]\label{thm:ponte-graduada}
A graduação do teorema anterior tem uma propriedade que só se vê quando
a obstrução \emph{não} é nula, e num domínio em estrela ela é sempre
nula. É a face que falta:
\[
\Lambda^{0}\xrightarrow{\;d\;}\Lambda^{1}\xrightarrow{\;d\;}\cdots,
\qquad
\operatorname{im}d\subseteq\ker d,
\qquad
H^{k}=\ker d/\operatorname{im}d .
\]
\textbf{(i) A imagem cabe no núcleo; a pergunta é se o PREENCHE.} Isso é
o que $d^{2}=0$ garante e não decide. Num domínio em estrela toda
fechada é exacta (Poincaré) e $H^{1}=0$ --- e é esse o regime dos
polinómios. O buraco precisa de um domínio com buraco.

\textbf{(ii) E a testemunha clássica não é polinomial.} A forma do
ângulo $(-y\,dx+x\,dy)/(x^{2}+y^{2})$ é fechada e não exacta em
$\mathbb{R}^{2}\setminus\{0\}$, e não vive no anel $\mathbb{Q}[x,y,z]$.
Registou-se em tempo próprio que \emph{a razão de não se ver o buraco
era a representação, não a matemática}.

\textbf{(iii) A realização que o mostra.} Sobre um grafo, o mesmo
complexo dá $\dim H^{1}=E-V+C$, e no ciclo $C_{n}$ a cocadeia
$\omega\equiv1$ é o análogo discreto exacto da forma do ângulo: soma
$n\neq0$ ao longo do ciclo, e toda exacta soma zero. \emph{O contínuo
descreve o buraco; o discreto exibe-o} --- e as duas descrições contam o
mesmo número, que é também o posto de $\pi_{1}$.

\textbf{(iv) A consequência metodológica.} \emph{O buraco não é
necessariamente invisível: ele pode ser invisível numa representação e
aparecer imediatamente noutra} --- e daqui sai a exigência de não basta
provar que existe obstrução, é preciso descobrir \emph{em que
representação} ela aparece. E o par completo: $\pi_1$ dá a versão
\textbf{não linear} (no oito é livre de posto $2$ e não abeliano) e
$H^1$ a \textbf{linearizada} ($\mathbb{Z}^2$, onde o comutador morre).
A passagem de uma à outra é apagar a \emph{ordem}.

\textbf{Estatuto.} A lei é do \emph{Corpo Algébrico} (localmente
contráctil $\not\Rightarrow$ globalmente contráctil); a face discreta é
do \emph{Corpo Topológico} ($\ker\partial/\operatorname{im}\partial$,
com o potencial como fibra e a contagem $E-V+C$); esta é a face
graduada, $\ker d/\operatorname{im}d$ --- o lugar onde a obstrução se define
e onde se diz \emph{porque} ela some nos polinómios. Nenhuma das três
afirma cohomologia nova: $\dim H^{1}=E-V+C$ é clássico, e o que se
acrescenta é que o buraco é \emph{um objecto só}, contado de dois modos.
\emph{Medido:} \code{tests/cohomologia\_ponte.c} ($7{:}0$).
$\square$
\end{teorema}

\setcounter{lei}{6}
\begin{lei}[o octonião dual --- ligar sem fundir]\label{lei:octoniao}
$\mathbb{O}=\mathbb{H}\times\mathbb{H}^{*}$, dois tecidos ligados pela estrela: $\nu\circ\nu=\mathrm{id}$,
resíduo $0$. Pela norma perde a associatividade; pela dualidade não perde nada.
\end{lei}

\noindent \textbf{É a sétima lei, e é operacional como as outras.} As seis descem $6\to1$ \emph{dentro} de
um corpo (Teorema~\ref{thm:estrelaseis}); a sétima \emph{liga dois}. O seu \textbf{operador} é a estrela
--- a bijeção dual (Def.~\ref{def:bijdual}), uma operação em dois sentidos ($\texttt{MOVE}(\cdot,\pm1)$)
---, e o seu \textbf{medidor} certifica-a a fechar por resíduo $0$: o circuito do tradutor --- dois
tetrais ligados pela interface, a volta $\text{tex}\to\text{pdf}\to\text{tex}$ exacta por ser inteira, e um
passo com perda (a norma que arredonda) quebra-a, $1968$ de $3000$ (\code{tests/circuito\_tradutor.c}). E
\emph{a interface é a hexal}: a costura não é o oito, é o seis --- o oito é o que se liga, o seis é onde
se liga. É o octonião que este sistema usa em toda a parte, a interface que não dissipa.

\subsection*{A torre dos tecidos: cada andar é o dual do de baixo}\label{sub:tecidos}

\textbf{A estrela não liga só dois corpos dados: ela \emph{gera} a torre inteira, um andar de cada
vez.} O primeiro tecido é a \textbf{contagem natural} $\mathbb{N}$ --- o discreto puro, o lado de Hurwitz no
seu grau mínimo. Todo andar acima sai do de baixo pela \emph{mesma} operação: o dual soma-se ao
original.

\begin{definicao}[tecido, e o passo de indução]\label{def:tecido}
Um \textbf{tecido} é um andar da torre dos sistemas numéricos. O passo que leva um andar ao seguinte é
\[
  T_{k+1} \;=\; T_k \;+\; T_k^{*},
\]
o par $(A,A^*)$ tomado a menos da \textbf{estaca} --- a involução $\nu$ com $\nu\circ\nu=\mathrm{id}$ e
ponto fixo \emph{na fronteira} --- e lido pela \textbf{cruz}, a medida que se conserva no passo. É a
bijeção dual (Definição~\ref{def:bijdual}) usada como \emph{construtor}.
\end{definicao}

\begin{teorema}[a torre dos tecidos é a estrela iterada; o salto $\mathbb{Q}\to\mathbb{R}$ é o de Hurwitz para Gentil]\label{thm:tecidos}
Partindo de $\mathbb{N}$ e aplicando o passo, a torre é
\[
  \mathbb{N} \;\xrightarrow{\;+\mathbb{N}^{*}\;}\; \mathbb{Z} \;\xrightarrow{\;+\mathbb{Z}^{*}\;}\; \mathbb{Q}
     \;\xrightarrow{\;+\mathbb{Q}^{*}\;}\; \mathbb{R} \;\xrightarrow{\;+\mathbb{R}^{*}\;}\; \mathbb{C} \;\longrightarrow\;\cdots,
\]
e três coisas, que são uma:
\begin{enumerate}
\item \textbf{é uma só operação, não quatro.} Classicamente são quatro construções distintas --- as
  diferenças ($\mathbb{Z}$), as razões ($\mathbb{Q}$), a completação ($\mathbb{R}$), os pares algébricos ($\mathbb{C}$). Aqui são o
  \emph{mesmo} passo $T{+}T^{*}$: a estaca dá a involução, a cruz conserva a medida, e o ponto fixo na
  fronteira é o que fica de \emph{fora} do par (o $0$ da diferença, o corte da completação).
\item \textbf{o salto qualitativo está num degrau só, e é o Teorema~\ref{thm:central}.} De $\mathbb{N}$ a $\mathbb{Q}$
  a estaca é \emph{algébrica} (opera, mas não completa); o degrau $\mathbb{Q}\to\mathbb{R}$ é onde \textbf{contar} vira
  \textbf{integrar} --- o discreto (Hurwitz) passa a contínuo (Gentil) pela soma reversível de Lebesgue.
  A torre não muda de mecanismo; muda de \emph{régua} nesse degrau, e é o único.
\item \textbf{a interface entre andares é uma dobra.} O bit atravessa de um tecido ao seguinte pelo
  \textbf{ponto fixo} da estrela, na dobra $\Delta=n^2+4$ do corpo (Teorema~\ref{thm:atravessa}) --- a
  estaca sobe, a cruz conserva, e a estrela fecha com resíduo $0$.
\end{enumerate}
\emph{E a indução não pára por dentro:} não há andar último, porque a régua é infinita
(Teorema~\ref{thm:atravessa}). O que a faz parar é \emph{externo} --- a \textbf{finitude do corpo} que
chega: enquanto o corpo não pára de chegar, a estaca não pára de subir e a estrela não pára de reverter.
\end{teorema}

\begin{proof}
Cada passo é a bijeção dual (Definição~\ref{def:bijdual}) instanciada: a involução $\nu$ é a diferença
($\mathbb{Z}$), a razão ($\mathbb{Q}$), o corte/limite ($\mathbb{R}$), a conjugação ($\mathbb{C}$); em todas $\nu\circ\nu=\mathrm{id}$
com resíduo $0$ e ponto fixo na fronteira, e a medida (a cruz) reparte sem resto --- é a conservação do
Teorema~\ref{thm:central}, a mesma já medida. Logo o passo está bem definido e é reversível em todos os
andares. O degrau $\mathbb{Q}\to\mathbb{R}$ é o único em que a estaca deixa de ser algébrica: $\mathbb{Q}$ opera mas não é
completo, e a completação é a soma reversível $\int f+\int f^{-1}=b\,f(b)-a\,f(a)$ (Teorema~\ref{thm:central}),
isto é, o cruzamento de Hurwitz para Gentil. Que a torre não tenha topo é o crescimento sem tecto da
dobra $n^2+4$ (Teorema~\ref{thm:atravessa}); a paragem, quando há, é a finitude do corpo, não da
régua. A construção completa, com $\mathbb{N},\mathbb{Z},\mathbb{Q},\mathbb{R}$ levantados um do outro pela estaca e conservados pela
cruz, é a do \emph{Catálogo}.
\end{proof}

\noindent \textbf{E é por isto que $\mathbb{N}$ é o primeiro tecido, não um pré-requisito.} A contagem natural
não está \emph{antes} da teoria: é o andar zero dela --- o lado de Hurwitz no grau mínimo ---, e tudo o
mais é a estrela aplicada a si mesma. \emph{Um tecido, o dual, e a subida: a torre inteira sai daí.}

\subsection*{O que gera os dois: um único bit}\label{sub:bit}

Falta dizer \emph{de que} os dois lados são feitos. E a resposta é operacional, não física: \textbf{um
único bit}. A \emph{Estrela} deste sistema é a \textbf{Estrela algébrica} --- a interface de involução
com resíduo zero ---, e não uma estrela do céu; daqui por diante escreve-se só \emph{Estrela}, e nada
nela é cosmologia.

\begin{definicao}[bit operacional e suporte neutro]\label{def:bit}
A estrutura dos slots é regida pela presença ou ausência de uma única unidade binária $b\in\{0,1\}$ sob
operadores reversíveis de \emph{convolução} e \emph{deconvolução}:
\begin{itemize}
\item $b=1$: \textbf{presença} --- o slot activo no reticulado discreto (a estrutura visível);
\item $b=0$: \textbf{ausência} --- o \textbf{suporte neutro}, o elemento neutro (o ``vazio'', apenas
como elemento neutro, sem física).
\end{itemize}
\end{definicao}

\begin{teorema}[um único bit desenha o discreto e o contínuo, e o dual é a sua ausência]\label{thm:bit}
Basta um único bit $b\in\{0,1\}$ sob o par reversível \textbf{convolução} (o gato, $\times\sigma$,
\emph{sobe}) e \textbf{deconvolução} (o esquilo, $\times\sigma'=\div\sigma$, \emph{desce}) para gerar os
dois lados do Teorema~\ref{thm:central}: o \textbf{discreto} é o próprio bit (a marca no reticulado); o
\textbf{contínuo} é o rasto da convolução (o fluxo, cuja razão $\to\sigma$); e o \textbf{dual} é a
\emph{ausência} do bit ($b=0$, o suporte neutro). E fecha: $\text{gato}\circ\text{esquilo}=\mathrm{id}$,
resíduo $0$ --- a Estrela, $(A\!\times\!b)\circ(A^{-1}\!\times\!b')=\mathrm{id}$.
\end{teorema}

\begin{proof}
Semeie-se um único bit --- $b=1$ num slot, $b=0$ (neutro) nos restantes: o \emph{delta} $[1,0]$. A
convolução ($\times\sigma$, a matriz companheira $A_m$) aplicada ao delta \textbf{desenha} o reticulado
métrico --- em $m=1$ dá $[1,0],[1,1],[2,1],[3,2],[5,3],\dots$, os Fibonacci ---, e a razão dos seus
termos converge a $\sigma$: o contínuo é o rasto do discreto. A deconvolução ($\times\sigma'$, $A^{-1}$)
devolve o bit único, resíduo $0$: a interface reverte. O dual não se constrói: é a ausência, $b=0$, dada
de graça. Logo \emph{um} bit gera o discreto (ele próprio) e o contínuo (a sua convolução), e o dual é o
bit que falta. \emph{É por isto que discretização e continuidade coexistem} --- não são dois objectos,
são um bit e o seu rasto.
\end{proof}

\medido{o bit único $[1,0]$ sob o gato desenha o metal --- Fibonacci em $m=1$, a razão $\to\varphi$ ao
oitavo passo ---, e o esquilo devolve-o com resíduo $0$; a volta $\text{gato}\circ\text{esquilo}=\mathrm{id}$
fecha em $\mathbb{R}^2$ a $\mathbb{R}^8$, e o dual é a ausência $b=0$ (\code{tests/neuronio.c}).}

\subsection*{E a bijeção dual é um oscilador --- que é o relógio}\label{sub:oscilador}

A estrela não é uma correspondência parada: \textbf{oscila}. E o oscilador não é um objecto novo ---
\textbf{é o relógio} (\S\ref{sec:renorm}), dito como dinâmica. As duas escadas do sistema correm ao
mesmo tempo, em lados duais --- as \emph{dimensões sobem} (a torre, a indução, o lado contínuo de
Gentil) e as \emph{interpretações descem} (a projeção, a meta-indução, o lado discreto de Hurwitz;
\S\ref{sec:corpo}) ---, e a estaca troca-os. E realiza-se em \textbf{dois passos, um por lei}:

\begin{center}\small
\begin{tabular}{@{}llll@{}}
\toprule
o passo & a lei & o período & o que faz \\
\midrule
reflectir & \textbf{Lei 1} ($1\dg=-1$) & \textbf{2} & o espelho: troca os dois lados \\
rodar & \textbf{Lei 2} ($T\dg=-T$, $T^2=-1$) & \textbf{$q$} & o relógio: avança uma marca, \emph{é} o oscilador \\
\bottomrule
\end{tabular}
\end{center}

\begin{teorema}[o oscilador é o relógio]\label{thm:osciladorrelogio}
Num relógio de $q$ marcas, avançar uma marca por instante é a Lei~2 --- o rotor que roda o par
$(\text{posição},\text{velocidade})$ $90^\circ$ em cada plano. As duas coordenadas ficam a um quarto
de volta uma da outra: \textbf{uma sobe enquanto a outra desce}, e \textbf{trocam}; a órbita fecha ao
fim de $q$ passos (o \emph{período}); e a norma $a^2+b^2$ --- a contagem de Hurwitz, o raio do relógio
--- \emph{não se move}. Logo o oscilador \emph{é} o relógio, e o período é o do relógio: $q=4$ no
mínimo (o rotor $J$, os quatro quadrantes), e o degrau seguinte é o metal.
\end{teorema}

\begin{proof}
O relógio vive no círculo (\S\ref{sec:renorm}): a velocidade máxima é meia volta, e o gerador que a
percorre é a Lei~2, $T^2=-1$, com período quatro na régua mínima. Posição e velocidade são a cruz do
relógio --- uma mede, a outra ordena ---, e a Lei~2 roda uma na outra: $\frac{d}{dt}\cos=-\sin$,
$\frac{d}{dt}\sin=\cos$, que é $J$. A norma $a^2+b^2$ é o invariante da rotação (o raio), imóvel por
construção; é a contagem de Hurwitz, e é a metade discreta. A metade contínua --- o próprio seno, a
espiral --- é o \emph{integral} desta contagem (Teorema~\ref{thm:central}, a soma reversível), e não
se avalia: conta-se, e a inversa veste a roupa no fim. \emph{Um relógio é um oscilador visto pela
contagem; um oscilador é um relógio visto pelo integral} --- e a estrela é a bijeção entre os dois.
\end{proof}

\subsection*{Mas a bijeção não é isócrona --- e é aí que mora o oito de Hurwitz}\label{sub:isocrona}

\textbf{Os dois relógios não correm o mesmo tempo.} A bijeção dual troca a ordem por ordem --- «duas
réguas com assinaturas diferentes marcam a \emph{mesma ordem} e correm \emph{tempos diferentes}»
(\emph{Teoria}) --- mas \textbf{não é isócrona}: não preserva o período, \emph{reparametriza} o tempo.
Isócrono, no sentido do pêndulo, é o relógio cujo período não vê a amplitude; e o isócrono verdadeiro
deste sistema é a \textbf{plena}, período seis. O relógio discreto, não.

\begin{teorema}[a bijeção dual não é isócrona; o oito de Hurwitz é o período do lado discreto]\label{thm:isocrona}
Combinar o \textbf{par} --- dois operandos, o dual --- custa, no lado \textbf{discreto} (Hurwitz, o
produto bilinear), a tábua inteira $n\times n=n^2$ produtos. Logo o seu \emph{tempo} cresce com a
dimensão --- não é constante, não é isócrono --- e no grau oito vale $8^2=\mathbf{64}$: \emph{dois
lances de um lado são sessenta e quatro do outro}. E é exactamente aí que esse relógio \textbf{pára}:
a norma bilinear fecha em oito e quebra em dezasseis. \textbf{O oito de Hurwitz é o período do relógio
discreto}, o tecto do tempo $n^2$ --- não uma lei do mundo. O lado \textbf{contínuo} (Gentil) não abre
a tábua: mede pela \emph{norma}, um escalar, outro tempo --- e por isso não tem tecto.
\end{teorema}

\begin{proof}
O produto bilinear de $\sum a_i e_i$ por $\sum b_j e_j$ é $\sum_{i,j} a_ib_j\,(e_ie_j)$: os $n^2$
cross-produtos da tábua, todos, e todos têm de ficar consistentes (a associatividade, a norma). A
consistência sobrevive a três quedas --- a ordem (grau dois), a comutatividade (grau quatro), a
associatividade (grau oito) --- e na quarta, grau dezasseis, cai a norma. Três quedas com norma, logo
o topo é $2^3=8$, onde a tábua tem $8^2=64$ entradas. O produto \emph{homogéneo} de Gentil não soma a
tábua: usa $\|z\|$, um número por operando, e a norma multiplicativa não pede a consistência da tábua
--- por isso não pára (§\ref{sec:completo}, medido de $\mathbb{R}^2$ a $\mathbb{R}^7$ em \code{nne.c}).
Os dois relógios marcam a mesma ordem e correm tempos diferentes: $n^2$ com tecto de um lado, a norma
sem tecto do outro. \emph{A bijeção liga-os, mas não os sincroniza.}
\end{proof}

\noindent \textbf{E isto dá a interpretação que faltava ao oito.} Dizer «Hurwitz classifica o
bilinear, logo não é lei geral» é verdade e não explica nada. O \emph{mecanismo} é temporal: o lado
discreto é um relógio cujo tempo é a tábua $n^2$, e uma tábua só fecha (a norma sobrevive) enquanto
há uma propriedade a perder antes dela --- e há exactamente três. O oito não é onde o mundo acaba: é
onde \emph{este relógio} completa o seu período. \emph{O contínuo corre noutro tempo, e continua.}

\medido{combinar o par custa $n^2$ produtos --- $1,4,16,64,256$ nos graus $1,2,4,8,16$ ---, e no grau
oito são $\mathbf{64}$: o tempo do relógio discreto não é constante (não é isócrono) e tem tecto no
oito, onde a norma bilinear fecha e a seguir quebra; o contínuo, pela norma, não (\code{tests/estelar\_completo.c}).}

\subsection*{O bit é o ponto fixo da estrela, e atravessa: as interfaces são as dobras}\label{sub:atravessa}

\textbf{A não-isocronia tem um ritmo, e o ritmo diz onde estão as interfaces --- mas o passo não é o
$\mathrm{lcm}$ de períodos inteiros: é a \emph{dobra} de cada corpo.} Um corpo da família metálica
$x^2=m\,x+1$ tem por período o seu \textbf{discriminante} $\Delta=m^2+4$ --- a \emph{dobra}: ouro
($m=1$) em $5$, prata ($m=2$) em $8$, bronze em $13$; e o membro dual, o ouro branco $x^2=4x-1$, em
$4^2-4=12$. São as dobras, calculadas e não tabeladas (\emph{Catálogo}), que marcam as interfaces. E o
que as atravessa é \emph{um} objecto que a teoria usa desde o início sem nunca o derivar: o
\textbf{ponto fixo} da estrela.

\begin{definicao}[o bit é o ponto fixo da estrela --- derivado]\label{def:bitfixo}
A estrela é a involução $\nu(x)=-1/x$ (a Lei~1 no elemento: $x^\dagger=-1/x$, $x\,x^\dagger=-1$). O seu
\textbf{ponto fixo} não se postula, deriva-se:
\[
  \nu(x)=x \;\iff\; x=-\tfrac1x \;\iff\; x^2=-1 \;\iff\; x=i .
\]
O \textbf{bit} é esse ponto fixo: $i$, na \textbf{fronteira} $|x|=1$ --- e \emph{não há ponto fixo
real}, pois $x^2=-1$ não tem raiz em $\mathbb{R}$. É por isto que o bit tem de \textbf{atravessar} para o
imaginário: a fronteira que a teoria chamava ``o ponto fixo'' é o círculo unitário, e o ponto é $i$.
\end{definicao}

\begin{teorema}[o bit atravessa por continuação analítica; as interfaces são as dobras metálicas]\label{thm:atravessa}
A estrela $-f=f^{-1}$ na forma de potência $f=b\,x^a$ tem \textbf{solução única} $a^2=1$, isto é
$f(x)=i\,x$ (a rotação, período quatro). A sua \emph{análise} --- $f^{(n)}=f^{-1}$ --- \textbf{continua}
o expoente:
\[
  f^{(n)}=f^{-1},\ f=b\,x^a \;\Longrightarrow\; a-n=\tfrac1a \;\Longrightarrow\; a^2=na+1
  \;\Longrightarrow\; a=\sigma_n=\tfrac{n+\sqrt{n^2+4}}2 .
\]
Três coisas, que são uma:
\begin{enumerate}
\item \textbf{o bit atravessa, e a medida fica intacta.} Na fronteira $|x|=1$ vale $|x^a|=|x|^a=1$ para
  \emph{todo} expoente $a$: a continuação analítica leva $a$ do inteiro $n$ (contar, o discreto de
  Hurwitz) à média metálica $\sigma_n$ (medir, o contínuo de Gentil) com a norma \textbf{presa em $1$}.
  Fora do círculo não: $\sigma>1$ cresce e o conjugado $\sigma^\dagger=-1/\sigma<1$ encolhe, e é o
  \emph{produto} $\sigma\sigma^\dagger=-1$ que repõe a fronteira.
\item \textbf{as interfaces são as dobras.} O período de cada corpo é a dobra $\Delta=n^2+4$ (o
  discriminante), não o $\mathrm{lcm}$: $5,8,13,20,\dots$ na família, $12$ no dual. É a dobra que
  cresce quando a curva se aplana --- ``o que a curva perde, o discriminante ganha'' (\emph{Teoria}).
\item \textbf{a subida é a torre metálica, e as voltas limpas são de Lucas.} $\sigma^k=F_{k-1}+F_k\sigma$
  com norma $N(\sigma^k)=(-1)^k$ --- a alternância preto-e-branco, a estaca; e $\sigma_n$ é potência
  \emph{pura} do ouro exactamente nos índices de Lucas ímpar, $\sigma_{L_k}=\varphi^{k}$
  ($n=1,4,11,29,\dots$).
\end{enumerate}
E a \textbf{régua é infinita} --- $n$ não tem tecto, a dobra $n^2+4$ cresce sem parar ---; o
\emph{finito} é o corpo que chega, e a sua finitude é a \textbf{condição de paragem}.
\end{teorema}

\begin{proof}
$-f=f^{-1}$ com $f=b\,x^a$ dá $b^{-1/a}x^{1/a}=-b\,x^a$: o expoente força $1/a=a$, logo $a^2=1$, e o
coeficiente $1/b=-b$, logo $b^2=-1$, $b=i$ --- solução única (medido em \code{tests/lei2.c}: dos
expoentes racionais, só os de $a^2=1$ servem). Para $f^{(n)}$, a $n$-ésima derivada de $b\,x^a$ tem
expoente $a-n$; igualá-la a $f^{-1}$ (expoente $1/a$) dá $a-n=1/a$, isto é $a^2=na+1$, cuja raiz $>1$ é
$\sigma_n$. O ponto fixo de $\nu(x)=-1/x$ é $x^2=-1$, isto é $i$ (Definição~\ref{def:bitfixo}), e
$|i|=1$; em $|x|=1$, $|x^a|=|x|^a=1$ para todo $a$, logo a continuação preserva a norma. A alternância
$N(\sigma^k)=(-1)^k$ e $\sigma_{L_k}=\varphi^k$ são as da \emph{Teoria} (a companheira e a família de
potências). Que não haja tecto é $n^2+4\to\infty$; a paragem, quando há, é a finitude do corpo.
\end{proof}

\noindent \textbf{E é por isto que o ponto fixo faltava.} A teoria dizia ``involução com ponto fixo na
fronteira'' em cada degrau e nunca dizia \emph{qual}: é $i$, e a fronteira é o círculo $|x|=1$. O bit
não passa por sincronia de relógios inteiros --- \emph{atravessa} por continuação analítica, e o único
sítio onde a medida sobrevive à travessia é o círculo onde ele mora.

\medido{o ponto fixo da estrela deriva-se --- $\nu(x)=x$ força $x^2=-1$, que em $\mathbb{Z}[i]$ é $i$
(exacto) e \emph{não} tem raiz real; está em $|x|=1$ ($|i|^2=1$); a norma é imune ao expoente na
fronteira --- $|i^k|=1$ e $|N(\sigma^k)|=1$ na torre metálica, dos dois lados ---; e a dobra é
$n^2+4$ ($5,8,13$) (\code{tests/estrela\_pontofixo.c}).}

\subsection*{E aplicar isto é a transformada universal: o Dirac toma o valor na transição}\label{sub:dourada}

\textbf{O ponto fixo não é só onde o bit atravessa --- é onde a transformada o lê.} A transformada
deste sistema \emph{é} a avaliação nas raízes da borda (Gelfand: o espectro de $\mathbb{Z}_p[x]/(\beta)$ são as
raízes de $\beta$), e essas raízes são os pontos fixos $\sigma$. O mecanismo é o \textbf{Dirac}: a soma
sobre a órbita da torção colapsa num ponto ---
\[
  \sum_{k}\chi_k(j)\,\chi_{-k}(j') \;=\; n\,\delta_{j,j'}
\]
--- e o ponto onde colapsa é a \textbf{transição do relógio}, o ponto fixo. Pela peneira do Dirac, a
transformada \emph{toma o valor da função na transição}: $F[f]=f(\sigma)$. A \textbf{dourada} ($m=1$,
$\sigma=\varphi$) é o caso mais simples --- ``o ouro não salta nó nenhum, vazamento zero'' ---, e
avalia em $\mathbb{Z}[\varphi]$ \emph{sem avaliar uma raiz}. É a mesma transformada universal já medida
(\code{tests/transformada.c}, §U1 o Dirac, §U5 a amostra): o ponto fixo e a transformada são o mesmo
objecto lido de dois lados --- \emph{atravessar} e \emph{ler}.

\subsection*{Fundir \emph{vários} corpos num só: o viveiro}\label{sub:viveiro}

A estrela liga \emph{dois} corpos. Mas compor é fundir \emph{muitos} num só --- um documento tem dezenas
de fontes, cada uma um corpo, e o resultado é um --- e essa é uma operação sua, complementar da estrela
e igualmente derivada. Os corpos não são uma lista: são um \textbf{viveiro}, e a lei que os funde não é
a soma nem o produto.

\begin{definicao}[o cruzamento --- a junção do viveiro]\label{def:viveiro}
Dados dois corpos $\mathbb{R}^a$ e $\mathbb{R}^b$, o \textbf{cruzamento} é
\[
  \mathbb{R}^a \vee \mathbb{R}^b \;=\; \mathbb{R}^{\,\mathrm{lcm}(a,b)},
\]
o \emph{menor} corpo que contém os dois. Ele \textbf{voa}: é ele próprio um corpo (tem inverso para todo
não-nulo), e não apenas um espaço com a dimensão certa.
\end{definicao}

\noindent \textbf{E não são as duas que pareciam servir.} A \emph{soma direta} $\mathbb{R}^a\oplus\mathbb{R}^b$ (dimensão
$a{+}b$) nunca voa --- $(1,0)$ e $(0,1)$ são não-nulos e o seu produto é zero, um divisor de zero: são
dois pássaros na gaiola, não um novo. O \emph{tensorial} $\mathbb{R}^a\otimes\mathbb{R}^b$ (dimensão $a{\cdot}b$) voa
\emph{só} quando $\gcd(a,b)=1$ --- e aí coincide com o cruzamento, porque $a{\cdot}b=\mathrm{lcm}$. A
\textbf{junção} $\vee$ voa \emph{sempre}, e não há corpos incruzáveis.

\begin{teorema}[fundir muitos é o \emph{fold} da junção, e está bem definido]\label{thm:viveiro}
O cruzamento é comutativo, associativo, \textbf{idempotente} ($\mathbb{R}^a\vee\mathbb{R}^a=\mathbb{R}^a$) e tem neutro
($\mathbb{R}^1$): é um \textbf{semirretículo}. Logo fundir $N$ corpos é o \emph{fold} da junção,
\[
  \mathbb{R}^{a_1}\vee\mathbb{R}^{a_2}\vee\cdots\vee\mathbb{R}^{a_N} \;=\; \mathbb{R}^{\,\mathrm{lcm}(a_1,\dots,a_N)},
\]
e o resultado \textbf{não depende da ordem} nem das repetições --- é \emph{um} corpo, o menor que contém
todos, e voa. \emph{É assim que muitos corpos viram um.}
\end{teorema}

\subsection*{O DTC na torre toda: sincronizar os relógios é o cruzamento do viveiro}\label{sub:dtctorre}

\textbf{Cada tecido é um relógio, e a torre são muitos --- o DTC generaliza-se sincronizando-os.} O
inversor (\S\ref{sub:inversor}) casa o factor \emph{num} corpo; na torre há um relógio por andar, cada
um com o seu período --- a \textbf{dobra} $\Delta$ do corpo (\S\ref{sub:atravessa}) ---, e esses
períodos \emph{não são lineares}. Sincronizá-los é uma operação que a teoria já tem: o \textbf{cruzamento
do viveiro}. Dois relógios de períodos $a$ e $b$ alinham no \emph{menor} instante comum,
$\mathrm{lcm}(a,b)$, que é exactamente $\mathbb{R}^a\vee\mathbb{R}^b=\mathbb{R}^{\mathrm{lcm}(a,b)}$ (Teorema~\ref{thm:viveiro}).

\begin{corolario}[o DTC na torre: o casamento é o fold do viveiro]\label{cor:dtctorre}
O controlo directo de torque generaliza-se da família $\star_s$ \emph{num} corpo (\code{tests/dtcn.c}) à
\textbf{torre inteira}: em cada interface (a dobra, Corolário~\ref{cor:dobra}) o inversor comuta para
casar $\mathrm{fp}=1$ pelo dual ($\nu\circ\nu=\mathrm{id}$), e a \textbf{sincronia} dos relógios de todos
os andares é o fold do cruzamento,
\[
  \bigvee_{k} \mathbb{R}^{a_k} \;=\; \mathbb{R}^{\,\mathrm{lcm}(a_1,\dots,a_N)} .
\]
O relógio da torre \emph{não é linear} --- bate no $\mathrm{lcm}$ das dobras, não num passo fixo ---, e é
por isso que quem o casa é a \textbf{modulação} (a comutação do inversor), e não um oscilador de passo
constante: \emph{um DTC por andar, sincronizados pelo viveiro}.
\end{corolario}

\noindent \textbf{E é a cascata de dobras (Corolário~\ref{cor:dobra}) lida como controlo.} Cada andar
dobra o relógio; o viveiro diz onde todos batem juntos; e o inversor, em cada dobra, empurra $s\to\pm1$
(o imposto $1-s^2$ a zero, \S\ref{sub:inversor}). A torre inteira casa-se andar a andar, sem oscilador e
sem armazenar --- o dual em cada degrau, o $\mathrm{lcm}$ a sincronizar.

\medido{a torre casa por controlo directo de torque, sincronizada pelo viveiro, e em INTEIRO: sobre
quatro andares de dobras $\Delta_n=n^2+4$ (períodos $\{5,8,13,20\}$), o imposto $\Pi(s)=(1-s^2)\|a\times
b\|^2$ anula exacto em $s=\pm1$ e o DTC escolhe do trial $\{-1,0,+1\}$ o $\pm1$ que o zera (fp$=1$ por
modulação); o viveiro sincroniza no $\mathrm{lcm}=520$, \emph{independente da ordem} (fold à esquerda $=$
à direita), contendo cada andar e o menor a fazê-lo; e a torre bate junta \textbf{pela primeira vez}
exactamente em $t=\mathrm{lcm}$, com cada andar já casado --- um DTC por andar, sincronizados pelo
viveiro. O torque é o cruzado: o imposto cavalga $\|a\times b\|^2$ e casá-lo ($s\to\pm1$) é casar fp$=1$.
Resíduo $0$ (\code{tests/dtc\_viveiro.c}).}

\begin{proof}
Comutatividade e associatividade fazem o fold independente de ordem e de agrupamento; a idempotência
colapsa as repetições. Como $\mathrm{lcm}$ tem exactamente estas três propriedades, o fold dá
$\mathbb{R}^{\mathrm{lcm}(a_1,\dots,a_N)}$ para qualquer parentização. Contém cada pai porque $\mathbb{R}^x\subset\mathbb{R}^n
\iff x\mid n$, e cada $a_i$ divide o $\mathrm{lcm}$; e é o \emph{menor} porque nenhum $n<\mathrm{lcm}$ é
divisível por todos. Voa porque o viveiro é fechado: o filhote é da mesma espécie dos pais, e por isso
se pode cruzar de novo, e de novo.
\end{proof}

\noindent \textbf{E o mecanismo não é o rótulo.} $\mathrm{lcm}$ é o \emph{quê}; o \emph{como} é o
relógio --- combinar dois relógios de períodos $a$ e $b$ fecha no passo $\mathrm{lcm}(a,b)$, que é onde
a figura de La~Hire volta a si. O viveiro é o relógio quando dois (ou muitos) se combinam; e diferente
da estrela, o cruzamento \textbf{não reverte} --- o filhote não devolve os pais, e é a idempotência que
o faz \emph{viveiro} e não fábrica. \textbf{É esta a fusão do compositor}: cada fonte é um corpo, e o
documento é a junção de todas --- o menor que as contém. \emph{A maioria dos corpos funde-se assim.}

\medido{fundir vários corpos num só é o fold de $\vee$: sobre sete listas, a junção dá
$\mathbb{R}^{\mathrm{lcm}}$ \emph{independente da ordem} (esquerda $=$ direita), o filhote único \textbf{voa} (é
corpo, medido em $\mathrm{GF}(2^n)$) e contém todos os pais; a soma directa nunca voa (divisor de zero
exibido) e o tensorial só entre coprimos, tudo por varredura exaustiva, resíduo $0$
(\code{tests/viveiro.c}, \code{tests/cruzamento.c}).}

\subsection*{E os dois lados não se acoplam: falam pela estrela}

\begin{teorema}[a ponte reverte, não funde]\label{thm:pontereversivel}
A ligação de dois corpos que \emph{reverte} não os funde num terceiro: um lado emite ($-1$), o
outro absorve ($+1$), e a estrela no meio troca-os sem os misturar --- liga, não cola.
\end{teorema}

\begin{proof}
Fundir dois corpos num só produto interno (o tensorial, ou colá-los bilinearmente) apaga a
distinção entre eles e, além do grau oito, mata a norma. A estrela não funde: é uma \emph{fronteira}
--- $a=1/a$, o ponto fixo (§\ref{sec:dois}) --- que os dois lados partilham sem partilhar corpo. O
que $A$ emite é o que $B$ absorve, passo a passo, e o par não se move ($\rho_A\rho_B$ imóvel). Ler
não altera quem escreveu; escrever não altera a leitura --- duas retas independentes, o tropical,
\emph{onde nada se desfaz}. E porque reverte, não apaga: paga zero (§\ref{sec:medir}).
\end{proof}

\noindent \textbf{E aqui a arquitetura e a álgebra são a mesma frase.} Um sistema que compõe (o
tradutor, o motor) e o que o hospeda são os dois lados; \emph{ligá-los copiando um para dentro do
outro é fundi-los} --- é apagar de um lado para refazer no outro, e apagar dissipa. A ligação certa é
a estrela: o hospedeiro emite os dados, o módulo absorve-os por $\texttt{MOVE}$, e devolve pela mesma
porta. \emph{Nenhum disco se copia para dentro de outro}; há um só, endereçado dos dois lados, e a
fronteira entre eles reverte. É por isso que o sistema não tem um \emph{buffer} de entrada: um
\emph{buffer} é a cópia que a estrela existe para não fazer --- e a estrela não a faz porque
\textbf{não oscila} (\S\ref{sub:mecanica}, \S\ref{sub:inversor}).

\subsection*{A conta que fecha}

\begin{center}\small
\begin{tabular}{@{}llll@{}}
\toprule
 & o grau & o que é & de que escada \\
\midrule
buraco branco & \textbf{$4,8,\dots$} & corpo que opera, emite & a torre, a subir sem topo \\
\textbf{a estrela} & \textbf{6} & a plena: soma $=$ produto, reversível & as interpretações, o topo \\
buraco negro & \textbf{$4,8,\dots$} & corpo que opera, absorve & a torre, a subir sem topo \\
\bottomrule
\end{tabular}
\end{center}

\noindent \textbf{As duas escadas do sistema encontram-se na estrela.} A torre sobe $1,2,4,8,\dots$
--- os lados, os corpos que operam, \emph{sem teto} --- e as interpretações descem $6,5,4,3,2,1$ ---
a plena no cimo. Essa descida $6\to1$ \emph{é} a progressão das seis leis (Lei $6$ hexal $\to$ Lei $1$
dual, a projeção que esquece uma fase por degrau --- \emph{Catálogo}, \code{obs:seis-leis}). A estrela é
o topo de uma e a ponte da outra, e é o mesmo objecto: o ponto fixo,
onde operar e ler deixam de se distinguir. \emph{Não foi ajustado para coincidir} --- coincide
porque a reversibilidade obriga o seis, a operação obriga a torre, e o único sítio onde as duas se
cruzam é a estrela. \textbf{E o processo não termina: acopla-se qualquer par de corpos, porque a
estrela reverte em vez de fundir.}

\medido{o seis é o único com soma $=$ produto dos divisores próprios em $20\,000$ inteiros;
$\mathbb{H}$ não comuta ($295/300$) e associa, $\mathbb{O}$ nem uma nem outra; o \textbf{lado
discreto} (Hurwitz, o bilinear $\sum x_i^2$) fecha em oito e falha em dezasseis --- o limite é da
\emph{contagem}, e Hurwitz \emph{classifica} o bilinear, não põe parede no objecto; a \textbf{bijeção
dual} (a estrela, pela medida) é reversível --- os dois cortes de $\int f$, o do domínio
(Riemann/Hurwitz) e o da imagem (Lebesgue/Gentil), dão a mesma contagem, e $\int f+\int f^{-1}$ fecha
(Young no reticulado), \emph{sem avaliar uma raiz}; o \textbf{oscilador é o relógio} --- a Lei~2 roda
(período $4$, com $a^2+b^2$ imóvel) e a Lei~1 reflecte (período $2$), um lado sobe enquanto o outro
desce e trocam; e a estrela \textbf{reverte}, $x\dg{}\dg=x$, resíduo $0$
(\code{tests/estelar\_completo.c}). A testemunha directa do lado contínuo, $\|xy\|=\|x\|\|y\|$ de
$\mathbb{R}^2$ a $\mathbb{R}^7$, está em \code{tests/nne.c}.}

%======================================================================
\section{O ALCANCE, e as três peças que o realizam}\label{sec:alcance}

\noindent\emph{Estas três vieram do centro} (\code{algebrico sec:camadas}), e vieram porque
são \textbf{realizações do alcance} e não fundamento: a completude, a construção do supremo
e o $\pi$ só existem onde o objecto \emph{não fecha} --- isto é, aqui, em $\tau=+1$. No
centro elas ocupavam vinte e duas páginas a fazer o que esta camada existe para fazer.

\emph{O que o centro guarda é a razão de elas serem precisas}: que $\mathbb{Q}$ não é
completo, e que a álgebra opera sem alcançar. O que ele não guarda é o alcance em si.

\section{A completude}\label{sec:completude}

\noindent\textbf{Primeiro, as duas caixas, que são duas.} O \S\ref{sec:corte}
decidiu $a/b<\sigma_m$: é o corte de um metálico \emph{específico}. O mecanismo
geral é outro, e enuncia-se agora:
\[
\boxed{
\begin{array}{ll}
\text{mecanismo \textbf{extremo explícito}} &
\text{a família metálica: } \sigma_m,\ t_k,\ \abs{\sigma\dg}^{k}\\[3pt]
\text{mecanismo \textbf{geral}} &
\text{CF arbitrária: } (a_k)\to(p_k/q_k)\to\abs{\det M_k}=1\to I_k
\end{array}}
\]

\begin{definicao}[os intervalos do mecanismo geral]\label{def:Ik}
Dada $(a_0;a_1,a_2,\ldots)$ com $a_k\geq1$, e $p_k/q_k$ os convergentes
(Def.~\ref{def:real}), põe-se para $k\geq0$
\[
I_k=\Bigl[\min\Bigl(\tfrac{p_k}{q_k},\tfrac{p_{k+1}}{q_{k+1}}\Bigr),\
\max\Bigl(\tfrac{p_k}{q_k},\tfrac{p_{k+1}}{q_{k+1}}\Bigr)\Bigr].
\]
Então, de $\abs{p_kq_{k-1}-p_{k-1}q_k}=1$:
\[
\boxed{\;I_{k+1}\subset I_k,\qquad \abs{I_k}=\frac{1}{q_kq_{k+1}}\;}
\]
--- uma cadeia \textbf{encaixada} cujo comprimento é dado em inteiros. E, por
$q_k\geq F_{k+1}$ (Teor.~\ref{thm:ouro}),
$\abs{I_k}\leq 1/(F_{k+1}F_{k+2})\leq2^{-k}$ para $k\geq1$: \emph{o comprimento
cai pelo menos como $2^{-k}$}, e é este o elo que a unicidade usa.
\end{definicao}

\medskip
\noindent\textbf{E aqui a distinção que decide tudo, porque é ela que se costuma
perder.} É verdade que \emph{intervalos encaixados de comprimento a tender para
zero não estabelecem, por si sós, a existência do ponto de intersecção}: em
$\mathbb{Q}$ há cadeias assim sem habitante. \textbf{Mas o objecto construído não é
o intervalo --- é o caminho}, e o caminho não espera por ponto nenhum. E o
\S\ref{sec:pontofixo} já disse \emph{o que} é esse ponto: o \textbf{ponto fixo da
Möbius}, que a órbita persegue e o andar de baixo não contém.
\[
\boxed{\;\begin{array}{ll}
\text{o \textbf{intervalo}} & \text{é a testemunha, e precisaria de um habitante que
o ocupasse}\\[3pt]
\text{o \textbf{caminho}} & \textbf{é} \text{ o habitante --- a sucessão } (a_k)
\text{ \emph{é} o real}
\end{array}\;}
\]
O \emph{Corpo Algébrico} (Teor.~\code{thm:central-continuo}) constrói
$\mathbb{R}\simeq\{\text{caminhos}\}/\!\sim$, e a completude é lá \textbf{interna}: a
sucessão de Cauchy dos truncados estabiliza \emph{bit a bit} dentro do próprio
objecto, sem sair dele. E o mecanismo deste paper diz \emph{porquê} ela aperta em
vez de vaguear --- o vazamento é zero porque a família é de Pisot,
$\abs{\sigma\dg}<1$ (\S\ref{sec:pisot}). \textbf{A completude não é aqui um
postulado acrescentado: é consequência do vazamento zero.} A identificação final
com $\mathbb{R}$, essa sim, é o teorema clássico de unicidade a menos de
isomorfismo, e cita-se como tal.

\medskip
\noindent\textbf{Mas o que a realização entrega diz-se com precisão, e são duas
coisas.}

\begin{teorema}[o mecanismo produz um corte, e um só]\label{thm:completude}
Seja $(a_k)_{k\geq1}$ uma sucessão \textbf{infinita} de quocientes parciais --- o
caso finito é o racional, e a sua representação termina (Def.~\ref{def:real}). A
classe define-se \textbf{directamente pelos convergentes}, sem passar pelo
\S\ref{sec:corte}, que é o caso metálico:
\[
\boxed{\;A_{(a_k)}=\bigl\{r\in\mathbb{Q}:\ r<c_{2j}\ \text{para algum convergente
de índice par}\bigr\}\;}
\]
--- os pares crescem e ficam todos abaixo, os ímpares descem e ficam todos acima
(Teor.~\ref{thm:encaixe}). A cadeia é
$(a_k)\to(p_k/q_k)\to I_k\to\text{corte}$, e o metálico fica como caso explícito. Então:

\emph{(i) $A$ é um corte de Dedekind.} $A$ e o seu complementar são não vazios, e
$A$ é fechado para baixo. E \textbf{$A$ não tem máximo}, com a testemunha dada pela
\emph{própria construção} --- não pelo \S metálico: os convergentes de índice
\textbf{par} crescem estritamente e ficam todos abaixo dos de índice ímpar, logo
todos em $A$; dado $r\in A$, \textbf{algum} convergente par \emph{posterior}
excede-o --- não necessariamente o seguinte, pois $r$ pode estar muito perto do
limite. O ponto não vive
em $\mathbb{Q}$: é por isso que \emph{o corte é o objecto}.

\emph{(ii) O habitante é único, pelo princípio do pombal, e a conta é inteira.}
Por Def.~\ref{def:Ik} tem-se $\abs{I_K}\leq2^{-K}$; logo, se $r_1/s_1$ e
$r_2/s_2$ habitam $I_K$ com $2^{K}>s_1s_2$, então
$\abs{I_K}\leq2^{-K}<1/(s_1s_2)$ e portanto
\[
\Bigl|\frac{r_1}{s_1}-\frac{r_2}{s_2}\Bigr|<\frac{1}{s_1s_2}
\quad\Longrightarrow\quad
\abs{r_1s_2-r_2s_1}<1 ,
\]
e $\abs{r_1s_2-r_2s_1}$ é um \textbf{inteiro} --- logo é zero, logo
$r_1/s_1=r_2/s_2$. \textbf{Não há dois racionais distintos em $I_K$.}

\emph{(iii) E o corte é determinado pela sucessão}: sucessões distintas dão cortes
distintos, com o racional separador exibido. $\square$
\end{teorema}
\colunas{a identificação dos cortes com o corpo ordenado completo}
        {o mecanismo que produz o corte, e a unicidade pelo pombal}
        {\code{geometria\_real.c} \S L12: $2\,624\,400$ pares sem violação;
quatro cláusulas nos oito $\sigma_m$; $28$ pares separados}

\medskip
\noindent Assim a divisão fica limpa, e cada peça faz exactamente o seu trabalho:
\[
\boxed{\;\begin{array}{l}
\text{o \textbf{mecanismo} entrega o corte;}\\
\quad\text{e o Algébrico identifica os cortes com o corpo completo}\\[3pt]
\text{o \textbf{pombal} dá a unicidade dos \textsc{aproximantes racionais}:}\\
\quad\text{não há dois racionais distintos em }I_K
\end{array}\;}
\]
\[
\boxed{
\begin{array}{l}
\text{pombal}\ \Longrightarrow\ \text{unicidade dos \textbf{aproximantes racionais}}\\[2pt]
\text{completude herdada}\ \Longrightarrow\ \text{existência do \textbf{habitante real}}
\end{array}}
\]
O pombal não estabelece, por si, a unicidade do \emph{real}: essa vem da estrutura
de corpo ordenado e da densidade de $\mathbb{Q}$, ambas herdadas. Nenhuma das
peças se obtém por varredura.

\medskip
\noindent\textbf{Correspondência estrutural herdada} (não é ingrediente da prova
de completude, que não depende dela). Como consequência adicional da estrutura
espectral já estabelecida no \emph{Corpo Algébrico} --- a convolução emerge da
transformada universal, $\delta$ é a sua identidade, e a deconvolução é a divisão
espectral, exacta onde o espectro não tem zeros ---, os elementos desta base
\emph{são} os Diracs transladados,
\[
\boxed{\;e_t=\delta_t,
\qquad \langle b,e_t\rangle=b_t\ \text{é a peneira},
\qquad (b\ast\delta_{t})\ast\delta_{-t}=b\;}
\]
--- e portanto, \emph{nessa leitura}, ler uma coordenada e localizar um ponto são
a mesma operação. O espectro de um $\delta$ não tem zeros, e é por isso que a
volta é exacta; onde o espectro zera --- o vector todo-uns, por exemplo --- a
informação perde-se e a deconvolução deixa de existir. \emph{A completude do
\S anterior não usa nada disto}; a correspondência regista-se por ser da mesma
estrutura.

\medskip
\noindent\textbf{E o teste verifica as operações realizadas; não mede a
completude:}
fecho, tricotomia, compatibilidade da ordem com $+$ e $\times$, e o $n$
arquimediano exibido. \emph{Verificado:} \code{tests/ordenado.c} ($7{:}0$),
\code{tests/convolucao\_universal.js}.

\noindent\textbf{E são duas rotas, que não se substituem:}
\[
\boxed{
\begin{array}{lcl}
\text{cortes} &\longrightarrow& \text{corpo ordenado completo}
\ \longrightarrow\ \text{unicidade}\ \longrightarrow\ \mathbb{R}\\[4pt]
(a_k) \ \longrightarrow\ (p_k/q_k) &\longrightarrow&
\text{intervalos encaixados}\ \longrightarrow\ \text{corte}
\end{array}}
\]
A primeira é a \emph{identificação}; a segunda é o \emph{mecanismo}. E a
formulação exacta, que é uma escolha metodológica e não uma tese sobre
fundamentos --- outra construção pode perfeitamente axiomatizar $\mathbb{R}$ por
sucessões de fracção contínua:
\[
\boxed{\;\begin{gathered}
\text{neste paper a CF \textbf{não} é a definição axiomática de }\mathbb{R};\\
\text{é o \textbf{mecanismo de realização} dos cortes}
\end{gathered}\;}
\]
É também o estatuto que o \emph{Teoria} já lhe dá: o corpo dual e as duas leis são
o \textbf{fundamento}; a Möbius de Fibonacci, o gato e os convergentes são a
\textbf{representação} --- ferramenta, não fonte.

\begin{observacao}[a identificação com $\mathbb{R}$]
Que um corpo ordenado completo seja isomorfo a $\mathbb{R}$, e único a menos de
isomorfismo, é teorema clássico, e cita-se como tal. \textbf{Mas a passagem a
$\mathbb{R}$ não é a identificação posterior de um subconjunto conveniente}: é
a realização da \emph{face alcançada} da própria construção. A face binária
permanece finita e não ordenável como corpo; a face alcançada é aquela em que
os cortes se tornam estrutura ordenada e completa. \emph{São duas faces do
mesmo eixo, não duas teorias concorrentes.}
\end{observacao}

%======================================================================
\section{A continuidade: dado $S$ limitado, o supremo constrói-se}\label{sec:supremo}

\noindent Até aqui a direcção foi sempre a mesma --- \emph{dado o habitante $x$, a
classe racional dele tem $x$ por supremo} (\S\ref{sec:completude},
Teor.~\ref{thm:completude}). Isso constrói o \textbf{corte a partir do habitante}.
A continuidade é a \textbf{volta}, e é ela que separa $\mathbb{R}$ de $\mathbb{Q}$
como estrutura ordenada:
\[
\boxed{\;S\neq\emptyset\ \text{limitado superiormente}
\ \Longrightarrow\ \exists\,\sup S\ \textbf{dentro do objecto}\;}
\]
E aqui não é axioma acrescentado: o caminho do supremo \textbf{constrói-se}, bit a
bit, em inteiros.

\begin{definicao}[o caminho do supremo]\label{def:supremo}
Para $S\neq\emptyset$ limitado superiormente,
\[
m_k:=\max\bigl\{m\in\mathbb{Z}:\ \tfrac{m}{2^{k}}<x\ \text{para algum }x\in S\bigr\}
\]
--- bem definido, porque $S$ não é vazio (há $m$ pequeno que serve) e é limitado
(há $m$ grande que não serve).
\end{definicao}

\begin{teorema}[$m_k$ é um caminho, e é o supremo]\label{thm:supremo}
\emph{(i) É um caminho}, e é o \textbf{passo} que o diz:
\[
\boxed{\;m_{k+1}\in\{2m_k,\ 2m_k+1\}\;}
\]
--- os $m_k$ descem a árvore binária \emph{sem saltar}, logo são um habitante do
objecto e não uma lista solta de inteiros.

\emph{(ii) É majorante:} $m_k+1$ já não serve, logo nenhum $x\in S$ ultrapassa o
caminho.

\emph{(iii) E é o \textbf{menor}} --- é esta a cláusula que trabalha, porque (ii)
sozinha daria um majorante qualquer. \textbf{E ela não se varre: prova-se}, e a prova é
a \emph{monotonia}. Ponha-se
\[
E_k(m):\quad \exists\,x\in S:\ \tfrac{m}{2^{k}}<x .
\]
Então $E_k$ é \textbf{decrescente} em $m$ --- se $m'\leq m$ e $E_k(m)$, o mesmo $x$
serve --- e daí, numa linha,
\[
\boxed{\;m'<m_k\ \Longrightarrow\ \tfrac{m'}{2^{k}}<\tfrac{m_k}{2^{k}}<x
\ \Longrightarrow\ E_k(m')\;}
\]
para \emph{todo} $m'$ abaixo, sem excepção e sem varredura: nenhum deles é majorante.
\emph{O que se verifica computacionalmente é o \textbf{passo}} $E_k(m+1)\Rightarrow
E_k(m)$, de onde a indução dá a cláusula inteira --- e não uma amostra dos $m'$, que
seria substituir a prova por uma janela. $\square$
\end{teorema}

\begin{corolario}[e o ÍNFIMO é a outra metade --- quem os troca é o
SWAP]\label{cor:infimo}
O Teor.~\ref{thm:supremo} constrói o supremo, e sozinho é meia lei: \emph{o par é
$(\sup,\inf)$}. A outra metade não pede construção nova --- deriva-se, e o que a produz é
uma peça que já está na fundação.

\textbf{E não é a involução.} Em $x>0$, $\nu(x)=-1/x$ tem derivada $1/x^{2}>0$, logo é
\emph{crescente}: leva supremo em \textbf{supremo}, e não em ínfimo. Quem inverte a ordem é
o \textbf{swap} da Def.~\ref{def:lei0},
\[
S=\begin{psmallmatrix}0&1\\1&0\end{psmallmatrix},\qquad [p:q]\mapsto[q:p],\qquad
\det S=-1,
\]
que em $x>0$ é $x\mapsto1/x$, \emph{decrescente}. Donde, para $A\subset\mathbb{Q}^{+}$
limitado e afastado de zero,
\[
\boxed{\;\inf A=\frac{1}{\sup\,S(A)},\qquad \sup A=\frac{1}{\inf\,S(A)}\;}
\]
--- e é o $\det=-1$ que diz que ele \textbf{troca}, o mesmo $-1$ do gato. \emph{O supremo e
o ínfimo não são duas construções: são uma, vista dos dois lados do swap.} $\square$
\end{corolario}
\colunas{o ínfimo, e quem troca o par}
        {$(\sup,\inf)$ é o par; quem o troca é o swap ($\det=-1$), e a involução $\nu$ NÃO
--- ela preserva o supremo}
        {\code{tests/transformada\_universal.c} \S T17: em $84$ conjuntos de racionais, «o
supremo do transformado é o ínfimo do original» em $84$, o mesmo nos dois sentidos, e
$\nu(x)=-1/x$ a PRESERVAR o supremo em $84$ --- que é o contraste que impede trocar de
operador; tudo por produto cruzado de inteiros}

\begin{proposicao}[e é aqui que $\mathbb{Q}$ falha, com testemunha]\label{prop:q-falha}
Para $S=\{r\in\mathbb{Q},\,r>0:\ r^{2}<2\}$ nenhum $p/q$ é o supremo, e a
testemunha é a \emph{mesma} nos dois casos:
\[
\boxed{\;(p,q)\longmapsto(3p+4q,\ 2p+3q),
\qquad \det\begin{psmallmatrix}3&4\\2&3\end{psmallmatrix}=1\;}
\]
--- a \textbf{unidade fundamental de $\mathbb{Z}[\sqrt2]$ ao quadrado},
$(1+\sqrt2)^{2}=3+2\sqrt2$. O que faz dela testemunha são duas identidades
inteiras, e são precisas as duas:
\[
\begin{array}{ll}
(3p+4q)^{2}-2(2p+3q)^{2}=p^{2}-2q^{2}
 & \textbf{preserva} \text{ a forma, logo o \textbf{lado}}\\[3pt]
(3p+4q)q-p(2p+3q)=2\,(2q^{2}-p^{2})
 & \text{e \textbf{move} na direcção desse lado}
\end{array}
\]
Juntas: quem está \emph{abaixo} sobe continuando abaixo --- fica em $S$, e é maior,
logo $p/q$ não era majorante; quem está \emph{acima} desce continuando acima ---
continua majorante, e é menor, logo $p/q$ não era o menor. E não há empate, pois
$p^{2}=2q^{2}$ é a descida proibida (Cor.~\ref{cor:habitante}). $\square$
\end{proposicao}

\begin{observacao}[e a mediante simples não serve]\label{obs:mediante}
A mediante de Farey $(p+2q)/(p+q)$ satisfaz
$(p+2q)^{2}-2(p+q)^{2}=-(p^{2}-2q^{2})$ e portanto \textbf{troca} o sinal: ela
\emph{atravessa} o corte e cai do lado oposto, em vez de ficar entre $p/q$ e ele.
A identidade é verdadeira e a testemunha seria falsa --- é preciso a matriz de
$\det=1$ acima, que é a mediante \emph{aplicada duas vezes}.
\end{observacao}
\colunas{a continuidade como propriedade do menor majorante}
        {o caminho do supremo, construído bit a bit em inteiros}
        {\code{tests/supremo.c} ($6{:}0$): $104$ passos por duas rotas concordantes;
$64$ níveis com a cláusula (a) e com o \textbf{passo} da monotonia sem uma violação em
$761\,859$ verificações --- e não numa janela dos últimos $40$, que era o que aqui
estava com a asserção a dizer «todos»; o gume com uma condição não monótona a violar o
passo as $429$ vezes previstas; $5490$ candidatos com as duas identidades da
testemunha e zero empates; e o supremo de $\{r^2<2\}$ dá $1\,482\,910/2^{20}$}

\medskip
\noindent\textbf{E é por isto que o sistema atravessa os andares com erro zero.} O
habitante \textbf{é} o caminho --- não há valor por trás à espera de ser
aproximado, logo não há erro a acumular: a ida e a volta são a involução, exacta
por construção (\S\ref{sec:dual}). A continuidade clássica precisa do limite
porque lá o ponto tem de \emph{aparecer} de fora; aqui o ponto \emph{é} a
sucessão dos passos, e o supremo é apenas mais um caminho que se desce.

%======================================================================
\section{E só então $\pi$}\label{sec:pi}

\begin{definicao}[dois predicados, e não um]\label{def:primo-ordem}
Seja $\mu_n\in\mathbb{Z}[x]$ o polinómio mínimo de $t_n=2\cos(\pi/n)$ sobre
$\mathbb{Q}$. Separam-se:
\[
\boxed{\;
\begin{array}{ll}
p\ \text{é \textbf{primo de realização} de } n
 & \iff\ \mu_n\ \text{tem raiz } \bar t_n\in\mathbb{F}_p\\[3pt]
(p,\bar t_n)\ \text{\textbf{realiza a ordem} } n
 & \iff\ \mu_n(\bar t_n)=0\ \text{e}\ \operatorname{ord}\bigl(C_n(\bar t_n)\bigr)=2n
\end{array}\;}
\]
O primeiro é condição sobre a \emph{existência} da realização modular do traço; o
segundo é propriedade do \textbf{par} $(p,\bar t_n)$ --- porque a companheira depende
da raiz escolhida --- e não decorre do primeiro. Os exemplos usados aqui cumprem ambos, e isso verifica-se:
\[
\begin{array}{lll}
n=4: & \mu_4=x^{2}-2,\ p=17 & \bar t_4\in\{11,6\}\\
n=6: & \mu_6=x^{2}-3,\ p=13 & \bar t_6\in\{9,4\}\\
n=3: & \mu_3=x-1 & \bar t_3=1\ \text{em toda a característica}
\end{array}
\]
Para $n=3$, $\mu_3=x-1$ e portanto \emph{toda} característica admite a raiz; mas
para a afirmação de ordem $6$ restringe-se às características em que $-I\neq I$,
em particular $p\neq2$ --- em característica $2$ a ordem colapsa. \textbf{Em geral},
nos casos usados $p$ é escolhido \emph{fora das características que degeneram a
relação algébrica relevante} --- em particular fora das que anulam o discriminante,
onde as folhas colapsam (Cor.~\ref{cor:dourada-completa}).
\end{definicao}

\begin{definicao}[$t_n$, pela ordem --- sem $\pi$]\label{def:tn}
$t_n$ é o traço da companheira
$C=\begin{psmallmatrix}t&-1\\1&0\end{psmallmatrix}$, de determinante $1$, tal que
\[
\boxed{\;C^{n}=-I\quad\text{e}\quad C^{k}\neq\pm I\ \text{ para } 0<k<n\;}
\]
--- pura condição de \textbf{ordem}, sem nenhuma constante. \textbf{Mas a ordem
não escolhe o ramo}: há mais de uma companheira de ordem $2n$, com traços
distintos. Fixa-se então o representante:
\[
\boxed{\;t_n:=\text{o \textbf{maior} traço real entre as companheiras de ordem }2n\;}
\]
--- e assim $t_n$ fica determinado sem que nenhuma constante entre.
\end{definicao}

\noindent E \emph{só então}, como \textbf{identificação posterior}, escreve-se
\[
t_n=2\cos\frac{\pi}{n},
\]
que é uma identificação e não a definição --- sem isso, $\pi$ entraria na
definição de $t_n$ e a cadeia
$\text{ordem}\to t_n\to A_n^{\rm in},A_n^{\rm circ}\to\pi_n\to\pi$
seria circular. E a companheira
$C_n=\begin{psmallmatrix}t_n&-1\\1&0\end{psmallmatrix}$ tem ordem $2n$. \textbf{Mas $t_n$ não é inteiro} --- $t_4=\sqrt2$, $t_6=\sqrt3$ ---,
e é aqui que está a passagem que dá o título a esta secção: nos modelos finitos
usados para a realização, \emph{o traço irracional é substituído pela raiz da
classe modular correspondente} --- uma raiz inteira, em $\mathbb{F}_p$, do mesmo
polinómio mínimo. E não se confundem os dois objectos:
\[
t_n=2\cos\frac{\pi}{n}\in\mathbb{R}
\qquad\text{contra}\qquad
\bar t_n\in\mathbb{F}_p ,
\]
sendo o segundo uma \textbf{realização modular de uma relação algébrica}, e não
«o mesmo número». E o alcance disso limita-se, porque as áreas usam
$\sqrt{4-t_n^{2}}$ e a passagem da raiz real para uma raiz em $\mathbb{F}_p$ exigiria
estrutura adicional:
\[
\boxed{\;\begin{gathered}
\text{a realização modular fecha a \textbf{relação algébrica do traço};}\\
\text{as áreas e }\pi_n\text{ permanecem na realização \textbf{real}}
\end{gathered}\;}
\]

\begin{teorema}\label{thm:pi}
\emph{(i) os elíticos.} Com $t_n=2\cos(\pi/n)$ realizado no primo da ordem,
\[
\begin{array}{llll}
n=3: & \mu_3=x-1 & \bar t_3=1\in\mathbb{F}_p & \operatorname{ord}C_3=6\\
n=4: & \mu_4=x^{2}-2 & \bar t_4=11\in\mathbb{F}_{17},\ \mu_4(11)\equiv0
     & \operatorname{ord}C_4=8\\
n=6: & \mu_6=x^{2}-3 & \bar t_6=9\in\mathbb{F}_{13},\ \mu_6(9)\equiv0
     & \operatorname{ord}C_6=12
\end{array}
\]
--- ordem exactamente $2n$ em cada caso. As duas raízes conjugadas realizam a
\emph{mesma relação algébrica} e, nestes exemplos, ambas produzem ordem $2n$ ($6$
em $\mathbb{F}_{17}$, $4$ em $\mathbb{F}_{13}$); a construção fixa uma delas como
representante do ramo escolhido.

\emph{(ii) os hiperbólicos} são os metais, o motor dos
\S\S\ref{sec:dual}--\ref{sec:todos}.

\emph{(iii) o parabólico.} Para $t=2$ o polinómio é $x^{2}-2x+1=(x-1)^{2}$, logo
$C=I+N$ com $N^{2}=0$; em característica $p$ ímpar $C^{k}=I+kN\neq I$ para
$1\leq k<p$ e $C^{p}=I$, portanto
\[
\boxed{\;\operatorname{ord}C=p:\ \text{o fecho parabólico é governado pela
\textbf{característica}}\;}
\]
e \emph{não} coincide genericamente com a ordem elíptica $2n$. Note-se que $p$
\textbf{pode} dividir $2n$ --- $n=3$ com $p=3$ dá $3\mid6$, e $n=5$ com $p=5$ dá
$5\mid10$ ---, pelo que a separação dos dois regimes é pela \emph{natureza do
fecho}, e não por uma divisibilidade universal. $\square$
\end{teorema}

\begin{definicao}[$\pi_n$ --- e ele é um intervalo]\label{def:pin}
Do traço $t_n=2\cos(\pi/n)$ do andar saem as duas áreas, ambas \textbf{algébricas
em $t_n$}:
\[
\boxed{\;
A_n^{\rm in}=\frac{n\,t_n}{4}\sqrt{4-t_n^{2}}
\qquad\text{e}\qquad
A_n^{\rm circ}=\frac{n\sqrt{4-t_n^{2}}}{t_n}\;}
\]
--- a do polígono inscrito e a do circunscrito, no círculo unitário, com notação
distinta da companheira $C_n$ do Teor.~\ref{thm:pi}, que é outra coisa. Elas
satisfazem, \emph{sem que $\pi$ entre}:
\[
A_n^{\rm in}<A_n^{\rm circ},\qquad
A_n^{\rm in}\uparrow,\qquad A_n^{\rm circ}\downarrow,\qquad
\abs{\pi_n}\to0 .
\]
Nos andares $3,4,6$ os quadrados são racionais --- $(A_3^{\rm in})^{2}=27/16$ e
$(A_3^{\rm circ})^{2}=27$; $A_4^{\rm in}=2$ e $A_4^{\rm circ}=4$;
$(A_6^{\rm in})^{2}=27/4$ e $(A_6^{\rm circ})^{2}=12$ --- e a duplicação de
Arquimedes $(A_{2n}^{\rm in})^{2}=A_n^{\rm in}A_n^{\rm circ}$ fecha exacta, dando
$(A_8^{\rm in})^{2}=8$. O encaixe $A_8^{\rm in}<333/106<22/7<A_8^{\rm circ}$
decide-se por comparação de inteiros.
\end{definicao}

\medskip
\noindent E daqui sai a definição que faltava, e ela usa o mecanismo do próprio
paper:
\[
\boxed{\;\pi_n:=\bigl[A_n^{\rm in},\,A_n^{\rm circ}\bigr],
\qquad
\pi_{2n}\subset\pi_n,
\qquad
\abs{\pi_n}\longrightarrow0\;}
\]
--- \textbf{$\pi_n$ é o intervalo do andar}, com as duas bordas exactas nele. O
encaixe está demonstrado ao longo da \emph{duplicação} $n\mapsto2n$
(Def.~\ref{def:pin}), pelo que a família encaixada é a cadeia
$\pi_{n_0}\supset\pi_{2n_0}\supset\pi_{4n_0}\supset\cdots$. Donde
\[
\boxed{\;\pi:=\bigcap_{j\geq0}\pi_{2^{j}n_0}\ \text{---\ o \textbf{único habitante}
do corte que a cadeia de duplicações determina.}\;}
\]
Não é constante importada nem aproximação numérica: é um corte, produzido pelo
mesmo mecanismo de \S\ref{sec:completude}, com bordas algébricas em cada andar.
\emph{O que é exacto por andar são as bordas}; $\pi$ continua a ser o corte. O que o quadro
dispensa é \emph{importar} uma constante: cada andar apresenta o seu.

\noindent Nenhuma constante foi importada: \textbf{cada andar apresenta a sua}.
E na realização modular adoptada, a ordem do operador elíptico fecha em
características adequadas, ao passo que a do parabólico é $p$ --- a mesma
fronteira do eixo, agora como propriedade da representação escolhida do
membro-limite, e não como propriedade universal da geometria circular. \emph{Assim o círculo não é usado para construir a reta: nesta realização, ele
reaparece como gume geométrico} --- e a palavra é \emph{reaparece}, não
\emph{portanto}: não há aqui um teorema causal entre a completude e o círculo.
\colunas{$\pi$ é o membro-limite da família (\code{thm:pi-familia})}
        {$\pi_k$ exacto por andar, realizado nos primos das ordens}
        {\code{geometria\_real.c} \S L10; \code{pi\_familia.js} ($7{:}0$)}

%======================================================================



%======================================================================
\section{E o que vem a seguir são as duas peças do ALCANCE que faltavam}\label{sec:todo-r}

\noindent\emph{Também estas vieram do centro}, e a segunda trazia o seu próprio atestado:
o \S do Teorema Central abre a dizer que a tríade Hurwitz--Gentil--Lebesgue \textbf{não
está na cadeia principal da construção}. Uma secção que declara não fundar nada é, por
definição, realização --- e realização do \emph{limite}, que é matéria deste lado.

E a primeira é o alcance dito por extenso: que a construção não se restringe à família
metálica, e que \emph{todo} $\mathbb{R}$ se atinge. O centro prova que a régua existe;
aqui mostra-se que ela chega a toda a parte.

\section{O Teorema Central funda o limite}\label{sec:central}

\noindent\textbf{E o estatuto desta secção diz-se à entrada, porque ele mudou.} A
tríade Hurwitz--Gentil--Lebesgue \emph{não} está na cadeia principal da construção ---
essa é a do \S\ref{sec:operacao}, $L_0,\ldots,L_7\to\star\to\tau\to E_{+}\oplus E_{-}\to
\operatorname{Dir},\operatorname{Cruz}\to B_8$, e nela não entra nenhum nome de fora.
\[
\boxed{\;\text{Hurwitz / Gentil / Lebesgue}\ \not\to\ \text{construção};
\qquad
\text{construção}\ \to\ \text{resultado}\ \to\ \text{comparação}\;}
\]
O que segue é \textbf{comparação com o objecto já construído}: três nomes clássicos
postos ao lado dele para dizer o que coincide e o que não, depois de a construção estar
feita e medida. \emph{Nenhuma peça do \code{analitico sec:supremo} ou do \code{analitico sec:completude}
depende do que se diz aqui.}

\begin{teorema}[contar pelo domínio e medir por níveis somam ao rectângulo]\label{thm:central}
\[
\boxed{
\begin{array}{ll}
\textbf{Hurwitz} & \text{a contagem do domínio \emph{assim denominada no Algébrico}}\\
\textbf{Gentil} & \text{a \textbf{identidade discreta} que estabelece a compatibilidade}\\
\textbf{Lebesgue} & \text{a \textbf{interpretação estrutural} da medida da imagem}
\end{array}}
\]
A soma de Gentil, em inteiros,
\[
0\leq x_n\leq q\ \Longrightarrow\
\sum_n x_n \;+\; \sum_{v=1}^{q} \#\{x_n<v\} \;=\; N\!\cdot\!q ,
\]
é --- \emph{pela identidade de Gentil estabelecida no Corpo Algébrico} --- a forma
discreta de $\int f+\int f^{-1}=b\,f(b)-a\,f(a)$: contar pelo \emph{domínio} e
medir por \emph{níveis} somam ao rectângulo. E o \emph{layer-cake}
$\sum_v\#\{x_n\geq v\}=\sum_n x_n$ fecha \textbf{sem esperar o limite}, com
resíduo zero em cada andar. $\square$
\end{teorema}

\noindent\textbf{E onde cada um entra, para que não fique dúvida:} \emph{Gentil} é
a identidade discreta, e é a única das três que se \emph{usa} na conta;
\emph{Lebesgue} entra como interpretação estrutural, e não como ingrediente. E o
que esta secção faz é \emph{fundar} o limite como ponto fixo --- a
\textbf{continuidade} propriamente dita, a existência do supremo, constrói-se no
\code{analitico sec:supremo} e não aqui:
\[
\boxed{\;\begin{gathered}
\text{\emph{nesta realização}, o limite lê-se como o \textbf{ponto fixo}}\\
\text{da compatibilidade entre as três contagens}
\end{gathered}\;}
\]

\begin{corolario}[a tricotomia distribui-se pelas três leituras]\label{cor:distribui}
\textbf{TERMINA} e \textbf{RODA} são o lado \emph{contável}, que Hurwitz conta:
fixado $D$, os racionais de denominador $\leq D$ são em número finito, e a quota
que ocupam nos $2^{K}$ intervalos de profundidade $K$ decresce estritamente com
$K$ \emph{na discretização diádica utilizada pelo medidor} --- a borda é magra.
\textbf{FOGE} é o que sobra. $\square$
\end{corolario}
\colunas{a tríade Hurwitz--Gentil--Lebesgue como eixo do Teorema Central}
        {a fundação do corte como ponto fixo, e a distribuição da tricotomia}
        {\code{geometria\_real.c} \S L11d: $60$ casos de Gentil e do layer-cake;
$22$ racionais, quota de $22/64$ a $22/4096$}

\noindent\emph{E o par cone/espiral tem a mesma forma da soma reversível}: duas
leituras do mesmo objecto, que têm de fechar uma na outra.

\medskip
\noindent\textbf{E o que Gentil preserva é uma norma, mas não a mesma.} É a metade que
fecha a tríade, e sem ela «Gentil casa» fica sem conteúdo operacional. O produto dele é
\emph{componente a componente}, $(a,b)\ast(c,d)=(ac,\,-bd)$, e a norma que ele preserva
é o \textbf{produto das coordenadas} --- o determinante de $\operatorname{diag}(x)$, e
não a euclidiana:
\[
\boxed{\;N(x\ast y)=\prod_i\abs{x_iy_i}=\prod_i\abs{x_i}\cdot\prod_i\abs{y_i}
=N(x)\,N(y)\qquad\text{em \textbf{toda} a dimensão}\;}
\]
Donde o par, e ele explica o tecto:
\[
\boxed{
\begin{array}{lll}
\textbf{Hurwitz} & \text{norma euclidiana}\ +\ \text{divisão} & \to\ \text{tecto em }8\\[3pt]
\textbf{Gentil} & \text{norma-produto, \emph{sem} divisão} & \to\ \textbf{sem tecto}
\end{array}}
\]
\textbf{O tecto de $8$ não é do objecto nem do número: é do par (norma, divisão).} Trocar
a norma tira o tecto e paga a divisão --- e o preço exibe-se: com uma coordenada nula o
produto tem norma zero sem que nenhum factor seja nulo, logo há \emph{divisores de zero},
que é exactamente o que a hipótese de Hurwitz proíbe.
\medskip
\noindent\textbf{E as operações, andar a andar.} A soma é grupo abeliano em \emph{toda} a
torre, dos dois lados: \textbf{tudo o que se perde perde-se do lado do produto}. Esse
perde um axioma por degrau --- comuta até $2$, associa até $4$, sem divisores de zero até
$8$ ---, e a \textbf{distributividade} é a fronteira dura: Cayley--Dickson tem-na em todos
os andares, e o lado de Gentil não. Donde os nomes, que dizem o que se pode fazer:
\[
\boxed{
\begin{array}{lll}
1,\ 2 & \mathbb{R},\ \mathbb{C} & \textbf{corpo}\\[2pt]
4,\ 8 & \mathbb{H},\ \mathbb{O} & \text{álgebra de divisão}\\[2pt]
16 & \text{sedeniões} & \text{anel --- há divisores de zero}\\[2pt]
3 & \text{Gentil} & \text{\emph{nem} anel --- falha a distributividade}
\end{array}}
\]
--- e \emph{«quase um corpo» não existe}: falhar um axioma não é gradação, é outro
objecto. \emph{Verificado:} \code{tests/torre\_corpo.c} ($4{:}0$).

\colunas{a tríade, e o tecto como propriedade da norma}
        {a norma que Gentil preserva, e o preço que ela cobra}
        {\code{tests/hurwitz.c} \S H8: a norma-produto multiplicativa em $2,4,8,16,32$,
e o divisor de zero exibido}

\medskip
\noindent\textbf{E por que é que a soma tem tecto $8$: está na base.} De
$N(xy)=N(x)N(y)$ com $N=\sum x_i^{2}$, polarizando em $y$, sai
\[
\boxed{\;\langle xy,\,xz\rangle=N(x)\,\langle y,z\rangle\;}
\]
--- \emph{o produto preserva a ortogonalidade}. Numa base ortonormal $e_0=1,\ldots,
e_{n-1}$ isso força $e_i^{2}=-1$ e $e_ie_j=-e_je_i$, que são as relações de Clifford:
\textbf{é a base que as carrega}, e é ela que a base ortonormal do \S\ref{sec:orto}
realiza.

\noindent\textbf{Mas a saturação não é o tecto --- e a torre não pára.} A tabela
$e_ie_j$ \emph{satura} a base em todos os andares medidos, até $64$: os produtos caem
todos nela, e cada andar sai \textbf{completo}. Saturar é o que manda \emph{dobrar}.
\[
\boxed{
\begin{array}{ll}
\text{preencher os andares \textbf{na ordem}, completos} & \text{é o que a construção
faz --- e basta}\\[3pt]
\text{o «tecto» de }8 & \text{é \textbf{uma exigência} posta de fora}
\end{array}}
\]
A exigência é que a norma \emph{euclidiana} seja multiplicativa e que não haja divisores
de zero. Ela não é a torre: é um pedido que se lhe faz. E cada propriedade que cai pelo
caminho --- a ordem em $2$, a comutatividade em $4$, a associatividade em $8$ --- é
\textbf{uma exigência a cair, e não o objecto a falhar}. O outro lado do par mostra-o:
com a norma de Gentil não há tecto nenhum.

\noindent E o que aparece em $16$, $x,y\neq0$ com $xy=0$, é \textbf{a mesma factorização
do zero} do Teor.~\ref{thm:divzero}: lá o polinómio a factorizar, aqui a álgebra a
passar dos oito, e nos dois casos a norma multiplicativa é a primeira exigência a cair.
\colunas{a base ortonormal das oito leis (\code{teoria thm:ortonormal})}
        {a polarização, os andares completos, e o tecto como exigência}
        {\code{tests/hurwitz.c} \S H9--H10: a base satura e os produtos caem nela nos
SETE andares até $64$, com a norma euclidiana a sobreviver em quatro e o divisor de zero
de dimensão $16$ exibido}

%======================================================================


\section{Todo $\mathbb{R}$ --- e o ouro é o extremo}\label{sec:todos}

A realização não se restringe à família metálica.

\begin{definicao}[e o domínio fecha-se dos dois lados]\label{def:real}
No corpo ordenado completo identificado pelo \emph{Corpo Algébrico}, cada
elemento admite a expansão canónica por fracção contínua. Os convergentes são
$p_k/q_k$ com
\[
q_{-1}=0,\qquad q_0=1,\qquad q_k=a_kq_{k-1}+q_{k-2},
\]
e a correspondência é esta, com as duas metades ditas:
\[
\boxed{
\begin{array}{ll}
\text{cada IRRACIONAL} & \text{uma sucessão INFINITA }(a_k)_{k\geq1},\ a_k\geq1\\[2pt]
\text{cada RACIONAL}   & \text{o caso em que o processo TERMINA}
\end{array}}
\]
Dizer só a primeira linha deixaria $\mathbb{R}$ a descoberto, e a afirmação
deste paper é sobre \textbf{todo} $\mathbb{R}$. Para $a/b$ o algoritmo termina com
resto zero, o último convergente dá $a/b$ \textbf{exacto} --- não aproximado
---, e a ambiguidade final $[a_0;\ldots,a_n]=[a_0;\ldots,a_n{-}1,1]$ resolve-se
por convenção canónica.
\end{definicao}

\begin{corolario}[e por isso a ENTREGA é a palavra, e não um
decimal]\label{cor:entrega}
Se a palavra $(a_k)$ \emph{é} o objecto, então nada obriga um programa a
entregar o seu resultado em decimal --- e entregá-lo assim é converter o
objecto numa das suas leituras, perdendo as outras. O que sai são os
$a_k$, inteiros; quem recebe reconstrói o corte que precisar:
\[
\text{texto}\ \longrightarrow\ \frac{p}{q}\ \text{exacto}\ \longrightarrow\
[a_0;a_1,\ldots,a_n]\ \longrightarrow\ \frac{p}{q}\ \longrightarrow\
\text{decimal, convergentes, ou o que for.}
\]
As duas pontas merecem ser ditas porque nenhuma é óbvia. \textbf{À entrada},
um decimal escrito não é um número aproximado: \code{3.14159} \emph{é}
$314159/100000$, e o que o estraga é convertê-lo a vírgula flutuante --- em
base dois $0{,}1$ não é $1/10$, e não é por imprecisão do formato mas porque
$10\nmid 2^k$. \textbf{À saída}, a mesma palavra dá \emph{todas} as
representações, e não uma: de \code{3.14159} saem $22/7$ e $355/113$ sem
ninguém os pedir, porque os convergentes são o que a recorrência do
$q_k$ produz.

\textbf{E o pipe fecha do outro lado com o MMC.} À entrada, cada texto lido isolado dá o
seu $p/q$ e um conjunto fica com denominadores diferentes --- o que obriga a cruzar
fracções a cada operação. O \textbf{mínimo múltiplo comum dos denominadores} é a
\emph{unidade} do conjunto, e leva-o todo a $\mathbb{Z}$ de uma vez:
\[
\text{textos}\ \longrightarrow\ \underbrace{\operatorname{mmc}(q_i)}_{\text{a unidade}}\
\longrightarrow\ \text{inteiros}\ \longrightarrow\ \text{operar}\ \longrightarrow\
[a_0;\ldots,a_n]\ \longrightarrow\ \text{o cliente}
\]
\emph{e entre a unidade e a palavra não há vírgula em sítio nenhum}. Normalizar, aqui, não
é dividir por coisa nenhuma --- é \textbf{escolher a régua}, e o que a autoriza é as
perguntas serem homogéneas: escalar por $\lambda$ multiplica os dois lados de uma
comparação, e Lagrange é homogénea de grau~4.

\emph{Verificado:} \code{tests/entrega.c} ($11{:}0$). A ida e volta pela
palavra é exacta nos $10\,000$ decimais de três casas e a comparação é por
produto cruzado --- a \emph{mesma} fracção, não uma que difere abaixo de uma
régua; pelo \code{double}, com os mesmos números, só $80$ voltam, e
é esse par de números lado a lado que dá conteúdo ao primeiro. Mede-se
ainda que a palavra \textbf{é} Euclides --- os seus quocientes são os do
algoritmo do mdc, um a um e na mesma ordem, em $3600$ pares --- e que os
dois tectos têm onde bater, com a palavra mais comprida a ser a do ouro:
dois Fibonacci consecutivos dão a fracção contínua toda de uns, que é a
Teorema~\ref{thm:ouro} lido como \emph{pior caso da representação}.
\end{corolario}

\begin{teorema}[a construção corre sobre quocientes arbitrários]\label{thm:todos}
Para uma sucessão \textbf{qualquer} --- não só a periódica ---:
$p_kq_{k-1}-p_{k-1}q_k=(-1)^{k-1}$, logo $\abs{\det}=1$ em todos os passos; os
convergentes alternam de lado; os intervalos formam uma cadeia
encaixada; e a monotonia lê-se directamente da recorrência:
\[
q_1\geq q_0,
\qquad
q_{k+1}=a_{k+1}q_k+q_{k-1}>q_k\quad(k\geq1)
\]
--- o estrito começa em $k\geq1$ porque, com $q_{-1}=0$, $q_0=1$ e $a_1=1$, tem-se
$q_1=q_0$.
\emph{O metálico é o caso de quociente parcial constante --- periódico, mas não
o único periódico, e muito menos o único caso.} $\square$
\end{teorema}

\begin{teorema}[a extremalidade da régua: $\varphi$ minimiza o crescimento]\label{thm:ouro}
Com $a_k\geq1$,
\[
\boxed{\;q_k\geq F_{k+1}\ \ (k\geq0),\qquad
\text{com igualdade em TODO }k\ \text{só na sucessão de uns,}\;}
\]
na convenção padrão $F_0=0$, $F_1=1$. \textbf{É esta a forma afiada}: com
$q_{-1}=0$ e $q_0=1$, a sucessão de uns dá exactamente $q_k=F_{k+1}$, e não apenas
$q_k\geq F_k$. Donde a resolução da régua,
\[
\abs{c_{k+1}-c_k}=\frac{1}{q_kq_{k+1}}\ \leq\ \frac{1}{F_{k+1}F_{k+2}}
\]
--- \textbf{a taxa de Fibonacci fornece o pior caso algébrico para o crescimento
dos denominadores} entre as expansões infinitas com $a_k\geq1$. A família metálica não é o escopo
--- é o caso \textbf{extremo}. $\square$

\medskip
\noindent\emph{E o escopo da palavra «lento» diz-se}, para não privilegiar
ninguém: o que está demonstrado é $q_k\geq F_{k+1}$ \emph{para cada $k$}, isto é,
$\varphi=[1;1,1,\ldots]$ \textbf{minimiza ponto a ponto o crescimento dos
denominadores $q_k$} entre as expansões infinitas com $a_k\geq1$ --- e, como
consequência assintótica, $F_k\sim\varphi^{k}/\sqrt5$. Não é uma
propriedade de $\varphi$ contra os outros reais em geral --- os racionais nem
sequer têm expansão infinita ---, é a \textbf{extremalidade da própria régua}:
o pior caso do mecanismo, e por isso o que o limita.
\end{teorema}
\colunas{a torre, o dual, e a família metálica}
        {a construção para TODOS os reais, com o ouro a limitar}
        {\code{geometria\_real.c} \S L8: $220$ sucessões, $3651$ pares, $0$
excepções; \S L8b: $3660$ racionais, terminação e exactidão em todos}

\begin{observacao}
A igualdade num índice isolado não caracteriza a sucessão de uns: a caracterização
é pela igualdade \emph{simultânea em todos} os índices.
\end{observacao}

\begin{teorema}[Série geométrica em $\mathbb{Q}(m\sqrt D)$]\label{thm:serie-quadratica}
Sejam $D,m,a,b\in\mathbb{Q}$, com $D$ livre de quadrados,
$m\neq0$, e seja
\[
    \alpha=a+b\,m\sqrt D
    \in\mathbb{Q}(m\sqrt D).
\]
Suponha
\[
    \alpha\neq1,
    \qquad
    |\alpha|>1
\]
na imersão real considerada. Então
\[
    \sum_{n=1}^{\infty}\alpha^{-n}
    =
    \frac{1}{\alpha-1},
\]
e, sempre que $(a-1)^2-b^2m^2D\neq0$,
\[
    \boxed{
    \sum_{n=1}^{\infty}\alpha^{-n}
    =
    \frac{a-1-bm\sqrt D}
         {(a-1)^2-b^2m^2D}
    }.
\]
Toda a operação algébrica ocorre exactamente em
$\mathbb{Q}(m\sqrt D)$, sem ponto flutuante.
\emph{Em particular}, com $D=5$, $m=1$ e $\alpha=\varphi=(1+\sqrt5)/2$ satisfazendo
$\varphi^2=\varphi+1$, obtém-se $\varphi(\varphi-1)=1$, logo $1/(\varphi-1)=\varphi$ e
$\sum_{n\geq1}\varphi^{-n}=\varphi$ --- o caso dourado é \emph{instância} deste teorema,
não fundamento separado.
\end{teorema}
\begin{proof}
Escreva
\[
    \alpha=a+b\,m\sqrt D.
\]
O conjugado quadrático é
\[
    \alpha^\ast=a-b\,m\sqrt D.
\]
Logo
\[
    (\alpha-1)(\alpha^\ast-1)
    =
    (a-1)^2-b^2m^2D.
\]
Se esse produto é não nulo, a inversa existe e é
\[
    \frac1{\alpha-1}
    =
    \frac{\alpha^\ast-1}
         {(\alpha-1)(\alpha^\ast-1)}
    =
    \frac{a-1-bm\sqrt D}
         {(a-1)^2-b^2m^2D}.
\]

Para $N\geq1$, a soma parcial é a progressão geométrica
\[
    S_N
    =
    \sum_{n=1}^{N}\alpha^{-n}
    =
    \frac{\alpha^{-1}(1-\alpha^{-N})}
         {1-\alpha^{-1}}.
\]
Como
\[
    \frac{\alpha^{-1}}
         {1-\alpha^{-1}}
    =
    \frac1{\alpha-1},
\]
obtemos
\[
    S_N
    =
    \frac{1-\alpha^{-N}}
         {\alpha-1}.
\]

A hipótese $|\alpha|>1$ implica
\[
    \lim_{N\to\infty}\alpha^{-N}=0.
\]
Portanto
\[
    \boxed{
    \sum_{n=1}^{\infty}\alpha^{-n}
    =
    \frac1{\alpha-1}
    }.
\]

\medskip\noindent\emph{Caso $\varphi$.}
Com $D=5$, $m=1$ e $\alpha=\varphi=(1+\sqrt5)/2$, a relação $\varphi^2=\varphi+1$ dá
$\varphi=1+1/\varphi$, logo $1/\varphi=\varphi-1$ e $\varphi(\varphi-1)=1$, donde
$1/(\varphi-1)=\varphi$. A convergência usa só $0<\varphi^{-1}<1$.
\end{proof}
\colunas{a série geométrica em $\mathbb{Q}(m\sqrt D)$}
        {$\sum_{n\geq1}\alpha^{-n}=1/(\alpha-1)$; $\varphi$ incluído no teorema}
        {\code{tests/serie\_phi.c}; e na transformada em \code{tests/universal.c} \S U5}

\begin{teorema}[AS DUAS PROPRIEDADES DE $\mathbb{Q}(m\sqrt D)$, e a passagem
PRIMO~$\leftrightarrow$~IRRACIONAL é a passagem DISCRETO~$\leftrightarrow$~CONTÍNUO]%
\label{thm:duas-propriedades}
O Teor.~\ref{thm:serie-quadratica} opera em $\mathbb{Q}(m\sqrt D)$ sem dizer o que sustenta
o espaço. São \textbf{duas} propriedades, e cada uma responde a uma pergunta diferente:
\[
\boxed{
\begin{array}{ll}
\textbf{(P1)} & \sqrt D\in\mathbb{Q}\iff D\ \text{é QUADRADO}\\[3pt]
\textbf{(P2)} & \sqrt D\notin\mathbb{Q}\ \text{e}\ m\neq0\ \Longrightarrow\
m\sqrt D\notin\mathbb{Q}
\end{array}}
\]
\textbf{(P1) diz que o espaço EXISTE} --- com $D$ não quadrado há alguma coisa fora de
$\mathbb{Q}$ para adjuntar, e \emph{só os quadrados têm raiz racional}. \textbf{(P2) diz
que o espaço tem RÉGUA} --- medir com $m\sqrt D$ em vez de $\sqrt D$ não devolve nada a
$\mathbb{Q}$, e é isto que autoriza o $m$ no nome do corpo: \emph{a mesma recta, com uma
régua $m$}. A prova de (P2) é uma divisão: se $m\sqrt D=a/b$ então $\sqrt D=a/(bm)$, e
$bm$ é inteiro não nulo.

\textbf{(1) E a passagem PRIMO $\to$ IRRACIONAL é a mesma indivisibilidade, lida dos dois
lados.} Para $p$ primo, $\sqrt p$ é irracional, e o que fecha a prova é o \textbf{lema de
Euclides} --- $p\mid a^{2}\Rightarrow p\mid a$ ---, que vale \emph{porque} $p$ é primo. Daí
sai a \textbf{descida}: de $a^{2}=p\,b^{2}$ vem $a=pk$, logo $b^{2}=pk^{2}$ com $b<a$, e em
$\mathbb{N}$ não há descida infinita. \emph{É a mesma descida do Teor.~\ref{thm:corte-fixo},
e é prova e não varredura: mede-se o PASSO.} Donde:
\[
\underbrace{p\ \text{não factoriza em }\mathbb{Z}}_{\text{DISCRETO}}
\qquad\Longleftrightarrow\qquad
\underbrace{\sqrt p\ \text{não é razão de inteiros}}_{\text{CONTÍNUO}}
\]

\textbf{(2) E quem autoriza a travessia é o TEOREMA CENTRAL.} O
Teor.~\ref{thm:central-energia} põe as três peças: \emph{Hurwitz conta} --- o lado exacto,
discreto ---, \emph{Gentil integra} --- o lado sem grau, contínuo --- e \emph{Lebesgue
transporta sem perder a ORDEM}, que é o degrau $\mathbb{Q}\to\mathbb{R}$. É precisamente a
ordem o que aqui atravessa: o corte do Teor.~\ref{thm:folhas-corte} decide-se por
$\operatorname{tr}$ e $\mathrm{N}$, ambas inteiras, e o que elas dizem do lado contínuo é de
que lado da recta o racional cai. \emph{A ponte não transporta valores: transporta ordem.}

\textbf{(2b) E O DUAL TEM AS SUAS DUAS, que são a VOLTA.} A
Def.~\ref{def:espaco} põe-nas: \textbf{(D1)} $\alpha=\alpha^{\ast}\iff\alpha\in\mathbb{Q}$
--- o que a conjugação fixa é exactamente $\mathbb{Q}$, porque
$\alpha-\alpha^{\ast}=2b\,m\sqrt D$ só se anula com $b=0$, e por (P2) não há outra maneira;
\textbf{(D2)} $\alpha\alpha^{\ast}=a^{2}-b^{2}m^{2}D\in\mathbb{Q}$ \emph{sempre}, porque a
parte com raiz cancela. São o \textbf{traço} e a \textbf{norma}, e são duais de (P1) e
(P2) uma a uma: \emph{(P1) diz quem entra, (D1) quem fica; (P2) diz que a escala não
regressa, (D2) que o dual sempre regressa}. Donde a volta é completa e faz-se
\emph{dentro} do corpo: a inversa é $\alpha^{\ast}/\mathrm{N}(\alpha)$, e o denominador é
a norma --- o produto das duas folhas, que é o primeiro passo da transformada.

\textbf{(2c) E V É ESPAÇO VECTORIAL, que é onde a construção do algébrico começa.}
$V=\mathbb{Q}\oplus\mathbb{Q}\,m\sqrt D$ tem dimensão $2$ sobre $\mathbb{Q}$, e as
coordenadas são \textbf{únicas} --- se $a_1+b_1m\sqrt D=a_2+b_2m\sqrt D$ com $b_1\neq b_2$
então $m\sqrt D=(a_1-a_2)/(b_2-b_1)$ seria racional, contra (P2). \emph{A unicidade das
coordenadas É a propriedade (P2)}, e é também o que torna invertível a matriz de avaliação
$\begin{psmallmatrix}1&\ \ s\\1&-s\end{psmallmatrix}$, de determinante $-2s$. Donde
$\{\operatorname{eval}_{+},\operatorname{eval}_{-}\}$ é uma \textbf{base do dual} $V^{\ast}$,
e a transformada universal é a passagem a essa base --- \emph{não uma construção
acrescentada}. O par traço/norma reparte-se pelos dois graus: o traço é
$\operatorname{eval}_{+}+\operatorname{eval}_{-}$, \textbf{linear}; a norma é
$\operatorname{eval}_{+}\cdot\operatorname{eval}_{-}$, \textbf{quadrática}. É por isso que
são as duas, e não uma.

\textbf{(3) E o SENTIDO diz-se, senão a passagem fica trocada.} Primo é \textbf{suficiente}
e \emph{não} necessário: $\sqrt6$ é irracional e $6$ não é primo. A condição exacta é (P1)
--- $D$ não ser quadrado ---, e é essa que o corpo usa. \emph{O que o primo dá não é a
irracionalidade: é a prova mais curta dela}, porque o lema de Euclides fecha a descida num
passo. $\square$
\end{teorema}
\colunas{as duas propriedades do espaço}
        {(P1) só os quadrados têm raiz racional --- o espaço existe; (P2) um irracional
vezes um inteiro não nulo é irracional --- o espaço tem régua; e primo~$\to$~irracional é
discreto~$\to$~contínuo, pelo lema de Euclides}
        {\code{tests/transformada\_universal.c} \S T12: «raiz racional $\iff$ quadrado» em
$119$ de $119$, com os $9$ quadrados a exibirem a raiz inteira, e $m\sqrt D$ fora de
$\mathbb{Q}$ em $1\,320$ de $1\,320$ pares; \S T11: o passo da descida a fechar nos $17$
primos, o quociente a ser corpo nos $17$, e os $36$ não-quadrados não-primos irracionais na
mesma --- que é o contraste que impede trocar suficiente por necessário}

\begin{teorema}[A TRANSFORMADA UNIVERSAL EM $\mathbb{Q}(m\sqrt D)$: a convolução e a
deconvolução]\label{thm:transformada-universal}
Há \textbf{uma} transformada, e o Teor.~\ref{thm:espectro} já a nomeou: \textbf{avaliação
nas folhas} $\sigma,\sigma\dg$ --- as raízes do polinómio mínimo, \emph{não} raízes da
unidade ---, realizada inteira pela companheira. Não há uma versão discreta e outra
contínua, nem uma rápida e outra lenta: \emph{ela é discreta porque avalia em duas folhas,
e é uma passagem porque só há duas}. Os qualificadores são redundância. Aqui escreve-se no
corpo do Teor.~\ref{thm:serie-quadratica}, sustentado pelas duas propriedades do
Teor.~\ref{thm:duas-propriedades}, onde as folhas são $\pm m\sqrt D$ --- as raízes
de $x^{2}-m^{2}D$ --- e o elemento $\alpha=a+b\,m\sqrt D$ é o polinómio $a+bx$ avaliado
nelas:
\[
\boxed{\;\operatorname{eval}_{\pm}(\alpha)=a\pm b\,m\sqrt D\;}
\]
E o que ela faz são \textbf{duas} operações, que são uma e a sua volta.

\textbf{(1) A IDA é a CONVOLUÇÃO, e a transformada diagonaliza-a.} Pelo
Cor.~\ref{cor:conv-desce} o produto é a convolução dos coeficientes seguida da redução ---
aqui, cíclica de comprimento $2$ com peso $m^{2}D$ no termo que dá a volta:
\[
(a_1,b_1)\circledast(a_2,b_2)=(a_1a_2+m^{2}Db_1b_2,\;a_1b_2+a_2b_1),
\]
e em cada folha
$\operatorname{eval}_\pm(\alpha\beta)=\operatorname{eval}_\pm(\alpha)\cdot
\operatorname{eval}_\pm(\beta)$. \emph{A convolução vira produto ponto a ponto} --- é o
Teor.~\ref{thm:espectro} lido neste corpo, e é a razão de a transformada existir.

\textbf{(2) A VOLTA é a DECONVOLUÇÃO, e é a do Teor.~\ref{thm:decon-andar}.} Dividir na
transformada é dividir folha a folha, e o que decide se isso se pode fazer é aquele
teorema: $w$ é invertível $\iff\gcd(w,D)=1$. Com $D$ livre de quadrados $x^{2}-m^{2}D$ é
irredutível, o quociente é \textbf{corpo}, e a condição \emph{degenera} em «não nulo» ---
por isso a norma sozinha decide. \emph{Não é uma condição forte: é o caso particular que
aquele teorema nomeia.} Onde ela tem conteúdo é com $D$ quadrado --- o polinómio factoriza,
o corpo parte-se, e aparecem os divisores de zero do Teor.~\ref{thm:divzero}, não nulos sem
inversa. \textbf{É aí que a deconvolução deixa de ser sempre possível}, e é esse contraste
que dá conteúdo ao caso irredutível.

\textbf{(2b) E QUEM NORMALIZA É A NORMA --- o primeiro passo já a entrega.} Não é preciso
reduzir representante nenhum antes de avaliar: a inversão em $\mathbb{Q}(m\sqrt D)$ é
$\beta^{\ast}/\mathrm{N}(\beta)$, e $\mathrm{N}(\beta)$ \emph{é o produto das duas folhas}
--- isto é, o primeiro passo da transformada. É o mesmo denominador que o
Teor.~\ref{thm:serie-quadratica} põe na caixa, $\mathrm{N}(\alpha-1)=(a-1)^{2}-b^{2}m^{2}D$.
\emph{Um $\gcd$ por cima não acrescenta nada}: dividir numerador e denominador pelo mesmo
inteiro é multiplicar por $1$, e o elemento do corpo é o mesmo. A normalização não é um
passo a fazer --- \textbf{é o que a avaliação já fez}.

\textbf{(3) E o traço e a norma são o registo inteiro das duas.} A soma das avaliações é
$2a$ --- \textbf{racional, porque a parte com raiz cancela} --- e o produto é
$a^{2}-b^{2}m^{2}D$. São as duas do Teor.~\ref{thm:fixo-dual}, agora com a razão à vista:
\emph{são as únicas funções das folhas que não dependem de qual delas se chama $\sigma$}.
A norma é multiplicativa, logo é o que a convolução preserva; e é ela que a deconvolução
consulta. \textbf{A raiz nunca se forma}: guardam-se os coeficientes.

\medskip
\noindent Nada disto é um teorema novo, e nenhuma das peças precisa de nome próprio: a
diagonalização é o Teor.~\ref{thm:espectro}, a volta é o Teor.~\ref{thm:decon-andar}, a
série é o Teor.~\ref{thm:serie-quadratica}, o par traço/norma é o Teor.~\ref{thm:fixo-dual}.
\emph{O que este enunciado faz é dizer que são a mesma máquina}, e que a máquina tem duas
operações e não quatro verificações. $\square$
\end{teorema}
\begin{corolario}[e é a mesma máquina em TODOS os objectos: as folhas são as raízes de
$x^{n}=w$]\label{cor:folhas-do-polinomio}
O Teor.~\ref{thm:espectro} diz \emph{«não em raízes da unidade»}, e no anel da Lei~8 a
avaliação é $F[x]_k=\sum_i x_i\,g^{ik}$ com $g$ de ordem $N$ --- que são raízes da unidade.
As duas coisas não se contradizem, e a cláusula não é uma proibição: \textbf{é o que as
folhas são num objecto e não noutro}. A máquina é a mesma.

\textbf{O polinómio é sempre a regra da volta.} Escreva-se o anel por
\[
x^{n}=w \qquad\text{(o que se paga ao dar a volta),}
\]
e tudo sai daí: as \textbf{folhas} são as $n$ raízes de $w$, e o produto é a convolução
cíclica de comprimento $n$ com peso $w$ no termo que volta. Os casos com nome são valores
de $(n,w)$:
\begin{center}
$(2,\;m^{2}D)$ --- o corpo quadrático, folhas $\pm m\sqrt D$ \qquad
$(n,\;1)$ --- o anel da Lei~8, folhas as raízes de $x^{n}-1$
\end{center}
\emph{A cláusula do Teor.~\ref{thm:espectro} lê-se então sem excepção}: as folhas são
raízes da unidade exactamente quando $w=1$, e não o são quando $w\neq1$. É a mesma
avaliação, o mesmo código, a mesma companheira --- muda o objecto, não a máquina.

\textbf{E a divisão de trabalho entre as duas operações fica à vista.} A
\textbf{convolução} diagonaliza em \emph{qualquer} folha que exista --- não precisa de
todas. Quem precisa das $n$ é a \textbf{deconvolução}: com as $n$ folhas o núcleo da
avaliação é trivial e a volta existe; quando faltam folhas o núcleo é não trivial e a
informação perde-se. É o Teor.~\ref{thm:decon-andar} dito na transformada. $\square$
\end{corolario}
\colunas{uma máquina só: as folhas são do polinómio}
        {$x^{n}=w$ dá as folhas e a convolução; $w\neq1$ é o corpo quadrático e $w=1$ o
anel da Lei~8; a ida basta-se com uma folha, a volta exige as $n$}
        {\code{tests/transformada\_universal.c} \S T8: em $76$ objectos $(n,w,P)$ --- $67$
com $w\neq1$ e $9$ com $w=1$ --- a convolução diagonaliza em $4\,896$ de $4\,896$
avaliações, com um código só; \S T8b: nos $23$ objectos com as $n$ folhas o núcleo é vazio
em $23$, e nos $26$ a que faltam folhas é não vazio em $26$}

\begin{corolario}[$\operatorname{Dir}$ é o que a transformada vê; $\operatorname{Cruz}$ é a
obstrução à diagonalização simultânea]\label{cor:dir-cruz-folhas}
A transformada diagonaliza \emph{um} operador de cada vez. Para diagonalizar \textbf{dois}
com a mesma avaliação é preciso que partilhem as folhas --- e o que mede essa falta é
$\operatorname{Cruz}$.

\textbf{(1) A conta é fechada.} Para $A_m,A_n$ com $\det=-1$,
\[
\boxed{\;\operatorname{Cruz}(A_m,A_n)=\tfrac12\bigl(A_mA_n-A_nA_m\bigr)
=\tfrac12(m-n)\begin{pmatrix}0&1\\-1&0\end{pmatrix}\;}
\]
--- \textbf{sempre antissimétrica}, e \emph{proporcional à diferença dos traços}.

\textbf{(2) Donde a leitura, em cadeia.} Com o determinante fixo, o traço determina as
folhas (Teor.~\ref{thm:fixo-dual}), logo
\[
\operatorname{Cruz}=0
\iff \operatorname{tr}A_m=\operatorname{tr}A_n
\iff \text{MESMAS FOLHAS}
\iff \text{a mesma avaliação diagonaliza os DOIS.}
\]

\textbf{(3) E é isto que autoriza o passo \textsc{opera}.} A convolução só vira produto
\emph{ponto a ponto} quando os operadores partilham as folhas; se não partilham, não há
transformada comum e o produto não separa. \emph{$\operatorname{Dir}$ é o que a
transformada vê --- o produto folha a folha; $\operatorname{Cruz}$ é o que sobra quando as
folhas não coincidem.}

\medskip
\noindent E fecha com a síntese: \emph{o traço mede} --- e $\operatorname{Cruz}$ mede a
\textbf{diferença} de traços. Quando ela anula, a operação faz-se na transformada com dois
escalares; quando não, o par não tem espectro comum. $\square$
\end{corolario}
\colunas{$\operatorname{Cruz}$ e a diagonalização simultânea}
        {$\operatorname{Cruz}(A_m,A_n)=\tfrac12(m-n)J$, antissimétrica e proporcional à
diferença dos traços; anula-se exactamente quando as folhas coincidem}
        {\code{tests/transformada\_universal.c} \S T6: nos $64$ pares $A_m,A_n$ com
$m,n\leq8$, «$\operatorname{Cruz}=0\iff$ mesmas folhas» em $64$, $\operatorname{Cruz}$
antissimétrica em $64$ e proporcional a $m-n$ em $64$}

\begin{teorema}[AS FOLHAS SÃO O CORTE: a transformada universal é o mecanismo que parte
o objecto, e o corpo é o que sobra da partição]\label{thm:folhas-corte}
\emph{A raiz disto é a propriedade} \textbf{(P1)} \emph{da Def.~\ref{def:espaco}}: só os
quadrados têm raiz racional, logo com $D$ não quadrado as folhas SEPARAM --- e é a
separação que produz o corte.
O Teor.~\ref{thm:transformada-universal} dá a máquina --- $\operatorname{eval}_\pm$, avaliação
nas duas folhas, inteira. Aqui diz-se o que ela \emph{faz}: ela não lê o objecto, ela
\textbf{corta-o}, e o corte que produz é o do Teor.~\ref{thm:corte}. As duas coisas são
uma.

\textbf{(1) O corte de $\sigma_m$ é a avaliação de um elemento nas suas duas folhas.}
Para $a/b\in\mathbb{Q}$ com $b\geq1$ ponha-se
\[
\beta \;=\; 2a-mb-b\sqrt D \;=\; 2b\Bigl(\tfrac ab-\sigma_m\Bigr),
\qquad D=m^{2}+4 ,
\]
que é o elemento $\xi+v\sqrt D$ de $\mathbb{Z}[\sqrt D]$ com $\xi=2a-mb$ e $v=-b$. As suas
duas folhas são $\operatorname{eval}_\pm(\beta)=\xi\pm v\sqrt D$, e as duas funções
simétricas delas --- as \emph{únicas} que não dependem de qual folha se chama $\sigma$
(Teor.~\ref{thm:fixo-dual}) --- são \textbf{ambas inteiras}:
\[
\operatorname{tr}\beta=2\xi ,\qquad \mathrm{N}\beta=\xi^{2}-b^{2}D .
\]
E é só com elas que o lado se decide:
\[
\boxed{\;\frac ab<\sigma_m \iff \xi<0 \ \ \text{ou}\ \ \mathrm{N}\beta<0\;}
\]
--- que é, letra por letra, a partição em dois casos do Teor.~\ref{thm:corte}. \emph{A raiz
nunca se forma}: a decisão inteira que ali aparecia como um artifício de cálculo é a
transformada a fazer o seu trabalho.

\textbf{(2) O corte é a SEPARAÇÃO das folhas, e o discriminante mede-a.} Para
$\alpha=u+v\sqrt D$ qualquer, a separação das duas folhas é
$\operatorname{eval}_+-\operatorname{eval}_-=2v\sqrt D$, e o seu quadrado calcula-se das
duas simétricas sem sair de $\mathbb{Z}$:
\[
\bigl(\operatorname{eval}_+-\operatorname{eval}_-\bigr)^{2}
=\operatorname{tr}^{2}\alpha-4\,\mathrm{N}\alpha=4v^{2}D .
\]
Logo \textbf{as duas folhas coincidem exactamente quando $\alpha\in\mathbb{Q}$}, que é o
Teor.~\ref{thm:descida} ($x=x\dg\iff x\in\mathbb{Q}$) dito na transformada. E daqui sai a
leitura que faltava:
\begin{center}
\emph{folha única $\Rightarrow$ o objecto é um \textbf{ponto}; folhas separadas
$\Rightarrow$ o objecto é um \textbf{corte}.}
\end{center}
Um racional não corta nada --- está lá. É a separação que produz as duas classes.

\textbf{(3) E a separação é genuína: nenhum irracional se disfarça.} Com $D$ não quadrado
e $v\neq0$, $4v^{2}D$ não é quadrado; logo a folha $\operatorname{eval}_+$ não é racional e
o par $(\operatorname{eval}_+,\operatorname{eval}_-)$ não colapsa por acidente aritmético.
Em particular $\mathrm{N}\beta=\xi^{2}-b^{2}D\neq0$ para todo $b\geq1$: \textbf{nenhum
racional cai em cima do corte}, que é a condição de corte do Teor.~\ref{thm:corte},
obtida aqui \emph{da transformada} e não de uma varredura.

\textbf{(4) Donde: o corpo é a partição.} Juntando, a transformada universal
factoriza $\mathbb{Q}(\sqrt D)$ em duas folhas, e essa factorização \emph{é} a construção
do corpo:
\begin{itemize}
\item o produto das folhas --- a norma --- dá o \textbf{lado} do corte (parte 1);
\item a diferença --- o discriminante --- dá a \textbf{abertura} (parte 2), e distingue
ponto de corte;
\item a soma --- o traço --- dá a \textbf{classe} de equivalência, que é o
Teor.~\ref{thm:descida}: $\mathbb{Q}$ são as classes dos pontos fixos sob o traço;
\item e a diagonalização simultânea (Cor.~\ref{cor:dir-cruz-folhas}) diz \emph{quando} dois
objectos são cortados pela \textbf{mesma} partição: quando os traços coincidem.
\end{itemize}
Nada disto sai de $\mathbb{Z}$. O corpo $\mathbb{Q}(\sqrt D)$ não é dado e depois medido:
ele é \textbf{o que sobra} quando a transformada parte $\mathbb{Q}$ nas folhas certas, e as
duas simétricas são o registo inteiro dessa partição. $\square$
\end{teorema}
\colunas{as folhas são o corte}
        {$\operatorname{tr}$ e $\mathrm{N}$, ambas inteiras, decidem o lado do corte de
$\sigma_m$ sem formar a raiz; o discriminante $\operatorname{tr}^{2}-4\mathrm{N}$ separa
ponto de corte}
        {\code{tests/transformada\_universal.c} \S T7: $9\,600$ racionais em cinco corpos,
o critério das folhas e a comparação em fracção contínua --- dois programas sem uma linha
em comum --- dão o mesmo lado em $9\,600$, com $7\,871$ abaixo e $1\,729$ acima e nenhum em
cima; \S T7b: em $5\,746$ elementos o discriminante anula-se exactamente sobre os
racionais, e nos $5\,304$ de folhas separadas nunca é quadrado}

\colunas{a construção dual desce: os pares derivados}
        {posicional\,/\,algébrica separadas por (P1); $\det=$ termo constante e Newton são
(D2) e (D1); convolução\,/\,deconvolução são a ida e a volta}
        {\code{tests/transformada\_universal.c} \S T16, por DOIS CAMINHOS: $x-b$ com uma
raiz racional em $19$ de $19$ e $\mu$ de discriminante não quadrado sem nenhuma em $42$ de
$42$; o determinante saído da MATRIZ a bater com o termo constante do POLINÓMIO em $110$ de
$110$; e $P_k$ da recorrência a bater com $\operatorname{tr}(A^{k})$ da matriz multiplicada
de facto em $770$ de $770$}

\colunas{a transformada universal em $\mathbb{Q}(m\sqrt D)$}
        {avaliação nas folhas $\pm m\sqrt D$; a CONVOLUÇÃO vira produto ponto a ponto, a
DECONVOLUÇÃO é a divisão folha a folha e degenera porque o corpo é irredutível; traço e
norma são o registo inteiro, e é a NORMA que normaliza --- nenhum $\gcd$ acrescenta}
        {\code{tests/transformada\_universal.c}: a convolução cíclica de $2$ com peso $m^{2}D$
e a avaliação a diagonalizá-la em $6561$ de $6561$; traço racional e norma a bater com a
fórmula fechada em $121$ de $121$; a inversão em $168$ de $168$ com o ÚNICO caso de norma
zero a ser o elemento nulo --- a degenerescência do corpo; e o contraste com $D$ quadrado,
onde aparecem $12$ não nulos de norma zero, que são os divisores de zero do
Teor.~\ref{thm:divzero}}

%======================================================================


\section{MELLIN --- o polo MULTIPLICATIVO da transformada}\label{sec:mellin}

\noindent\emph{Vem do centro:} a transformada é UMA só, a avaliação nas folhas
(\code{algebrico thm:espectro}), e ela decompõe-se nos \textbf{dois polos duais}. Este
documento é o polo \textbf{multiplicativo}; o aditivo está no
\code{topologico sec:fourier}.

\[
\boxed{
\begin{array}{lll}
\text{no } \tau=-1 & \textbf{FOURIER} & \text{grupo ADITIVO, o CÍRCULO, } \abs{t}\leq2\\[2pt]
 & & \text{ELÍPTICO: ordem finita, tudo volta}\\[6pt]
\textbf{aqui} \ (\tau=+1) & \textbf{MELLIN} & \text{grupo MULTIPLICATIVO, caracteres } x^{-s}\\[2pt]
 & & \text{as folhas RECÍPROCAS, } \sigma\sigma\dg=-1\\[2pt]
 & & \text{HIPERBÓLICO: o traço CRESCE, e NÃO fecha}
\end{array}}
\]

\noindent\textbf{Por que é que Mellin é deste lado.} Porque os seus caracteres são
\emph{potências} $x^{-s}$ --- as exponenciais vestidas de função-potência --- e enquanto o
grupo é $\mathbb{R}^{+}$ o índice $s$ é \textbf{contínuo}: não há nada de discreto. As
folhas do metal \emph{não estão no círculo}; são recíprocas, com $\abs{\sigma}>1$, e o
operador \textbf{não volta a si}. É exactamente a propriedade que define $\tau=+1$
(\code{algebrico thm:cuspide}), e é a mesma pela qual a completude, o supremo e o alcance
de $\mathbb{R}$ vivem aqui.

\textbf{E o que torna a dourada DISCRETA é a borda}, não Mellin: enquanto o grupo é
$\mathbb{R}^{+}$ o índice é contínuo, e é ao passar à borda que ele se fecha num número
finito de passos --- o factor é $N$ e não $\sqrt N$
(\code{algebrico thm:dourada-discreta}). \emph{O contínuo é deste lado; o que o discretiza
é a passagem, e ela é canónica.}

\medskip
\noindent\emph{A tríade obedece à dualidade:} a transformada fica no centro porque é
\textbf{uma}; os dois polos descem para os lados porque são \textbf{dois}, e a reflexão
$\tau\mapsto-\tau$ troca-os --- soma $\leftrightarrow$ produto, círculo
$\leftrightarrow$ folhas, fechar $\leftrightarrow$ não fechar.

\section{As realizações do polo MULTIPLICATIVO}\label{sec:realiz-mult}

\begin{teorema}[A CONTINUAÇÃO TEM IDA E VOLTA, e as duas são exactas]\label{thm:continuacao}
\emph{Os polos de que aqui se fala são as folhas da \code{algebrico def:espaco}}, e a ida e a
volta deste teorema são as duas metades que lá estão: sair de $\mathbb{Q}$ e regressar.
A continuação analítica costuma dizer-se numa direcção: a série vive num disco, a forma
fechada vive no plano, e o que sai do disco não se perde --- \emph{muda-se de fórmula, não
de objecto}. Mas isso é meia lei. A outra metade é a \textbf{volta}, e é ela que faz da
continuação uma \emph{bijecção} e não uma projecção:
\[
\boxed{
\begin{array}{ll}
\textbf{IDA} & m\ \longrightarrow\ t_k=\sigma^{k}+(\sigma\dg)^{k}\ \longrightarrow\
L(x)=\sum t_k x^{k}/k\ \longrightarrow\ -\log(1-mx-x^{2})\\[3pt]
\textbf{VOLTA} & L\ \longrightarrow\ c_k\ \longrightarrow\ t_k=k\,c_k\ \longrightarrow\
\text{o OPERADOR}
\end{array}}
\]

\textbf{(1) A volta faz-se sem sair de $\mathbb{Q}$.} Multiplicam-se os coeficientes da
forma fechada por $k$, e o que sai \emph{tem} de ser a sequência de traços: \textbf{inteira},
a satisfazer a recorrência do operador, e com $t_1=m$. \emph{O operador lê-se de volta na
série} --- a ida não perde nada.

\textbf{(2) E ela é ÚNICA, que é o que lhe dá conteúdo.} Dois operadores diferentes dão
séries diferentes; se não dessem, «recuperar $m$ da série» não queria dizer nada, porque
dois objectos com a mesma imagem tornariam a inversa uma \emph{escolha}. Mede-se par a
par, em $\mathbb{Q}$ e sem régua.

\textbf{(3) E a volta É o \code{algebrico thm:residuo-fixo}.} Extrair o coeficiente é somar
sobre os polos, e os polos \emph{são} os pontos fixos: $t_k=\sigma^{k}+(\sigma\dg)^{k}$ é a
soma dos resíduos. \emph{A continuação de volta não devolve um número: devolve o par
dual.} É por isso que os metais aparecem na continuação --- dentro do disco eles não se
veem, e nos polos são tudo o que há (\code{algebrico thm:corte-fixo}).

\textbf{(4) E há uma segunda volta, que é a do DOMÍNIO.} A continuação analítica prolonga
o \emph{valor}; a projectiva prolonga o \emph{lugar}: em $\mathbb{P}^1$ o polo em $0$ é
ponto ordinário, e $S$ troca $0\leftrightarrow\infty$ com $S^{2}=I$
(\code{algebrico thm:aditiva}). \emph{Uma acrescenta plano à função, a outra acrescenta ponto
ao espaço} --- e as duas dizem «o buraco era da fórmula, não do objecto». $\square$
\end{teorema}\begin{corolario}[a cúspide na CONTINUAÇÃO: o polo duplo, e o $k$ que
aparece]\label{cor:cuspide-continuacao}
O Teor.~\ref{thm:continuacao} põe a ida e a volta, e os \textbf{polos} de
$L(x)=-\log(1-mx-x^{2})$ são as folhas. A cúspide é onde eles \textbf{colidem}, e a
consequência lê-se nas duas direcções.

\textbf{A recorrência dos traços é a mesma nos três regimes} ---
$t_{k+1}=\operatorname{tr}t_k-\det t_{k-1}$, com $t_0=2$, $t_1=\operatorname{tr}$, tudo
inteiro ---, \emph{mas a base das suas soluções muda ao passar por $\tau=0$}:
\[
\begin{array}{lll}
\tau=+1 & \{\sigma^{k},\,(\sigma\dg)^{k}\} & \text{polos separados: os traços CRESCEM}\\
\tau=\ \ 0 & \{\lambda^{k},\,k\,\lambda^{k}\} & \textbf{polo DUPLO: aparece o }k\\
\tau=-1 & \text{folhas conjugadas} & \text{polos no círculo: OSCILA, e fecha}
\end{array}
\]
Na cúspide $t_k=2\lambda^{k}$ exactamente, e a segunda solução da recorrência é
$k\,\lambda^{k}$ --- ela \emph{só existe} porque o polo é duplo.

\textbf{E o $k$ que aparece é o mesmo $k$ da volta.} O Teor.~\ref{thm:continuacao}(1)
recupera o operador por $t_k=k\,c_k$; aqui o $k$ entra na própria base das soluções. Não são
dois $k$: é a \textbf{ordem do polo} a mostrar-se no coeficiente --- na ida como o
factor que a segunda solução carrega, na volta como o factor que a extracção repõe.
\emph{A cúspide é o único sítio onde as duas direcções da torre se veem uma à outra.}

\textbf{E as três peças da tríade dizem isto cada uma na sua coluna.} A continuação tem ida
e volta (Teor.~\ref{thm:continuacao}); o que muda de peça para peça é \emph{onde} a
reflexão se realiza:
\[
\begin{array}{lll}
\textbf{CANTOR} & \text{a escrita} & \text{dígitos }\{-1,0,+1\}:\ \text{reverter troca os}
\ \pm1\ \text{e deixa o }0\\[2pt]
\textbf{JULIA} & \text{o operador} & \nu^{2}=\mathrm{id}:\ \text{a reflexão em pessoa,
ordem }2\\[2pt]
\textbf{VIVIANI} & \text{o meio-ângulo} & i^{2}=-1:\ \text{a reflexão é a METADE dele,}
\ \operatorname{ord}(i)=2\operatorname{ord}(\nu)
\end{array}
\]
Cantor reflecte \emph{onde se escreve}, Julia reflecte \emph{o que se aplica}, e Viviani é
a sua raiz quadrada --- por isso a ordem $4$. \emph{As três reflectem o mesmo
$\{-1,0,+1\}$}, e as três fixam o zero: a cúspide é o que nenhuma delas move, e é por isso
que o teste $\operatorname{rev}(\tau)=\tau$ vale nas três colunas sem mudar de forma.
$\square$
\end{corolario}\begin{teorema}[o corte É o ponto fixo --- com ida e volta, que é a
dualidade]\label{thm:corte-ponto-fixo}
O \code{algebrico thm:real-caminho} diz que \emph{a folha não é nó}, e fecha-o com uma
impossibilidade concreta: $(2m+N)^{2}=5N^{2}$. Essa impossibilidade tem nome, e
generaliza-se a toda a família --- \textbf{ela é a equação do ponto fixo}.

\emph{(i) O andar tem três pontos, e a Lei 0 é que os fecha.} Em $\mathbb{P}^1$ são
$0=[0:1]$, $1=[1:1]$ e $\infty=[1:0]$; e por $0^{\dagger}=\infty$
(\code{teoria cor:zero-infinito}) o infinito \textbf{não é um número grande}: é um ponto
como os outros. A acção de $A_m$ é então a matriz vezes o vector, \emph{sem uma
divisão},
\[
[p:q]\ \longmapsto\ [\,m\,p+q\ :\ p\,],
\]
e a órbita de $\infty$ é a sucessão dos convergentes. \textbf{O ponto de partida só
existe por causa da Lei 0}: num corpo sem ela, «começar no infinito» não quer dizer
nada.

\emph{(ii) O ponto fixo é o vector próprio, e são dois.} $x$ é fixo por $g(x)=m+1/x$
exactamente quando $[x:1]$ é vector próprio de $A_m$, isto é $x^{2}=mx+1$; as duas
raízes têm soma $\operatorname{tr}A_m=m$ e produto $\det A_m=-1$ --- \emph{inteiros},
e é o par dual $(\sigma,\sigma^{\dagger})$ do \code{algebrico thm:gato}.

\emph{(iii) E o corte é a falta dele.} Em coordenadas homogéneas o ponto fixo pede
$p^{2}-mpq-q^{2}=0$, e \textbf{completando o quadrado} isso é
\[
\boxed{\;4\bigl(p^{2}-m\,p\,q-q^{2}\bigr)=(2p-mq)^{2}-D\,q^{2},
\qquad D=m^{2}+4\;}
\]
--- identidade em $\mathbb{Z}$. Logo
\[
\boxed{\;\text{o ponto fixo vive em }\mathbb{P}^1(\mathbb{Q})
\iff D\ \text{é quadrado perfeito}\;}
\]
e a questão deixa de ser sobre racionais para ser sobre um inteiro. \textbf{E aí
decide-se em duas linhas, para todo $m\geq1$}: $m^{2}<D<(m+2)^{2}$ --- a segunda
desigualdade é $0<4m$ ---, logo o único quadrado possível no meio é $(m+1)^{2}$; e
$D=(m+1)^{2}$ pede $2m=3$, que não tem solução em $\mathbb{Z}$. \emph{É o passo, e
não a lista.}

Portanto: a órbita corre toda em $\mathbb{P}^1(\mathbb{Q})$ e o ponto para onde ela
vai não está lá. \textbf{O corte não é uma construção acrescentada por fora: é o
ponto fixo da Möbius visto do andar que não o contém}, e o
$(2m+N)^{2}=5N^{2}$ do \code{algebrico thm:real-caminho} é o caso $m=1$.

\emph{(iv) E a face finita di-lo pelo contrário.} A \emph{mesma} equação
$x^{2}=mx+1$ tem raiz em $\mathbb{F}_p$ exactamente quando $D$ é resíduo quadrático
--- e aí o ponto fixo \textbf{está} no andar, a órbita cai nele, e não há corte
nenhum a fazer. Donde o enunciado na forma que o torna estrutural:
\[
\boxed{\;\text{o corte é a equação do ponto fixo a \textbf{não fechar} no corpo onde
a órbita corre}\;}
\]
--- e não um defeito de $\mathbb{Q}$.

\emph{(v) E há VOLTA --- é a metade dual, e sem ela isto seria meia lei.} A ida é a
órbita; a volta existe porque $\det A_m=-1$, e é o \textbf{fator de potência unitário}
que a torna exacta:
\[
\boxed{\;A_m^{-1}=\frac{1}{\det}\operatorname{adj}A_m
=\begin{pmatrix}0&1\\1&-m\end{pmatrix}\ \text{---\ \textbf{inteira}},
\qquad
[p:q]\ \longmapsto\ [\,q\ :\ p-mq\,]\;}
\]
E essa acção é, \emph{letra por letra}, a descida do ponto (iii). \textbf{A descida não
era um truque de teoria dos números: é a órbita para trás} --- e é por isso que ela
encolhe, porque a ida cresce como $\sigma^{k}$ e a volta desce pelo mesmo caminho, a
chegar ao $\infty$ exacto.

\noindent E o par $0\leftrightarrow\infty$ é o caso mais simples de todos: a troca
$S=\begin{pmatrix}0&1\\1&0\end{pmatrix}$ tem $S^{2}=I$, logo é a \emph{sua
própria} inversa --- \textbf{ida e volta pela mesma matriz}, que é o que a Lei 0 diz.
Donde a forma dual do enunciado:
\[
\boxed{
\begin{array}{ll}
\text{a \textbf{ida}} & \text{a órbita de }\infty\text{: os convergentes, a crescer}\\[3pt]
\text{a \textbf{volta}} & A_m^{-1}\text{ inteira: a descida, a encolher}\\[3pt]
\text{e o que a \textbf{fecha}} & \lvert\det A_m\rvert=1
\end{array}}
\]

\emph{Medido:} \code{tests/corte\_ponto\_fixo.c} ($6{:}0$) --- $264$ passos da órbita
com $\lvert\det\rvert=1$; Cayley--Hamilton nos metais; a forma canónica em $31\,212$
triplos; o passo nos seus dois membros; $62$ dos $126$ metais de
$\mathbb{F}_{127}$ com ponto fixo, por duas rotas concordantes, contra \emph{zero}
em $\mathbb{Q}$ para todo $m$; e a volta: $A_mA_m^{-1}=I$, a órbita a regressar ao
$\infty$ exacto, $S^{2}=I$, e $16\,128$ pares (metal, ponto) a fechar na face finita.
$\square$
\end{teorema}\begin{teorema}[E O QUE O FUNDAMENTA É O TEOREMA CENTRAL: Hurwitz conta, Gentil integra,
Lebesgue casa-os]\label{thm:central-energia}
O \code{algebrico thm:conservacao} diz \emph{que} a energia se conserva. Falta dizer
\textbf{por que é que ela pode ser conservada} --- e a resposta é o \textbf{teorema
central Gentil--Hurwitz}, que vive hoje no \emph{Corpo Analítico}
(\code{analitico sec:central}) por não estar na cadeia principal, e que se usa em
\code{teoria thm:parseval-multi} --- \emph{o Parseval multidimensão, no anel escolhido} ---
para derivar a energia; e cujo corolário
(\code{algebrico cor:verifica-energia}) já dizia a frase toda: \emph{a indução legítima CONSERVA
$E$; a que viola o invariante muda $E$, e a meta-indução acusa}.

\textbf{(1) HURWITZ CONTA --- e é o lado exacto.} A órbita completa da raiz fecha:
\[
\boxed{\;\sum_{k=0}^{N-1}\omega^{(i-j)k}=N\cdot\delta_{ij}\;}
\]
--- $N$ na diagonal, \textbf{zero} fora. Os cruzados não são pequenos: cancelam-se
\emph{exactamente}, e é daqui que Parseval sai. \emph{É o corte discreto.}

\textbf{(2) GENTIL INTEGRA --- e é o lado que não tem grau.} A mesma linha, integrada,
é $\int\abs{f}^{2}=\int\abs{F}^{2}$ com a medida de Lebesgue.

\textbf{(3) LEBESGUE CASA-OS, e o que ele preserva é a ORDEM.} A $\sigma$-aditividade
transporta a contagem para o integral \emph{sem perder a ordem} --- é o degrau
$\mathbb{Q}\to\mathbb{R}$. \emph{A soma reversível é o que autoriza a ponte}, e
reversível é a mesma palavra do \code{algebrico thm:conservacao}(3): conservar e poder voltar.

\textbf{(4) E O TECTO DE OITO É DA NORMA, NÃO DO OBJECTO.} Hurwitz limita a
$n\in\{1,2,4,8\}$ quando se exige norma \emph{euclidiana} multiplicativa \textbf{e}
operação bilinear. A recta não paga esse tecto porque a energia que ela conserva
\textbf{não é a euclidiana}: é o \textbf{determinante}, que é multiplicativo em todo $n$.
\emph{Foi por isso que $\abs{\det}=1$ pôde ser a lei de conservação em todos os andares.}
Gentil escapa pelo mesmo sítio, e paga com a bilinearidade --- \emph{o que ele cede não é
a norma, é a bilinearidade}. $\square$
\end{teorema}

\noindent\emph{Veio do centro, e veio para este polo} (\code{analitico sec:mellin}): o
Teorema Central do \textbf{contínuo}, que é realização específica do lado que não fecha.
A régua vem com ele.

\begin{teorema}[o Teorema Central do contínuo --- a classe racional
encaixota o irracional por baixo]\label{thm:central-continuo}
Os quatro pontos do gerente, cada um com medida inteira exata:
\begin{enumerate}
\item \textbf{existência e unicidade do limite}: dois habitantes de
todos os intervalos até a profundidade $K$ coincidem ---
$|rs'-r's|$ é inteiro $<ss'/2^{K}<1$, logo \emph{zero} (o pombal);
distintos são separados na profundidade $\log_{2}(ss')$.
\item \textbf{sobrejetividade, na classe comparável}: todo $x$ com
teste de chão exato tem caminho --- racionais, quadráticos ($1/\varphi$,
$\sqrt2-1$) e $\pi-3$ até à profundidade do certificado ($23$ bits) ---
e a \emph{volta} põe $x$ em todos os intervalos.
\item \textbf{unicidade módulo a borda}: os diádicos --- e só eles ---
têm \emph{dois} caminhos ($1000\ldots=0111\ldots$, com a soma
geométrica exata a fechar no nó); o não-diádico tem um.
$\mathbb{R}\cong\{\text{caminhos}\}/\!\sim$, e $\sim$ identifica
exatamente \emph{os nós da árvore}: a borda é a ramificação.
\item \textbf{completude e o encaixotamento}: a classe racional
$\{m_{k}/2^{k}\}$ \emph{cresce}, fica \emph{toda abaixo} de $x$, e
\emph{ultrapassa todo racional} $<x$ --- $x$ é o \textbf{supremo
exato} da sua classe (o corte de Dedekind em três cláusulas
verificadas); e a sequência de Cauchy dos truncados estabiliza bit a
bit \emph{dentro} do objeto.
\end{enumerate}
O Teorema Central lê o conjunto: Hurwitz \emph{conta} a borda (os
nós, magros: $0{,}246\to0{,}015\to0{,}001$), Gentil \emph{casa} o
endereço, Lebesgue \emph{mede} o que sobra. \emph{A reta não é onde
algo está sentado --- é o limite dos caminhos da torre.}
\emph{Medido:} \code{tests/central\_continuo.js} ($7{:}0$; assinado
por $\mathcal{M}$).
$\square$
\end{teorema}

\section{A estrela}\label{sec:estrela}

Tudo isto corre, e não corre numa máquina: \textbf{corre numa estrela}. Não é uma maneira de falar
--- é o que a única instrução diz.

\subsection*{Uma instrução, dois sentidos}

\[
  \texttt{MOVE(slot, sentido)} : \qquad
  \underbrace{+1 \ \text{lê}}_{\textbf{absorve --- o lado negro}}, \qquad
  \underbrace{-1 \ \text{escreve}}_{\textbf{emite --- o lado branco}} .
\]

\noindent \textbf{É a Lei 1 aplicada à entrada e à saída}: ler e escrever são a mesma operação com o
sinal trocado. Não há instrução de leitura e instrução de escrita --- há uma, e um sinal. E os dois
sentidos são exactamente os dois buracos: \emph{absorver é o negro, emitir é o branco}.

\smallskip
\noindent \textbf{Logo o sistema é a estrela.} Metade dele absorve e metade emite; um buraco negro
só faria a primeira, um buraco branco só a segunda, e nenhum dos dois seria reversível. A estrela é
a \emph{interface} entre eles --- e o sistema é a interface porque tem \code{MOVE} nos dois sentidos
e mais nada.

\begin{center}\small
\begin{tabular}{@{}llll@{}}
\toprule
 & o que faz & na instrução & é reversível? \\
\midrule
\textbf{buraco branco} & só emite & só \code{MOVE($\cdot$,\,$-1$)} & não \\
\textbf{a estrela} & \textbf{emite e absorve} & \textbf{\code{MOVE} nos dois} & \textbf{sim} \\
\textbf{buraco negro} & só absorve & só \code{MOVE($\cdot$,\,$+1$)} & não \\
\bottomrule
\end{tabular}
\end{center}

\subsection*{Sem memória, e sem estado}

O sistema é \textbf{reversível}: se toda representação tem dual, todo passo se desfaz, e um passo
que se desfaz não precisa de guardar nada para voltar. Daí a regra de projecto, e é a ideia central:
\emph{não há RAM, e não há estado}. A estrela é uma \textbf{interface de transformação} --- o corpo
entra, atravessa, sai, e \emph{nada fica}; \textbf{não guarda nada de ninguém}. O que parece residir
é a \textbf{régua}, a assinatura do corpo (dezenas de bytes, não megabytes), mas ela \emph{cavalga o
corpo} --- entra discreta e sai contínua --- e não é estado retido pela estrela. \emph{E não é um
banco}: o banco são os arquivos pessoais de quem opera, e vivem do lado do cliente, não da estrela.

\subsection*{E o estado é o canal do caos --- por isso não há \emph{malloc}}

\noindent Guardar estado não é só \emph{redundante}: é abrir um canal por onde um vazamento propaga.
Uma máquina com estado é uma \textbf{recorrência} $s_{n+1}=f(s_n,u_n)$ --- o estado de um passo
alimenta o seguinte. Um vazamento (um bit corrompido no dado, uma escrita selvagem no estado)
\textbf{entra na recorrência} e contamina tudo o que vem depois; e se $f$ \textbf{expande} --- $|f'|>1$,
como o dobrador $s\mapsto 2s$, o \emph{shift} de Bernoulli --- o vazamento \textbf{amplifica} a cada
passo, $\delta_n=\bigl(\prod f'\bigr)\delta_0$. É \textbf{dependência sensível às condições iniciais},
com expoente de Lyapunov $\lambda=\log|f'|>0$: \textbf{caos}. Um erro no bit invisível torna-se, em
poucos passos, o bit dominante.

\smallskip
\noindent A estrela \textbf{não tem essa recorrência}. Trata cada dado \emph{independente} e
\emph{reversível} --- $\nu\circ\nu=\mathrm{id}$, resíduo $0$ (Teorema~\ref{thm:pontereversivel}): não há
$s_n$ a alimentar $s_{n+1}$, logo o vazamento \textbf{não tem onde acumular} e fica no dado que o
sofreu. E a involução é \textbf{neutra}: $|\nu'|=1$ no ponto fixo $|x|=1$, logo $\lambda=0$ --- nem
cresce nem se guarda. \textbf{Não usar \code{malloc} não é higiene de código: é remover o canal --- o
estado --- que transforma um vazamento num sistema caótico.} É o mesmo facto por dois lados: a órbita
que acumula não fecha (\S da medida, $0{,}\overline9\ne1$); ser o ponto fixo fecha, resíduo $0$,
memória nenhuma.

\medido{acumular estado propaga e amplifica um vazamento: numa recorrência expansiva $s\mapsto 2s+u$,
corromper \emph{um} dado contamina os $35$ de $35$ passos a jusante e a contaminação \textbf{dobra} a
cada passo ($1,2,4,8,16,32,\dots$, Lyapunov $\log2>0$: caos); sem estado, com a estrela
$\nu(x)=-1/x$ por dado, o mesmo vazamento afecta \textbf{só} a saída corrompida ($1$ de $40$) e toda
a outra tem resíduo $0$; e $\nu\circ\nu=\mathrm{id}$ com resíduo $0$ ($\lambda=0$, sem canal, sem
propagação): \code{tests/estado\_caos.c}.}

\subsection*{Sem tempo, e o que se mede em vez dele}

Um sistema reversível \textbf{não tem tempo}: o que se mede como tempo é latência, e a latência não é
uma propriedade do algoritmo. O que se mede é \textbf{bits apagados}, porque apagar é a única
operação que não se desfaz.

\smallskip
\noindent \textbf{E é aqui que o buraco negro paga.} Absorver sem emitir \emph{é} apagar --- não por
semelhança: é a mesma operação. Um passo que só vai num sentido não tem inverso, e o que não se
desfaz custa. A mesma tarefa, com a mesma entrada e a mesma saída, nas duas formas:

\begin{center}\small
\begin{tabular}{@{}lrrl@{}}
\toprule
a forma & bits apagados & o custo mínimo & é qual dos três \\
\midrule
destrutiva & $131\,072$ & $3{,}763\times10^{-16}$ J & \textbf{o buraco negro} \\
involutiva & $\mathbf{0}$ & $\mathbf{0}$ \emph{exacto} & \textbf{a estrela} \\
\bottomrule
\end{tabular}
\end{center}

\noindent E o zero é \emph{exacto}, não um número pequeno: \textbf{sem bit apagado não há mínimo
termodinâmico nenhum a pagar}. A estrela não é mais eficiente que o buraco negro --- está noutro
regime, e é o mesmo trial: um só absorve e paga, o outro faz os dois e não paga.

\smallskip
\noindent \emph{E a régua é essa, não a memória ocupada}: um vector grande que nunca é escrito não
dissipa nada, e um único inteiro reescrito mil milhões de vezes paga mil milhões de vezes. Uma
memória que se refresque dissipa \textbf{por existir}; um disco escrito uma vez fica quieto.

\medido{a única instrução chega: os dois sentidos de \code{MOVE} são inversos em $134$ casos com $0$
falhas --- o que se escreve lê-se de volta ---, o salto é a mesma \code{MOVE} com o \emph{pc} por
destino, e a mesma \code{MOVE} serve $500$ corpos diferentes sem uma instrução nova, porque o corpo é
\emph{campo} e não \emph{opcode}; e o \textbf{controlo} diz o que isso custa: a máquina de duas
precisa de \textbf{dois códigos}, esta de \textbf{um código e um argumento} --- e argumento é dado, e
dado mora no disco (\code{tests/move.c}). E o custo: $131\,072$ bits apagados na forma destrutiva
contra $\mathbf{0}$ na involutiva, o que a $2{,}87098\times10^{-21}$ J/bit dá
$3{,}763\times10^{-16}$ J contra \textbf{zero exacto}; e o \textbf{controlo}: a involutiva \emph{é}
involutiva --- repetida, devolve o original em $4096$ de $4096$, e é por isso que os seus bits não se
perderam (\code{tests/dissipa.c}).}

\section{Como se mede aqui: o medidor é $0$}\label{sec:medir}

Uma asserção do tipo $x=42$ \emph{compara} contra um valor posto por quem a escreve --- e comparar é
ler metade, e metade dissipa. \textbf{A medida não compara: reverte, e lê o resíduo.}

\begin{center}\small
\begin{tabular}{@{}ll@{}}
\toprule
resíduo $0$ & reverte --- e \emph{valida os dois lados ao mesmo tempo} \\
resíduo $\neq0$ & dissipa --- e aí sim é preciso raciocínio: o que se perdeu? \\
\bottomrule
\end{tabular}
\end{center}

\noindent \textbf{E é a prova dos nove.} Não se resolve e depois se confere: o resto módulo nove é um
morfismo, logo a conta e a prova correm na \emph{mesma} passagem e a diferença é zero. Não há passo
de verificação --- há um segundo lado a andar ao lado do primeiro. \emph{Se der diferente, está
errado.} E querer dissipação é só \textbf{tirar a reversão}.

\smallskip
\noindent E o que separa isto de um estilo: \textbf{uma asserção que compara contra um valor escrito
passa no objecto certo \emph{e} no trocado} --- ela verifica a aritmética de quem a escreveu, não o
objecto. A reversão separa-os.

\subsection*{E mede-se \emph{pela metade}}

Um resíduo zero sozinho não mede: não diz que é \emph{só} ali. Um resíduo não-zero sozinho também
não: não diz que existe. \textbf{A medida são as duas metades da mesma passagem}:

\begin{center}\small
\begin{tabular}{@{}ll@{}}
\toprule
o resíduo onde \emph{tem} de ser zero & diz que \textbf{existe} \\
o resíduo onde \emph{não pode} ser zero & diz que é \textbf{único} \\
\bottomrule
\end{tabular}
\end{center}

\noindent É a cruz outra vez, aplicada à própria medição: uma coordenada mede, a outra ordena, e
nenhuma sozinha é o par. \emph{Uma asserção que só olha para um dos lados está pela metade} --- e
estar pela metade é exactamente o que dissipa.

\subsection*{E daí a ferramenta: promover}

Invertendo o sinal da reversão obtém-se \emph{duas vezes} o objecto. Daí o protocolo, numa passagem:
\[
  \boxed{\;S=\frac{x+x^{\dagger}}{2},\qquad A=\frac{x-x^{\dagger}}{2},\qquad S+A=x\;}
\]
\textbf{Cortar ao meio não quebra nada porque as metades reconstroem} --- o que se perde é ficar com
\emph{uma}. E isto não é uma verificação: \textbf{é a operação}. Um valor vira um par, $\dim$ dobra
--- é o $\dim A_{n+1}=2\dim A_n$ da torre --- e encadeia, $1\to2\to4$.

\smallskip
\noindent \textbf{E é a dual da escrita.} Escrever directo põe o novo por cima do velho, e o velho
não volta: isso é apagar, e apagar custa. Escrever o \emph{par} não apaga, porque o $S$ é o centro.
\emph{Quem escreve o par já leu} --- uma passagem onde escrever-depois-ler gasta duas. Em \code{MOVE}:
$-1$ sozinho só emite, $+1$ sozinho só absorve, cada um meia operação; promover faz as duas na mesma
passagem, \textbf{e é por isso que é a estrela}.

\smallskip
\noindent A divisão por dois é \emph{exacta}, e não por sorte: $x+x^{\dagger}=2c$ e
$x-x^{\dagger}=2(x-c)$ são sempre pares porque a involução o garante. \textbf{Num objecto que não
reverte ela falha} --- e é assim que a ferramenta acusa em vez de devolver um número errado.

\medido{sem um único número esperado escrito: estaca, cifra e ordenação-com-caminho dão resíduo $0$;
o período de $J$ sai \emph{sem} se escrever qual é --- quatro passos fecham, dois não; a prova dos
nove fecha em $23\,200$ pares e um erro de uma unidade acusa em todos; retirada a reversão o resíduo
deixa de ser zero em todos os casos; $S+A$ reconstrói $x$ em $48\,861$ casos e nenhuma das $97\,722$
somas é ímpar; promover custa $801$ passagens contra $1\,602$ de escrever-e-ler, e o valor anterior
recupera-se em \emph{todas} as escritas promovidas e perde-se em \emph{todas} as directas; e o
\textbf{controlo}: num impostor a paridade quebra em $7\,813$ de $15\,025$
(\code{tests/reversao.c}, \code{tests/meia\_passagem.c}, \code{lib/promove.h}, \code{tests/promove.c}).}

\section{A ordenação é o corpo estelar em operação}

\subsection*{Uma estrela não compara}

Comparar é medir, e medir é interferir --- \emph{o Sol não se mede para operar: opera}. Uma
comparação entre dois elementos devolve um bit e deita fora tudo o resto; um algoritmo que compara
paga em bits apagados aquilo que a estrela não paga.

\smallskip
\noindent \textbf{E uma estrela não tem fila}: \emph{irradia} (\S\ref{sec:estrela}). É por isso que a
ordenação deste sistema atravessa o disco com \code{MOVE} e não acumula nada pelo caminho.

\subsection*{As peças são as que já havia}

\begin{center}\small
\begin{tabular}{@{}lll@{}}
\toprule
a peça & no corpo estelar & na ordenação \\
\midrule
a \textbf{estaca} & a involução que troca os lados & marca cada valor no seu lugar \\
a \textbf{cruz} & mede o que não se moveu & lê a marca, sem perguntar \\
o \textbf{trial} & negro / estrela / branco & menor / igual / maior, por \emph{sinal} \\
$J$ & o rotor de período $4$ & os quatro quadrantes da saída \\
\code{MOVE} & ler é escrever com sinal trocado & todo o acesso ao disco \\
\bottomrule
\end{tabular}
\end{center}

\noindent \textbf{A referência são os naturais}, e eles aparecem na \emph{transição} entre
dimensões: cada involução do relógio deixa uma marca, e as marcas são a ordem. Nada se produz ---
\textbf{a ordem já estava lá}, e o percurso apenas a lê.

\subsection*{E o custo}

\begin{center}\small
\begin{tabular}{@{}lrl@{}}
\toprule
 & bits apagados & porquê \\
\midrule
uma comparação & $127$ de $128$ & devolve um bit e deita fora o resto \\
uma escrita & $64$ & o valor anterior perde-se \\
uma \textbf{troca} & $\mathbf{0}$ & é involutiva: desfaz-se \\
\bottomrule
\end{tabular}
\end{center}

\noindent Com o caminho guardado --- que é a permutação, e a permutação é o dual --- o sistema
\textbf{volta} ao ponto de partida, e o custo é \emph{zero}.

\medido{$200$ listas saem crescentes e completas descendo pela estaca; o percurso faz $0$ perguntas
de ordem; entram $2000$ e saem $2000$ com a soma conservada; $500$ valores repetidos cabem em $416$
bytes contra $74\,912$ de $500$ distintos; o estado residente são $48$ bytes e \emph{não vê} o $n$;
a volta é exacta em $200\,000$ de $200\,000$ e com o caminho guardado apagam-se $\mathbf{0}$ bits
contra $25\,600\,000$ sem ele; e o \textbf{controlo}: percorrida ao contrário não sobe um único par
(\code{tests/dualsort.c}, \code{tests/dualsort\_banco.c}, \code{tools/bench\_dissipa.c}).}

\section{A especificação, numa página}

\begin{center}\footnotesize
\setlength\tabcolsep{4pt}
\setlength\tabcolsep{4pt}
\begin{tabular}{@{}p{3.1cm}p{4.4cm}p{7.4cm}@{}}
\toprule
a camada & o que a define & o que ela obriga \\
\midrule
a exigência & o que conserva na órbita conserva no espaço & tudo o resto \\
a Lei 1 & $1\dg=-1$ & o que separa; o ponto fixo \\
a Lei 2 & $T\dg=-T$, logo $T^2=-1$ & o passo; o período $4$; a bidualidade \\
a estaca & $x\mapsto x\dg$, involutiva & o movimento \\
a cruz & $(x\oplus x\dg,\,x\otimes x\dg)$ & a medida e a ordem \\
o trial & $\{-1,0,+1\}$ & três regimes, e não há quarto \\
a torre & $\dim A_{n+1}=2\dim A_n$ & as dimensões; a plena em $6$ \\
o corpo estelar & $a=\text{empurra}/\text{puxa}-1$ & \textbf{a trindade}: branco, estrela, negro \\
a geometria & $\Delta=\operatorname{tr}^2-4\det$ & três curvaturas, pelo sinal \\
a equação de campo & conservar & $\alpha=\tfrac12$; $d=4$; $\Lambda=3/D$; $8\pi=2\cdot4\pi$ \\
a expansão & $H^2=1/D$ & quatro equações, duas independentes \\
a cosmologia & $\rho\,a^{3(1+w)}=C$ & três conteúdos; órbita de quatro \\
os dois escuros & um ponto fixo por involução & a matéria e o vácuo; a radiação vê-se \\
as analíticas & o relógio e a família & $\pi$, $e$, $\varphi$, $\sigma_m$, $\Lambda=3/D$ \\
as constantes & $(\text{passo})^a\times$ puro & $c=\sigma$; $G$: $8\pi$; $k$: $4\pi$; $k_B$: câmbio \\
o $c^2$ & $\sigma^j=F_j\sigma+F_{j-1}$ & sai da espiral: duas coordenadas chegam \\
a massa & $M_{ij}=\sum(x_i-c_i)(x_j-c_j)$ & \textbf{vectorial}: o escalar é só o traço \\
o grau & $1$: momento; $2$: massa & duas raízes trazem o par; uma não \\
a renormalização & $\max_p\min(p,q-p)=q/2$ & o número é da régua; os degraus são os metais \\
a roupa & três escalas, uma por sinal & o dimensional sai; o puro não \\
\textbf{a estrela} & \code{MOVE(slot,\,sentido)} & $+1$ absorve, $-1$ emite: \emph{é a interface} \\
a dissipação & apagar é o que não se desfaz & o negro paga; a estrela paga $0$ exacto \\
os dois universos & $a_B=1/a_A$ & Parseval: $\rho_A\rho_B$ não se move \\
a antimatéria & a terceira involução & quatro conteúdos, todos biduais \\
a temperatura & $T=1/r$ & a involução no único comprimento; $S$ é a área \\
a segunda lei & $S_{\text{n}}\!\cdot\!S_{\text{b}}=1$ & vale de \emph{um} lado; o par não se move \\
medir & resíduo $0$, não comparação & a prova dos nove: resolver \emph{e} provar \\
e pela metade & zero \emph{e} não-zero & uma diz que existe, a outra que é único \\
promover & $x\mapsto(S,A)$ & sobe um andar; quem escreve o par já leu \\
a ordenação & irradiar em vez de comparar & $0$ bits apagados \\
\bottomrule
\end{tabular}
\end{center}

\bigskip
\noindent \textbf{E é uma coisa só.} O ponto fixo do trial, o coeficiente $\tfrac12$ da equação de
campo, o patamar da escada, a borda onde a estrela vive e o zero entre os dois sinais são
\emph{o mesmo objecto} visto de cinco sítios. A dimensão quatro que o traço escolhe é o período de
$J$. O $w=-1$ que a integração deixa é o ponto fixo onde a órbita degenera. Nada disto foi
ajustado para coincidir: coincide porque \textbf{a exigência é uma}.

\vfill
{\footnotesize\color{cinza}\noindent
Teoria, catálogo, medidores e o resto do sistema em
\href{https://goldenkingdom.patriatechnology.com}{goldenkingdom.patriatechnology.com}.
Proprietário --- uso mediante contrato e pagamento (ver \texttt{LICENSE}).\par}

\end{document}
